% Lectura y comentario : « El mundo laboral hoy » ; ayudas, subvenciones, acuerdos bilaterales y multilaterales — Espagnol Tle A — Cours signé OnBuch+
% Programme officiel MINESEC. Compiler avec : tectonic E16-mundo-laboral-ayudas-desarrollo.tex
\newcommand{\DOCMATIERE}{Espagnol}
\newcommand{\DOCNIVEAU}{Tle A}
\newcommand{\DOCMODULE}{Módulo 4 : Activités économiques et monde du travail}
\newcommand{\DOCLECON}{E16}
\newcommand{\DOCTITRE}{Lectura y comentario : « El mundo laboral hoy » ; ayudas, subvenciones, acuerdos bilaterales y multilaterales}
% OnBuch+ — préambule « cours signés » ENRICHI (compilé avec tectonic / XeLaTeX)
% Système visuel « Pop » d'OnBuch+ : fond crème (jamais blanc), bordures encre
% épaisses, ombres « dures » décalées, titres Archivo Black en capitales.
% Version MATHÉMATIQUES : graphes (pgfplots/tikz), boîtes pédagogiques
% (définition, propriété, méthode, attention, le savais-tu, exercice résolu,
% exercice, TP, bilan), page de garde et signature OnBuch+ ancrée en bas.
%
% Macros attendues dans main.tex AVANT \input{preamble} :
%   \DOCMATIERE  \DOCNIVEAU  \DOCMODULE  \DOCLECON (numéro)  \DOCTITRE
\documentclass[11pt]{article}

\usepackage{fontspec}
\usepackage[spanish,french]{babel} % français = langue principale ; l'espagnol via \esp{..} / espagnol
\usepackage[a4paper,margin=2cm,top=3.4cm,bottom=2.8cm,headheight=1.4cm,footskip=1.3cm]{geometry}
\usepackage{amsmath,amssymb,amsfonts}
\usepackage{mathtools} % \xmapsto, \coloneqq... souvent utilisés par les modèles
\usepackage{cancel} % simplifications barrées (bilans de forces)
% \abs{z} (module, valeur absolue) et \norm{u} (norme), utilisés par les modèles
\providecommand{\abs}{}\let\abs\relax\DeclarePairedDelimiter{\abs}{\lvert}{\rvert}
\providecommand{\norm}{}\let\norm\relax\DeclarePairedDelimiter{\norm}{\lVert}{\rVert}
\usepackage[table]{xcolor} % [table] : fournit aussi \rowcolors (alternance de lignes)
\usepackage{pagecolor}
\usepackage{tikz}
\usetikzlibrary{arrows.meta,positioning,calc,shapes.geometric,shapes.misc,shapes.arrows,decorations.pathmorphing,decorations.pathreplacing,decorations.markings,patterns,shadows,mindmap,trees,fit,backgrounds,matrix,intersections,angles,quotes}
\usepackage{pgfplots}
\usepackage[version=4]{mhchem} % équations nucléaires \ce{^{235}_{92}U}
\usepackage[european,siunitx]{circuitikz} % schémas électriques
\pgfplotsset{compat=1.18}
\usepgfplotslibrary{fillbetween,groupplots} % aires entre courbes (calcul intégral)
\usepackage{enumitem}
\usepackage{fancyhdr}
% [most] charge toutes les bibliothèques tcolorbox (listings, minted, raster,
% documentation...) et gonfle inutilement la mémoire TeX sur un document long
% (~300 boîtes) : on ne charge que ce qui est réellement utilisé.
\usepackage[skins,breakable]{tcolorbox}
\usepackage{graphicx}
\usepackage{colortbl}
\usepackage{booktabs}
\usepackage{tabularx}
\usepackage{diagbox} % case d'en-tête diagonale des tableaux à double entrée
% Colonnes extensibles centrée / gauche / droite (C, L, R), souvent utilisées par les modèles
\newcolumntype{C}{>{\centering\arraybackslash}X}
\newcolumntype{L}{>{\raggedright\arraybackslash}X}
\newcolumntype{R}{>{\raggedleft\arraybackslash}X}
\usepackage{array}
\usepackage{multicol}
\usepackage{multirow}
\usepackage{siunitx}
% parse-numbers=false : accepte aussi \SI{2,5 \times 10^{-3}}{...}
\sisetup{locale=FR,output-decimal-marker={,},group-minimum-digits=4,parse-numbers=false}
\DeclareSIUnit\ohm{\ensuremath{\symup{\Omega}}} % U+2126 absent de la police
\DeclareSIUnit\division{div} % réglages d'oscilloscope (ms/div, V/div)
\DeclareSIUnit\tr{tr} % tours (vitesse de rotation, ex. \SI{1500}{\tr\per\minute})
\DeclareSIUnit\molar{M} % concentration molaire (ex. \SI{3}{\molar}, \SI{1}{\micro\molar})
\DeclareSIUnit\an{an} % année (le modèle confond parfois avec \year, un compteur TeX) — ex. \SI{5}{\kilo\meter\per\an}
\DeclareSIUnit\FCFA{FCFA} % franc CFA (ex. \SI{500}{\FCFA\per\kilo\gram})
\DeclareSIUnit\degreeBrix{\textdegree Brix} % teneur en sucre d'un sirop/jus (réfractométrie)
\DeclareSIUnit\month{mois} % le modèle utilise parfois \month, un compteur TeX, comme unité de durée
\usepackage{qrcode}
% \degree : commande fréquemment utilisée par les modèles pour un angle ou une
% température en dehors de \SI{}{} (non fournie par les paquets chargés).
\providecommand{\degree}{^{\circ}}
% \qed (fin de démonstration), utilisé par les modèles sans amsthm
\providecommand{\qed}{\hfill$\square$}
% \barre : double barre des valeurs interdites dans un tableau de variations (tkz-tab)
\providecommand{\barre}{\ensuremath{\|}}
% environnement proof (amsthm non chargé) utilisé par les modèles
\ifdefined\proof\else\newenvironment{proof}[1][Démonstration]{\par\noindent\textit{#1.}\ }{\qed\par}\fi
\usepackage{lastpage}
\usepackage{hyperref}
\hypersetup{hidelinks}

% ============================================================
% Palette « Pop » OnBuch+ — mêmes hex que la classe P (plus_ui.dart)
% ============================================================
\definecolor{poporange}{HTML}{F59321}
\definecolor{popdark}{HTML}{E07A0C}
\definecolor{popcream}{HTML}{FFF6E8}
\definecolor{popink}{HTML}{1C1714}
\definecolor{poppurple}{HTML}{7A5AE0}
\definecolor{popgreen}{HTML}{2E9E6A}
\definecolor{poppink}{HTML}{E0457B}
\definecolor{popblue}{HTML}{2F7DE1}
\definecolor{popgold}{HTML}{C98A12}
\definecolor{popmuted}{HTML}{8A7F76}
\definecolor{popline}{HTML}{EADFD2}
\definecolor{popgray}{HTML}{8A7F76}
\colorlet{popgrey}{popgray}
% Alias tolérants (le modèle nomme parfois les couleurs différemment)
\colorlet{popwhite}{white}
\colorlet{popblack}{popink}
\colorlet{popred}{poppink}
\colorlet{popyellow}{popgold}
% Teintes claires (fonds de tableaux/encadrés) — le modèle nomme parfois
% les couleurs avec un suffixe L (ex. popblueL) sans les avoir définies.
\colorlet{poporangeL}{poporange!20}
\colorlet{popdarkL}{popdark!20}
\colorlet{poppurpleL}{poppurple!20}
\colorlet{popgreenL}{popgreen!20}
\colorlet{poppinkL}{poppink!20}
\colorlet{popblueL}{popblue!20}
\colorlet{popgoldL}{popgold!20}
\colorlet{popredL}{popred!20}
\colorlet{popgrayL}{popgray!20}
% Teintes claires pour fonds de tableaux / zones
\colorlet{popblueL}{popblue!10!white}
\colorlet{poporangeL}{poporange!14!white}
\colorlet{popgreenL}{popgreen!12!white}
\colorlet{poppurpleL}{poppurple!12!white}
\colorlet{poppinkL}{poppink!12!white}
\colorlet{popgoldL}{popgold!14!white}

\pagecolor{popcream}

% ============================================================
% Typographie
% ============================================================
\defaultfontfeatures{Ligatures=TeX}
\setmainfont{PlusJakartaSans-Medium.ttf}[Path=../../../fonts/,BoldFont=PlusJakartaSans-Bold.ttf]
% Symboles, API et emojis absents de la police principale : repli sur DejaVu Sans (caractères actifs)
\newfontface\symfont{DejaVuSans.ttf}[Path=../../../fonts/]
\makeatletter
\newcommand\symactive[1]{\catcode#1=13 \begingroup\lccode`\~=#1 \lowercase{\endgroup\def~}{{\symfont\char#1}}}
\newcommand\symblank[1]{\catcode#1=13 \begingroup\lccode`\~=#1 \lowercase{\endgroup\def~}{}}
\newcommand\symhyph[1]{\catcode#1=13 \begingroup\lccode`\~=#1 \lowercase{\endgroup\def~}{\mbox{-}}}
\newcount\symc
\newcommand\symrange[2]{\symc=#1 \@whilenum\symc<#2 \do{\expandafter\symactive\expandafter{\the\symc}\advance\symc by 1 }}
\newcommand\symblankrange[2]{\symc=#1 \@whilenum\symc<#2 \do{\expandafter\symblank\expandafter{\the\symc}\advance\symc by 1 }}
\makeatother
\symrange{"0250}{"02FF}\symactive{"2713}\symactive{"2714}\symactive{"2717}\symactive{"2718}\symactive{"2610}\symactive{"2611}\symactive{"2612}
\symhyph{"2011}\symhyph{"2E1A}\symblankrange{"1F300}{"1FAFF}\symblank{"FE0F}
% Maths en sans-serif (Fira Math) avec chiffres et lettres droites de Plus
% Jakarta, pour que les formules se fondent dans le texte.
\usepackage[math-style=ISO,bold-style=ISO]{unicode-math}
\setmathfont{FiraMath-Regular.otf}[Path=../../../fonts/]
\setmathfont{PlusJakartaSans-Medium.ttf}[Path=../../../fonts/,range={up/num,up/{Latin,latin},\mathup}]
\usepackage{icomma} % virgule décimale sans espace en mode math
% Caractères mathématiques Unicode absents des polices de texte (Plus Jakarta,
% Archivo Black) : rendus par leur équivalent mathématique, en texte comme en
% formule — sinon ils s'affichent en carré vide (ex. « Arithmétique dans ℤ »).
\usepackage{newunicodechar}
\newunicodechar{ℝ}{\ensuremath{\mathbb{R}}}
\newunicodechar{ℤ}{\ensuremath{\mathbb{Z}}}
\newunicodechar{ℂ}{\ensuremath{\mathbb{C}}}
\newunicodechar{ℕ}{\ensuremath{\mathbb{N}}}
\newunicodechar{ℚ}{\ensuremath{\mathbb{Q}}}
\newunicodechar{𝔻}{\ensuremath{\mathbb{D}}}
\newunicodechar{α}{\ensuremath{\alpha}}
\newunicodechar{β}{\ensuremath{\beta}}
\newunicodechar{γ}{\ensuremath{\gamma}}
\newunicodechar{δ}{\ensuremath{\delta}}
\newunicodechar{ε}{\ensuremath{\varepsilon}}
\newunicodechar{θ}{\ensuremath{\theta}}
\newunicodechar{λ}{\ensuremath{\lambda}}
\newunicodechar{μ}{\ensuremath{\mu}}
\newunicodechar{σ}{\ensuremath{\sigma}}
\newunicodechar{τ}{\ensuremath{\tau}}
\newunicodechar{φ}{\ensuremath{\varphi}}
\newunicodechar{ψ}{\ensuremath{\psi}}
\newunicodechar{ω}{\ensuremath{\omega}}
\newunicodechar{Σ}{\ensuremath{\Sigma}}
% majuscules produites par \MakeUppercase dans les titres (α-aminés -> Α)
\newunicodechar{Α}{\ensuremath{\alpha}}
\newunicodechar{Β}{\ensuremath{\beta}}
\newunicodechar{Γ}{\ensuremath{\Gamma}}
\newunicodechar{Δ}{\ensuremath{\Delta}}
\newunicodechar{Ε}{\ensuremath{\varepsilon}}
\newunicodechar{Θ}{\ensuremath{\Theta}}
\newunicodechar{Λ}{\ensuremath{\Lambda}}
\newunicodechar{Μ}{\ensuremath{\mu}}
\newunicodechar{Τ}{\ensuremath{\tau}}
\newunicodechar{Φ}{\ensuremath{\Phi}}
\newunicodechar{Ψ}{\ensuremath{\Psi}}
\newunicodechar{Ω}{\ensuremath{\Omega}}
\newunicodechar{↦}{\ensuremath{\mapsto}}
\newunicodechar{∈}{\ensuremath{\in}}
\newunicodechar{∉}{\ensuremath{\notin}}
\newunicodechar{∧}{\ensuremath{\wedge}}
\newunicodechar{⇔}{\ensuremath{\Leftrightarrow}}
\newunicodechar{⟺}{\ensuremath{\Longleftrightarrow}}
\newunicodechar{⇒}{\ensuremath{\Rightarrow}}
\newunicodechar{⟹}{\ensuremath{\Longrightarrow}}
\newunicodechar{∘}{\ensuremath{\circ}}
\newunicodechar{∪}{\ensuremath{\cup}}
\newunicodechar{∩}{\ensuremath{\cap}}
\newunicodechar{≡}{\ensuremath{\equiv}}
\newunicodechar{⊂}{\ensuremath{\subset}}
\newunicodechar{⊥}{\ensuremath{\perp}}
\newunicodechar{∃}{\ensuremath{\exists}}
\newunicodechar{∀}{\ensuremath{\forall}}
\newunicodechar{∛}{\ensuremath{\sqrt[3]{\,}}}
\newunicodechar{∜}{\ensuremath{\sqrt[4]{\,}}}
\newunicodechar{⃗}{\ensuremath{{}^{\to}}}
\newunicodechar{ⁿ}{\ifmmode^{n}\else\textsuperscript{\ensuremath{n}}\fi}
\newunicodechar{ˣ}{\ifmmode^{x}\else\textsuperscript{\ensuremath{x}}\fi}
\newunicodechar{ᵇ}{\ifmmode^{b}\else\textsuperscript{\ensuremath{b}}\fi}
\newunicodechar{ʳ}{\ifmmode^{r}\else\textsuperscript{\ensuremath{r}}\fi}
\newunicodechar{ᵘ}{\ifmmode^{u}\else\textsuperscript{\ensuremath{u}}\fi}
\newunicodechar{ʸ}{\ifmmode^{y}\else\textsuperscript{\ensuremath{y}}\fi}
\newunicodechar{ᵃ}{\ifmmode^{a}\else\textsuperscript{\ensuremath{a}}\fi}
\newunicodechar{ᵉ}{\ifmmode^{e}\else\textsuperscript{\ensuremath{e}}\fi}
\newunicodechar{ᵏ}{\ifmmode^{k}\else\textsuperscript{\ensuremath{k}}\fi}
\newunicodechar{ᵖ}{\ifmmode^{p}\else\textsuperscript{\ensuremath{p}}\fi}
\newunicodechar{ᴺ}{\ifmmode^{N}\else\textsuperscript{\ensuremath{N}}\fi}
\newunicodechar{ᶜ}{\ifmmode^{c}\else\textsuperscript{\ensuremath{c}}\fi}
\newunicodechar{ʰ}{\ifmmode^{h}\else\textsuperscript{\ensuremath{h}}\fi}
\newunicodechar{ᵠ}{\ifmmode^{q}\else\textsuperscript{\ensuremath{q}}\fi}
\newunicodechar{ₙ}{\ifmmode_{n}\else\textsubscript{\ensuremath{n}}\fi}
\newunicodechar{ₐ}{\ifmmode_{a}\else\textsubscript{\ensuremath{a}}\fi}
\newunicodechar{ᵢ}{\ifmmode_{i}\else\textsubscript{\ensuremath{i}}\fi}
\newunicodechar{ᵣ}{\ifmmode_{r}\else\textsubscript{\ensuremath{r}}\fi}
\newunicodechar{ₖ}{\ifmmode_{k}\else\textsubscript{\ensuremath{k}}\fi}
\newunicodechar{ᵦ}{\ifmmode_{\beta}\else\textsubscript{\ensuremath{\beta}}\fi}
\newunicodechar{ₓ}{\ifmmode_{x}\else\textsubscript{\ensuremath{x}}\fi}
\newfontfamily\popdisplay{ArchivoBlack-Regular.ttf}[Path=../../../fonts/]
\newfontfamily\poplogo{SpaceGrotesk-Bold.ttf}[Path=../../../fonts/]
\newfontfamily\popsemi{PlusJakartaSans-SemiBold.ttf}[Path=../../../fonts/]
\newfontfamily\popxbold{PlusJakartaSans-ExtraBold.ttf}[Path=../../../fonts/]
\newfontfamily\popxboldls{PlusJakartaSans-ExtraBold.ttf}[Path=../../../fonts/,LetterSpace=3.0]

\setlist[enumerate]{leftmargin=1.7em,itemsep=5pt,topsep=4pt}
\setlist[itemize]{leftmargin=1.7em,itemsep=4pt,topsep=4pt,label={\color{poporange}\textbullet}}
\setlength{\parindent}{0pt}
\setlength{\parskip}{5pt}
\renewcommand{\arraystretch}{1.25}

% pgfplots : style Pop par défaut
\pgfplotsset{
  every axis/.append style={axis line style={popink,line width=1pt},tick style={popink},
    grid style={popline,line width=0.5pt},label style={font=\small},tick label style={font=\footnotesize}},
  popcurve/.style={poporange,line width=1.8pt,no markers,smooth},
}
% Même style disponible dans un \draw TikZ simple (hors pgfplots)
\tikzset{popcurve/.style={poporange,line width=1.8pt}}
\tikzset{popgrid/.style={popline,very thin}}
\colorlet{popcurve}{poporange} % le nom du style est parfois utilisé comme couleur

% ============================================================
% Pill / squiggle
% ============================================================
\newcommand{\poppill}[3]{%
  \tikz[baseline=-0.55ex]\node[draw=popink,line width=1.2pt,rounded corners=2.6pt,%
    fill=#2,inner xsep=6.5pt,inner ysep=3pt]{%
    \popxboldls\fontsize{7}{7}\selectfont\color{#3}\MakeUppercase{#1}};%
}
\newcommand{\popsquiggle}[1]{%
  \tikz[baseline=-0.4ex]{
    \draw[#1,line width=2.6pt,line cap=round]
      plot [smooth] coordinates {(0,0.09) (0.34,0.01) (0.68,0.21) (1.02,0.01) (1.36,0.10)};
  }%
}
% Mot-clé surligné (marqueur orange sous le texte)
\newcommand{\cle}[1]{{\popxbold\color{popdark}#1}}

% ============================================================
% En-tête / pied
% ============================================================
\pagestyle{fancy}
\fancyhf{}
\renewcommand{\headrulewidth}{0pt}
\renewcommand{\footrulewidth}{0pt}
\lhead{\raisebox{-0.15ex}{\poplogo\fontsize{13}{13}\selectfont\color{popink}OnBuch\color{poporange}+}}
\rhead{\poppill{\DOCMATIERE\ \textbullet\ \DOCNIVEAU\ \textbullet\ Leçon \DOCLECON}{white}{popink}}
\lfoot{\popxbold\fontsize{7.6}{7.6}\selectfont\color{popmuted}\MakeUppercase{Cours signé OnBuch+}\ \textbullet\ {\popsemi\DOCTITRE}}
\rfoot{\tikz[baseline=-0.6ex]\node[rounded corners=8.5pt,fill=popink,minimum height=17pt,minimum width=17pt,inner xsep=5pt,inner sep=0]{\popxbold\fontsize{8}{8}\selectfont\color{popcream}\thepage\,/\,\pageref*{LastPage}};}

% ============================================================
% Titres
% ============================================================
\newcommand{\chaptitre}[2]{%
  \begin{tcolorbox}[enhanced,colback=poporange,colframe=popink,boxrule=2.6pt,arc=11pt,
      drop shadow=popink,left=15pt,right=15pt,top=12pt,bottom=11pt,before skip=3pt,after skip=18pt]
    \poppill{#2}{popink}{popcream}\par\vspace{9pt}
    {\raggedright\hyphenpenalty=10000\exhyphenpenalty=10000\popdisplay\fontsize{22}{24}\selectfont\color{popcream}\MakeUppercase{#1}\par}
  \end{tcolorbox}
}
% Section numérotée : pastille orange avec le numéro + titre Archivo
\newcounter{popsec}
\newcommand{\coursec}[1]{%
  \stepcounter{popsec}\setcounter{popsub}{0}%
  \par\vspace{16pt}\noindent
  \begin{minipage}{\linewidth}
  \tikz[baseline=-0.7ex]\node[draw=popink,line width=1.6pt,fill=poporange,rounded corners=4pt,minimum size=22pt,inner sep=0]{\popdisplay\fontsize{12}{12}\selectfont\color{popcream}\thepopsec};%
  \hspace{8pt}{\popdisplay\fontsize{15}{17}\selectfont\color{popink}\MakeUppercase{#1}}\par
  \vspace{4pt}\noindent{\color{poporange}\rule{36pt}{3.6pt}}
  \end{minipage}\par\nopagebreak\vspace{8pt}
}
\newcounter{popsub}
\newcommand{\courssub}[1]{%
  \stepcounter{popsub}%
  \par\vspace{9pt}\noindent{\popxbold\fontsize{11.5}{13}\selectfont\color{popink}{\color{poporange}\thepopsec.\thepopsub}\hspace{6pt}#1}\par\nopagebreak\vspace{4pt}
}

% ============================================================
% Boîtes pédagogiques — même recette : fond blanc, bordure épaisse colorée,
% ombre dure encre, badge plein. Titre optionnel [#1].
% ============================================================
\tcbset{popbase/.style={breakable,enhanced,colback=white,boxrule=2.3pt,arc=9pt,
  shadow={3pt}{-3pt}{0pt}{popink},left=14pt,right=14pt,top=11pt,bottom=11pt,before skip=11pt,after skip=15pt,
  fontupper=\popsemi\fontsize{10.6}{14.5}\selectfont\color{popink},
  fontlower=\popsemi\fontsize{10.6}{14.5}\selectfont\color{popink}}}
% Coche dessinée (le glyphe ✓ n'existe pas dans les polices du document)
\newcommand{\popcheck}[1]{\tikz[baseline=-0.5ex]\draw[#1,line width=1.6pt,line cap=round,line join=round](-0.13,0.0)--(-0.04,-0.09)--(0.13,0.1);}
\providecommand{\coche}{\popcheck{popgreen}}
\providecommand{\resultat}[1]{\par\smallskip\noindent\fcolorbox{popgreen}{popgreenL}{\parbox{\dimexpr\linewidth-2\fboxsep-2\fboxrule\relax}{\textbf{Résultat.} #1}}\par\smallskip}
\newcommand{\popboxhead}[3]{\poppill{#1}{#2}{white}\ifx\relax#3\relax\else\hspace{6pt}{\popxbold\fontsize{10.6}{12}\selectfont\color{popink}#3}\fi\par\vspace{7pt}}

\newtcolorbox{aretenir}[1][]{popbase,colframe=popgreen,before upper={\popboxhead{À retenir}{popgreen}{#1}}}
% Texte en espagnol (document de lecture, dialogue, poème, extrait) : cadre
% d'étude. Un titre optionnel (source, auteur) s'affiche dans l'en-tête.
\newtcolorbox{textebox}[1][]{popbase,colframe=popblue,before upper={\popboxhead{Texto}{popblue}{#1}\begin{otherlanguage}{spanish}},after upper={\end{otherlanguage}}}
% Marques ✓ / ✗ (forme correcte / fautive) souvent utilisées par les modèles
\providecommand{\cross}{\textcolor{poppink}{\ensuremath{\times}}}
\providecommand{\xmark}{\cross}\providecommand{\crossmark}{\cross}\providecommand{\cmark}{\ensuremath{\checkmark}}
% Ligne de réponse / justification à compléter (inventée par les modèles)
\providecommand{\justification}[1]{\par\noindent\textit{Justification~:} \dotfill\par}
% \textipa (package tipa non chargé) : la police ne couvre pas l'API -> simple passage
\providecommand{\textipa}[1]{#1}
% Soulignement (ulem non chargé) : \uline, \ul
\providecommand{\uline}[1]{\underline{#1}}\providecommand{\ul}[1]{\underline{#1}}
% \glossaire{\item ...} inventée par les modèles : liste de glossaire
\providecommand{\glossaire}[1]{\begin{itemize}#1\end{itemize}}
% \alert (beamer) inventée par les modèles : texte mis en évidence
\providecommand{\alert}[1]{\textbf{\textcolor{poppink}{#1}}}
% variantes ASCII d'ordinaux écrites en macro
\providecommand{\iere}{\textsuperscript{re}}\providecommand{\ieme}{\textsuperscript{e}}\providecommand{\ere}{\textsuperscript{re}}\providecommand{\eme}{\textsuperscript{e}}\providecommand{\er}{\textsuperscript{er}}\providecommand{\ie}{\textsuperscript{e}}
% Ordinaux français écrits en macro par les modèles : 1\ière, 3\ième
\newcommand{\ière}{\textsuperscript{re}}\newcommand{\ième}{\textsuperscript{e}}
% \exemple (inventée par les modèles) : étiquette « Exemple : »
\providecommand{\exemple}{\textbf{Exemple~:}\ }
% Mot ou phrase en espagnol à l'intérieur d'un paragraphe français
\newcommand{\esp}[1]{\ifmmode\text{\textit{#1}}\else\begingroup\selectlanguage{spanish}\itshape #1\endgroup\fi}
\newtcolorbox{exemplebox}[1][]{popbase,colframe=poporange,before upper={\popboxhead{Exemple}{poporange}{#1}}}
\newtcolorbox{definition}[1][]{popbase,colframe=popblue,before upper={\popboxhead{Définition}{popblue}{#1}}}
\newtcolorbox{propriete}[1][]{popbase,colframe=poppurple,before upper={\popboxhead{Propriété}{poppurple}{#1}}}
\newtcolorbox{methode}[1][]{popbase,colframe=popgold,before upper={\popboxhead{Méthode}{popgold}{#1}}}
\newtcolorbox{attention}[1][]{popbase,colframe=poppink,before upper={\popboxhead{Attention}{poppink}{#1}}}
\newtcolorbox{experience}[1][]{popbase,colframe=popblue,colback=popblueL,before upper={\popboxhead{Expérience}{popblue}{#1}}}
% « Le savais-tu ? » : carte violette pleine (contexte, culture, Cameroun)
\newtcolorbox{savaistu}[1][]{popbase,colback=poppurple,colframe=popink,
  fontupper=\popsemi\fontsize{10.4}{14.2}\selectfont\color{white},
  before upper={\poppill{Le savais-tu\,?}{white}{poppurple}\ifx\relax#1\relax\else\hspace{6pt}{\popxbold\color{white}#1}\fi\par\vspace{7pt}}}
% Exercice résolu : énoncé en haut, \tcblower puis la solution sur fond crème
\newcounter{popexo}\newcounter{popexr}
\newtcolorbox{exoresolu}[1][]{popbase,colframe=poporange,colbacklower=popcream,
  segmentation style={popink,line width=1.2pt,dashed},
  before upper={\stepcounter{popexr}\popboxhead{Exercice résolu \thepopexr}{poporange}{#1}},
  before lower={\poppill{Solution}{popink}{popcream}\par\vspace{6pt}}}
% Exercice (non résolu en place) — la correction est dans la partie Corrigés
\newtcolorbox{exercice}[1][]{popbase,colframe=popink,
  before upper={\stepcounter{popexo}\popboxhead{Exercice \thepopexo}{popink}{#1}}}
% Corrigé
\newtcolorbox{corrige}[1][]{popbase,colframe=popgreen,colback=popgreenL,
  before upper={\popboxhead{Corrigé}{popgreen}{#1}}}
% Objectifs / prérequis / bilan
\newtcolorbox{objectifs}{popbase,colframe=popink,colback=poporangeL,
  before upper={\popboxhead{Objectifs de la leçon}{popink}{}}}
\newtcolorbox{prerequis}{popbase,colframe=popink,colback=popblueL,
  before upper={\popboxhead{Prérequis}{popblue}{}}}
\newtcolorbox{bilan}{popbase,colframe=popink,colback=white,boxrule=2.8pt,
  before upper={\popboxhead{Fiche bilan}{poporange}{L'essentiel à connaître pour l'examen}}}
% Figure encadrée avec légende
\newtcolorbox{popfigure}[1][]{enhanced,breakable=false,colback=white,colframe=popline,boxrule=1.4pt,arc=8pt,
  left=10pt,right=10pt,top=10pt,bottom=8pt,before skip=10pt,after skip=12pt,halign=center,
  lower separated=false,halign lower=center,
  fontlower=\popsemi\fontsize{9}{11.5}\selectfont\color{popmuted},
  #1}
\newcommand{\legende}[1]{\par\vspace{6pt}{\popsemi\fontsize{9}{11.5}\selectfont\color{popmuted}#1}}

% ============================================================
% Signature OnBuch+
% ============================================================
\newcommand{\poplogoinline}[1]{{\poplogo\fontsize{#1}{#1}\selectfont\color{popink}OnBuch\color{poporange}+}}
\newcommand{\signatureonbuch}{%
  \begin{tcolorbox}[enhanced,colback=popink,colframe=popink,boxrule=2.2pt,arc=10pt,
      drop shadow=poporange,left=14pt,right=14pt,top=10pt,bottom=10pt,before skip=0pt,after skip=0pt]
    \tikz[baseline=-0.7ex]\node[circle,fill=popgreen,draw=popcream,line width=1.3pt,minimum size=22pt,inner sep=0]{\popcheck{white}};%
    \hspace{9pt}\begin{minipage}[c]{0.62\linewidth}
      {\popxbold\fontsize{12}{13}\selectfont\color{popcream}Signé\ \textbullet\ {\poplogo OnBuch\color{poporange}+}}\par\vspace{2pt}
      {\raggedright\popsemi\fontsize{8}{10.5}\selectfont\color{popline}\MakeUppercase{Équipe pédagogique OnBuch+}\\\MakeUppercase{Conforme au programme officiel MINESEC}\par}
    \end{minipage}\hfill
    {\poplogo\fontsize{22}{22}\selectfont\color{popcream}OnBuch\color{poporange}+}
  \end{tcolorbox}%
}

% ============================================================
% Page de garde
% #1 titre · #2 matière · #3 niveau · #4 module · #5 n° leçon · #6 durée ·
% #7 description / objectifs (texte ou itemize)
% ============================================================
\newcommand{\pagegarde}[7]{%
  \thispagestyle{empty}
  \newgeometry{a4paper,margin=1.7cm,top=1.6cm,bottom=1.4cm}%
  \begin{tikzpicture}[remember picture,overlay]
    \fill[poporange!18] ([xshift=-3.2cm,yshift=-3.6cm]current page.north east) circle (4.6cm);
    \fill[poppurple!14] ([xshift=2.4cm,yshift=7.6cm]current page.south west) circle (3.8cm);
  \end{tikzpicture}%
  {\poplogo\fontsize{30}{30}\selectfont\color{popink}OnBuch\color{poporange}+}\hfill
  \raisebox{0.6ex}{\poppill{Cours signé \textbullet\ Premium}{popink}{popcream}}\par
  \vspace{1pt}\popsquiggle{poppurple}\par
  \vspace{8pt}
  \begin{tcolorbox}[enhanced,colback=poporange,colframe=popink,boxrule=2.8pt,arc=13pt,
      drop shadow=popink,left=18pt,right=18pt,top=12pt,bottom=13pt,after skip=10pt]
    \poppill{#2 \textbullet\ #3}{popink}{popcream}\hspace{5pt}\poppill{#4}{white}{popink}\par\vspace{8pt}
    {\popdisplay\fontsize{14}{14}\selectfont\color{popink}LEÇON #5}\par\vspace{4pt}
    {\raggedright\hyphenpenalty=10000\exhyphenpenalty=10000\popdisplay\fontsize{25}{27}\selectfont\color{popcream}\MakeUppercase{#1}\par}
    \vspace{8pt}
    {\popxbold\fontsize{9.5}{9.5}\selectfont\color{popink}\MakeUppercase{Durée conseillée : #6}}
  \end{tcolorbox}
  \begin{tcolorbox}[enhanced,colback=white,colframe=popink,boxrule=2.2pt,arc=10pt,
      drop shadow=popink,left=16pt,right=16pt,top=10pt,bottom=10pt,after skip=8pt]
    \poppill{Dans cette leçon}{poppurple}{white}\par\vspace{5pt}
    \popsemi\fontsize{9.3}{12.6}\selectfont\color{popink}#7
  \end{tcolorbox}
  \noindent\begin{minipage}[t]{0.487\linewidth}
    \begin{tcolorbox}[enhanced,colback=popgreen,colframe=popink,boxrule=2.2pt,arc=10pt,
        drop shadow=popink,left=11pt,right=11pt,top=10pt,bottom=10pt,height=3.1cm,valign=center]
      \tcbox[colback=white,colframe=popink,boxrule=1.3pt,arc=5pt,boxsep=0pt,nobeforeafter,
        left=3pt,right=3pt,top=3pt,bottom=3pt,valign=center]{\qrcode[height=1.9cm]{https://play.google.com/store/apps/details?id=com.luvvix.onbuch&hl=fr}}%
      \hspace{8pt}\begin{minipage}[c]{0.5\linewidth}
        {\popdisplay\fontsize{11}{12.5}\selectfont\color{white}\MakeUppercase{Télécharge OnBuch}}\par\vspace{4pt}
        {\raggedright\popsemi\fontsize{8.4}{11}\selectfont\color{white}Gratuit \textbullet\ Google Play\\Tuteur IA, annales, concours, résultats\par}
      \end{minipage}
    \end{tcolorbox}
  \end{minipage}\hfill
  \begin{minipage}[t]{0.487\linewidth}
    \begin{tcolorbox}[enhanced,colback=poppurple,colframe=popink,boxrule=2.2pt,arc=10pt,
        drop shadow=popink,left=11pt,right=11pt,top=10pt,bottom=10pt,height=3.1cm,valign=center]
      \tikz[baseline=-0.7ex]\node[circle,fill=white,draw=popink,line width=1.3pt,minimum size=1.6cm,inner sep=0]{\popdisplay\fontsize{24}{24}\selectfont\color{poppurple}+};%
      \hspace{8pt}\begin{minipage}[c]{0.6\linewidth}
        {\popdisplay\fontsize{11}{12.5}\selectfont\color{white}\MakeUppercase{Passe à OnBuch+}}\par\vspace{4pt}
        {\raggedright\popsemi\fontsize{8.4}{11}\selectfont\color{white}Tous les cours signés, Tuteur IA illimité, fiches PDF — onglet OnBuch+ de l'app\par}
      \end{minipage}
    \end{tcolorbox}
  \end{minipage}
  \vspace{8pt}
  \popcontenu
  \vfill
  \signatureonbuch
  \restoregeometry
  \newpage
}

% Rangée « Ce cours contient » (page de garde)
\newcommand{\poptile}[3]{% #1 couleur #2 gros texte #3 légende
  \begin{tcolorbox}[enhanced,colback=#1,colframe=popink,boxrule=2pt,arc=9pt,shadow={2.5pt}{-2.5pt}{0pt}{popink},
      left=6pt,right=6pt,top=9pt,bottom=9pt,halign=center,valign=center,height=1.75cm,width=\linewidth,nobeforeafter]
    {\popdisplay\fontsize{12.5}{13}\selectfont\color{white}\MakeUppercase{#2}}\par\vspace{3pt}
    {\centering\popsemi\fontsize{7.8}{9.5}\selectfont\color{white}#3\par}
  \end{tcolorbox}}
\newcommand{\popcontenu}{%
  \par\noindent{\popxboldls\fontsize{7.5}{7.5}\selectfont\color{popmuted}CE COURS SIGNÉ CONTIENT}\par\vspace{7pt}
  \noindent\begin{minipage}[t]{0.235\linewidth}\poptile{popblue}{Cours}{complet, illustré, pas à pas}\end{minipage}\hfill
  \begin{minipage}[t]{0.235\linewidth}\poptile{poporange}{Méthodes}{et exercices résolus}\end{minipage}\hfill
  \begin{minipage}[t]{0.235\linewidth}\poptile{poppink}{Exercices}{type examen + corrigés}\end{minipage}\hfill
  \begin{minipage}[t]{0.235\linewidth}\poptile{popgreen}{Fiche bilan}{l'essentiel à retenir}\end{minipage}\par}

% Sommaire
\newtcolorbox{sommaire}{enhanced,colback=white,colframe=popink,boxrule=2.2pt,arc=10pt,
  drop shadow=popink,left=18pt,right=18pt,top=13pt,bottom=13pt,before skip=6pt,after skip=20pt,
  before upper={\poppill{Sommaire}{poppurple}{white}\par\vspace{9pt}\popsemi\fontsize{10.6}{14.5}\selectfont\color{popink}}}

% Signature de fin de cours — poussée tout en bas de la dernière page
\newcommand{\findecours}{%
  \par\vfill\nopagebreak
  \begin{center}\popsquiggle{poporange}\end{center}\vspace{4pt}
  \signatureonbuch
}
\AtBeginDocument{\def\llbracket{\mathopen{[\mkern-3.5mu[}}\def\rrbracket{\mathclose{]\mkern-3.5mu]}}} % intervalles d'entiers (glyphe ⟦ absent de la police)
% Glyphes absents de Fira Math : redéfinis à partir de glyphes disponibles
\AtBeginDocument{%
  \def\setminus{\mathbin{\backslash}}\let\smallsetminus\setminus
  \def\vdots{\mathord{\vbox{\baselineskip3.2pt\lineskiplimit0pt\kern4pt\hbox{.}\hbox{.}\hbox{.}}}}%
  \def\ddots{\mathinner{\mkern1mu\raise6.5pt\hbox{.}\mkern2mu\raise3.6pt\hbox{.}\mkern2mu\raise0.7pt\hbox{.}\mkern1mu}}%
  \def\triangle{\mathord{\scalebox{0.9}{$\symup{\Delta}$}}}%
  \def\nabla{\mathord{\raisebox{\depth}{\scalebox{1}[-1]{$\symup{\Delta}$}}}}%
  \def\checkmark{\popcheck{popdark}}%
  \def\star{\mathbin{\ast}}\def\bigstar{\mathord{\ast}}%
  \def\diamond{\mathbin{\scalebox{0.6}{$\Diamond$}}}%
  \def\bigcup{\mathop{\mathchoice{\raisebox{-0.3em}{\scalebox{1.7}{$\cup$}}}{\scalebox{1.25}{$\cup$}}{\cup}{\cup}}\displaylimits}%
  \def\bigcap{\mathop{\mathchoice{\raisebox{-0.3em}{\scalebox{1.7}{$\cap$}}}{\scalebox{1.25}{$\cap$}}{\cap}{\cap}}\displaylimits}%
  \def\lesseqqgtr{\mathrel{\lessgtr}}\def\bigoplus{\mathop{\oplus}\displaylimits}%
}
\providecommand{\ssi}{si et seulement si} % abréviation fréquente
\AtBeginDocument{\providecommand{\wideparen}{\overparen}} % arc de cercle

\begin{document}

\pagegarde{Lectura y comentario : « El mundo laboral hoy » ; ayudas, subvenciones, acuerdos bilaterales y multilaterales}{Espagnol}{Tle A}{Módulo 4 : Activités économiques et monde du travail}{E16}{2 h}{Cette leçon de lecture-commentaire propose l'étude d'un texte original sur le monde du travail actuel et l'aide au développement. Elle fournit le lexique thématique indispensable (subventions, accords bilatéraux et multilatéraux, coopération, ONG) et la méthode complète du commentaire de texte exigée au Baccalauréat.
\vspace{4pt}
\begin{multicols}{2}\small
\begin{itemize}
\item À la fin de cette leçon, je sais lire et comprendre globalement un texte économique en espagnol
\item je sais identifier les idées directrices d'un texte et les hiérarchiser
\item je sais employer correctement le lexique de l'aide au développement (subvención, ayuda, acuerdo, inversión, cooperación, ONG)
\item je sais distinguer un acuerdo bilateral d'un acuerdo multilateral et donner des exemples
\item je sais répondre à des questions de compréhension en justifiant par le texte
\item je sais construire un commentaire de texte structuré (introduction, développement, conclusion)
\end{itemize}\end{multicols}}

\begin{sommaire}
\begin{multicols}{2}
\begin{enumerate}
\item Texte support : « El mundo laboral hoy »
\item Compréhension du texte
\item Lexique thématique : l'aide au développement et le monde du travail
\item Méthode du commentaire de texte
\item Commentaire modèle guidé
\item Production guidée
\item Activité d'intégration
\item Méthodes et exercices résolus
\item Exercices
\item Corrigés des exercices
\item Fiche bilan
\end{enumerate}
\end{multicols}
\end{sommaire}

\begin{objectifs}
\begin{itemize}
\item À la fin de cette leçon, je sais lire et comprendre globalement un texte économique en espagnol
\item je sais identifier les idées directrices d'un texte et les hiérarchiser
\item je sais employer correctement le lexique de l'aide au développement (subvención, ayuda, acuerdo, inversión, cooperación, ONG)
\item je sais distinguer un acuerdo bilateral d'un acuerdo multilateral et donner des exemples
\item je sais répondre à des questions de compréhension en justifiant par le texte
\item je sais construire un commentaire de texte structuré (introduction, développement, conclusion)
\item je sais exprimer mon opinion sur les enjeux du développement économique en Afrique
\item je sais rédiger une synthèse courte d'un texte économique
\end{itemize}
\end{objectifs}

\begin{prerequis}
\begin{itemize}
\item Lexique général du travail vu en Première (trabajo, empleo, empresa, trabajador)
\item Techniques de lecture : repérage du thème, du titre, des connecteurs logiques
\item Expression de l'opinion (pienso que, en mi opinión, me parece que + indicatif/subjonctif)
\item Présent de l'indicatif et connecteurs (sin embargo, además, por eso, aunque)
\item Notions de géographie économique : pays développés / pays en développement
\end{itemize}
\end{prerequis}

\newpage

\coursec{Texte support : « El mundo laboral hoy »}

Chaque matin à Yaoundé, des jeunes diplômés consultent les offres d'emploi sur leurs téléphones, assis dans les cybercafés du quartier Nlongkak ou dans un bendskin qui file vers le centre-ville. Au même moment, la radio annonce un nouvel accord de coopération entre le Cameroun et l'Espagne pour financer des projets agricoles dans la région du Centre. Derrière ces réalités quotidiennes se cachent des mots clés : \esp{subvención}, \esp{ayuda al desarrollo}, \esp{acuerdo bilateral}... Comprendre ces termes en espagnol, c'est pouvoir lire la presse hispanophone, débattre des enjeux du développement de l'Afrique et réussir l'épreuve de commentaire de texte au Baccalauréat.

\textbf{Problématique :} comment lire, comprendre et commenter un texte économique en espagnol portant sur le monde du travail et l'aide au développement ? Cette leçon vous donne les clés, étape par étape.

\courssub{Présentation du texte}

\begin{definition}[Fiche d'identité du texte]
\begin{itemize}
\item \textbf{Titre :} \esp{El mundo laboral hoy} (« Le monde du travail aujourd'hui »).
\item \textbf{Source :} Didáctica V, unité 4 (manuel d'espagnol de Terminale, programme MINESEC).
\item \textbf{Type de texte :} article de presse économique à visée argumentative et informative.
\item \textbf{Thème annoncé :} l'évolution du monde du travail en Afrique et les mécanismes de l'aide au développement (subventions, investissements, accords, ONG).
\item \textbf{Structure :} 4 paragraphes — (1) le travail aujourd'hui ; (2) l'aide au développement ; (3) les accords bilatéraux et multilatéraux ; (4) le rôle des ONG et l'avenir de l'Afrique.
\end{itemize}
\end{definition}

\courssub{Lecture du texte}

\textbf{Étape 1 — Lecture silencieuse (3 minutes).} Lisez le texte seul, crayon en main. Soulignez les mots que vous comprenez sans dictionnaire (les \cle{mots transparents}) et entourez les mots nouveaux.

\textbf{Étape 2 — Lecture à voix haute par l'enseignant.} Écoutez la prononciation : notez l'accent tonique de \esp{subvención}, \esp{inversión}, \esp{organización} (accent écrit sur la dernière syllabe, \emph{agudas}) et la prononciation de \esp{ONG} [o-ené-gé].

\begin{textebox}[El mundo laboral hoy — article de presse économique (texte original)]
El mundo del trabajo ha cambiado mucho en los últimos años. Hoy, muchos jóvenes africanos buscan un empleo digno, pero las empresas no siempre pueden contratarlos. Por eso, la ayuda al desarrollo sigue siendo necesaria.

Los países ricos conceden subvenciones y realizan inversiones en África: construyen escuelas, hospitales y carreteras. Estas ayudas se organizan mediante acuerdos bilaterales, entre dos países, como el acuerdo entre España y Camerún, o acuerdos multilaterales, entre varios países, como los de la Unión Europea o las Naciones Unidas.

Además, muchas ONG trabajan sobre el terreno para que la cooperación llegue a la población. Sin embargo, algunos economistas piensan que África no necesita solamente ayuda, sino también un comercio justo y más inversiones productivas.

El verdadero desarrollo debe crear empleo duradero para los jóvenes.
\end{textebox}

\textbf{Traduction du texte.} « Le monde du travail a beaucoup changé ces dernières années. Aujourd'hui, beaucoup de jeunes Africains cherchent un emploi digne, mais les entreprises ne peuvent pas toujours les embaucher. C'est pourquoi l'aide au développement reste nécessaire. Les pays riches accordent des subventions et réalisent des investissements en Afrique : ils construisent des écoles, des hôpitaux et des routes. Ces aides s'organisent au moyen d'accords bilatéraux, entre deux pays, comme l'accord entre l'Espagne et le Cameroun, ou d'accords multilatéraux, entre plusieurs pays, comme ceux de l'Union européenne ou des Nations unies. De plus, de nombreuses ONG travaillent sur le terrain pour que la coopération parvienne à la population. Cependant, certains économistes pensent que l'Afrique n'a pas seulement besoin d'aide, mais aussi d'un commerce équitable et de davantage d'investissements productifs. Le véritable développement doit créer un emploi durable pour les jeunes. »

\textbf{Glossaire des mots difficiles.}

\begin{center}
\begin{tabularx}{\linewidth}{|X|X|}\hline
\rowcolor{popblueL} \textbf{Español} & \textbf{Français} \\ \hline
\esp{un empleo digno} & un emploi digne, décent \\ \hline
\esp{contratar} & embaucher, recruter \\ \hline
\esp{conceder} & accorder, octroyer \\ \hline
\esp{mediante} & au moyen de, par le biais de \\ \hline
\esp{sobre el terreno} & sur le terrain \\ \hline
\esp{sin embargo} & cependant, pourtant \\ \hline
\esp{el comercio justo} & le commerce équitable \\ \hline
\esp{duradero, -a} & durable \\ \hline
\esp{sino} & mais (après une négation) \\ \hline
\end{tabularx}
\end{center}

\courssub{Repérage des mots transparents}

Avant tout dictionnaire, exploitez les \cle{mots transparents} : mots espagnols très proches du français, qui donnent immédiatement accès au sens global.

\begin{center}
\begin{tabularx}{\linewidth}{|X|X|X|}\hline
\rowcolor{popblueL} \textbf{Español} & \textbf{Français} & \textbf{Remarque} \\ \hline
\esp{cooperación} & coopération & \esp{-ción} = -tion \\ \hline
\esp{inversión} & investissement & attention au sens précis (faux ami) \\ \hline
\esp{multilateral} & multilatéral & même famille latine \\ \hline
\esp{organización} & organisation & \esp{z} se prononce [θ] en Espagne \\ \hline
\esp{hospital} & hôpital & le \esp{h} est muet \\ \hline
\esp{economista} & économiste & épicène : \esp{el/la economista} \\ \hline
\esp{productivo, -a} & productif, -ive & adjectif variable \\ \hline
\esp{población} & population & \esp{-ción} = -tion \\ \hline
\esp{desarrollo} & développement & racine \esp{desarrollar} \\ \hline
\end{tabularx}
\end{center}

\begin{attention}[Faux amis et pièges du texte]
\begin{itemize}
\item \esp{inversión} ne signifie pas « inversion » (renversement) dans ce contexte économique, mais « \textbf{investissement} » (placement de capitaux). Le sens de « renversement » existe aussi (\esp{inversión de la tendencia}), mais ici c'est exclusif.
\item \esp{las empresas} = « les entreprises », jamais « les emprises ».
\item \esp{conceder} = « accorder, octroyer » (acte positif) ; ne pas traduire par « concéder » (qui implique souvent une concession à regret ou un aveu).
\item \esp{actual} = « actuel » (qui a lieu maintenant), \textbf{jamais} « réel » (\esp{real}). \esp{Actualmente} = « actuellement », pas « en fait ».
\item \esp{realizar} = « réaliser » au sens de « concrétiser, effectuer » (\esp{realizar una inversión}), pas « prendre conscience de » (\esp{darse cuenta de}).
\end{itemize}
\end{attention}

\courssub{Lexique clé de l'aide au développement}

\begin{definition}[Subvención]
\esp{Subvención} (nom féminin) : aide financière accordée par un État ou une organisation, sans obligation de remboursement. En français : « subvention ».\\
\esp{El gobierno concede una subvención a los agricultores.} « Le gouvernement accorde une subvention aux agriculteurs. »
\end{definition}

\begin{definition}[Acuerdo bilateral / acuerdo multilateral]
\esp{Acuerdo bilateral} : accord signé entre \textbf{deux} pays (\esp{bi-} = deux). Exemple : \esp{el acuerdo entre España y Camerún}.\\
\esp{Acuerdo multilateral} : accord signé entre \textbf{plusieurs} pays ou organisations (\esp{multi-} = plusieurs). Exemple : \esp{los acuerdos de la Unión Europea o las Naciones Unidas}.\\
Le nom \esp{acuerdo} est masculin : \esp{un acuerdo, el acuerdo}.
\end{definition}

\begin{definition}[ONG]
\esp{ONG} = \esp{Organización No Gubernamental} (féminin, car \esp{organización} est féminin). On dit : \esp{una ONG, la ONG, muchas ONG}. L'acronyme ne change pas au pluriel à l'écrit, mais se prononce [o-ené-gés].
\end{definition}

\begin{attention}[Le genre de « ayuda » et « ONG »]
\begin{itemize}
\item \esp{Ayuda} est \textbf{féminin} : on dit \esp{la ayuda}, \esp{una ayuda}, \esp{mucha ayuda}, \esp{esta ayuda}. Ne confondez pas avec \esp{el agua} : \esp{agua} est féminin mais prend \esp{el} devant \emph{a} tonique (phonétique) — on dit quand même \esp{mucha agua}. Pour \esp{ayuda}, aucune exception phonétique : toujours \esp{la ayuda}. Erreur typique : *\esp{el ayuda}, *\esp{un ayuda}.
\item \esp{ONG} est féminin (\esp{la/una ONG}) car il signifie \esp{Organización No Gubernamental}. Erreur typique : *\esp{el ONG}, *\esp{un ONG}.
\end{itemize}
\end{attention}

\begin{propriete}[Sino vs Pero : la correction après négation]
Après une négation (\esp{no}), on utilise \textbf{\esp{sino}} (et non \esp{pero}) pour introduire une idée opposée qui \textbf{corrige} ou \textbf{remplace} la première.
\begin{itemize}
\item \esp{No necesita ayuda, \textbf{sino} comercio justo.} (Il n'a pas besoin d'aide, \textbf{mais} [il a besoin] de commerce équitable.) $\rightarrow$ Correction/Remplacement.
\item \esp{Es caro, \textbf{pero} bueno.} (C'est cher, \textbf{mais} c'est bon.) $\rightarrow$ Simple opposition/ajout, pas de négation avant.
\end{itemize}
\emph{Dans le texte :} \esp{« no necesita solamente ayuda, \textbf{sino} también un comercio justo »}.
\end{propriete}

\begin{propriete}[Para que + Subjonctif : la finalité]
\esp{Para que} (afin que) introduit une subordonnée de \textbf{finalité} et exige le \textbf{subjonctif} (car l'action n'est pas encore réalisée, elle est visée).
\begin{center}
\begin{tabular}{|l|l|l|}\hline
\textbf{Sujet principal} & \textbf{Para que} & \textbf{Sujet secondaire + Subjonctif} \\ \hline
\esp{Muchas ONG trabajan} & \esp{para que} & \esp{la cooperación \textbf{llegue} (subj. présent) a la población.} \\ \hline
\esp{El gobierno invierte} & \esp{para que} & \esp{los jóvenes \textbf{encuentren} (subj. présent) empleo.} \\ \hline
\end{tabular}
\end{center}
\emph{Rappel :} Subjonctif présent 3e pers. sg. de \esp{llegar} $\rightarrow$ \esp{llegue} ; de \esp{encontrar} (changement o$\rightarrow$ue) $\rightarrow$ \esp{encuentre}.
\end{propriete}

\begin{savaistu}[L'Espagne et le Cameroun : une coopération vivante]
L'Espagne entretient une coopération active avec le Cameroun, notamment dans l'éducation (programmes de bourses pour les étudiants camerounais, formation des enseignants d'espagnol via l'AECID) et la santé. L'espagnol est d'ailleurs la première langue vivante étrangère apprise dans les lycées camerounais : le voisin hispanophone, la Guinée équatoriale, seul pays d'Afrique subsaharienne de langue officielle espagnole, partage une frontière avec le Cameroun. Apprendre l'espagnol, c'est aussi pouvoir communiquer avec ce voisin immédiat et accéder aux opportunités de la \esp{Cooperación Sur-Sur} !
\end{savaistu}

\courssub{Structure du texte : la carte mentale}

Le texte s'organise en quatre paragraphes qui progressent du constat (le chômage des jeunes) vers les solutions (aides, accords, ONG) puis vers une ouverture critique (le commerce équitable).

\begin{popfigure}
\begin{tikzpicture}[mindmap, grow cyclic, every node/.style={concept, align=center}, text=white,
root concept/.append style={concept color=popblue, font=\bfseries},
level 1/.append style={sibling angle=90, level distance=3.6cm},
concept color=popblue]
\node[root concept]{El mundo laboral hoy}
  child[concept color=poporange]{ node[align=center]{1. El trabajo hoy\\ \footnotesize empleo digno\\ \footnotesize jóvenes africanos} }
  child[concept color=popgreen]{ node[align=center]{2. Ayuda al desarrollo\\ \footnotesize subvenciones\\ \footnotesize inversiones} }
  child[concept color=poppurple]{ node[align=center]{3. Acuerdos\\ \footnotesize bilaterales\\ \footnotesize multilaterales} }
  child[concept color=poppink]{ node[align=center]{4. ONG y futuro\\ \footnotesize comercio justo\\ \footnotesize empleo duradero} };
\end{tikzpicture}
\legende{Carte mentale de la structure du texte « El mundo laboral hoy »}
\end{popfigure}

\courssub{Première formulation orale du thème et de la thèse}

\begin{methode}[Dégager le thème et la thèse d'un texte]
\begin{enumerate}
\item \textbf{Le thème} (de quoi parle le texte ?) se formule en un groupe nominal : « le monde du travail et l'aide au développement en Afrique » — \esp{el mundo laboral y la ayuda al desarrollo en África}.
\item \textbf{La thèse} (que dit l'auteur sur ce thème ?) se formule en une phrase complète. Repérez les connecteurs d'opinion : \esp{sin embargo}, \esp{algunos economistas piensan que}, \esp{el verdadero desarrollo debe...}.
\item \textbf{Formulation orale :} entraînez-vous à dire à voix haute : \esp{El texto trata del mundo laboral en África. El autor defiende la idea de que la ayuda es necesaria, pero no suficiente: el verdadero desarrollo debe crear empleo duradero para los jóvenes.}
\end{enumerate}
\end{methode}

\begin{aretenir}[L'essentiel de cette étape]
\begin{itemize}
\item Le texte est un article économique en 4 paragraphes : constat, aides, accords, ONG et avenir.
\item Les mots transparents (\esp{cooperación, inversión, multilateral, organización, población, desarrollo}) ouvrent la porte du sens.
\item Le lexique clé : \esp{subvención, ayuda al desarrollo, acuerdo bilateral/multilateral, inversión, cooperación, ONG}.
\item Points de grammaire du texte : \esp{sino} (correction après négation), \esp{para que + subjonctif} (finalité), genre de \esp{ayuda} et \esp{ONG}.
\item La thèse : l'aide est nécessaire mais insuffisante ; l'Afrique a besoin de commerce équitable et d'emplois durables.
\end{itemize}
\end{aretenir}

\courssub{Entraînement à la traduction (Thème / Version)}

\begin{methode}[Méthode du thème (Français $\rightarrow$ Espagnol) pour un texte économique]
\begin{enumerate}
\item \textbf{Analyse :} Repérez les termes techniques (ne les traduisez pas littéralement : \emph{investissement} $\neq$ \esp{inversión} ? Si, ici oui ! Mais \emph{subvention} = \esp{subvención}).
\item \textbf{Grammaire :} Vérifiez les genres (\esp{la ayuda, la ONG, el acuerdo}), les régimes prépositionnels (\esp{conceder \textbf{a} alguien}, \esp{invertir \textbf{en} algo}), les temps (présent pour faits actuels, \esp{ha cambiado} pour passé récent).
\item \textbf{Reformulation :} Évitez le calque. Ex: « Les entreprises ne peuvent pas les embaucher » $\rightarrow$ \esp{Las empresas no pueden contratarlos} (pas *\esp{embauchar}).
\item \textbf{Relecture :} Accords, accents, \esp{¿ ¡}, \esp{ñ}.
\end{enumerate}
\end{methode}

\begin{exemplebox}[Mini-thème : phrases du texte]
Traduisez en espagnol :
\begin{enumerate}
\item Les pays riches accordent des subventions et réalisent des investissements en Afrique.
\item Ces aides s'organisent au moyen d'accords bilatéraux ou multilatéraux.
\item Les ONG travaillent sur le terrain pour que la coopération parvienne à la population.
\item Cependant, l'Afrique a besoin non seulement d'aide, mais aussi de commerce équitable.
\end{enumerate}
\tcblower
\textbf{Solution :}
\begin{enumerate}
\item \esp{Los países ricos conceden subvenciones y realizan inversiones en África.}
\item \esp{Estas ayudas se organizan mediante acuerdos bilaterales o multilaterales.}
\item \esp{Las ONG trabajan sobre el terreno para que la cooperación llegue a la población.}
\item \esp{Sin embargo, África necesita no solo ayuda, sino también comercio justo.} (Ou : \esp{no solamente ayuda, sino también...})
\end{enumerate}
\end{exemplebox}

\begin{exoresolu}[Questions de compréhension sur le texte]
Répondez en espagnol, par des phrases complètes :
\begin{enumerate}
\item \esp{¿Qué buscan muchos jóvenes africanos hoy?}
\item \esp{¿Qué construyen los países ricos en África?}
\item \esp{¿Cuál es la diferencia entre un acuerdo bilateral y un acuerdo multilateral?}
\item \esp{Según algunos economistas, ¿qué necesita África además de ayuda?}
\item \esp{¿Para qué trabajan las ONG sobre el terreno? (Utilisez \textbf{para que} + subjonctif)}
\end{enumerate}
\tcblower
\textbf{Solutions :}
\begin{enumerate}
\item \esp{Muchos jóvenes africanos buscan un empleo digno.}
\item \esp{Los países ricos construyen escuelas, hospitales y carreteras.}
\item \esp{Un acuerdo bilateral se firma entre dos países, mientras que un acuerdo multilateral se firma entre varios países u organizaciones.}
\item \esp{Según algunos economistas, África necesita también un comercio justo y más inversiones productivas.}
\item \esp{Las ONG trabajan sobre el terreno para que la cooperación llegue a la población.}
\end{enumerate}
\end{exoresolu}

\begin{exercice}[À vous de jouer]
\begin{enumerate}
\item Trouvez dans le texte trois mots transparents supplémentaires (non listés dans le tableau) et donnez leur équivalent français.
\item Complétez avec l'article défini ou indéfini correct (\esp{el / la / un / una}) : \esp{... ayuda al desarrollo}, \esp{... acuerdo bilateral}, \esp{... subvención}, \esp{... ONG}, \esp{... inversión productiva}.
\item Formulez oralement, en deux phrases en espagnol, le thème et la thèse du texte.
\item \textbf{Version :} Traduisez en français : \esp{« El verdadero desarrollo debe crear empleo duradero para los jóvenes, no solo recibir ayudas exteriores. »}
\item \textbf{Thème :} Mettez en espagnol : « Le véritable développement doit créer un emploi durable pour les jeunes, et non pas seulement recevoir des aides extérieures. »
\end{enumerate}
\end{exercice}

\begin{corrige}[Exercice « À vous de jouer »]
\begin{enumerate}
\item Exemples possibles : \esp{hospital} (hôpital), \esp{población} (population), \esp{desarrollo} (développement), \esp{necesario} (nécessaire), \esp{justo} (juste/équitable), \esp{verdad} (vérité/vrai).
\item \esp{la ayuda al desarrollo} ; \esp{un/el acuerdo bilateral} ; \esp{una/la subvención} ; \esp{una ONG} (féminin : \esp{organización no gubernamental}) ; \esp{una/la inversión productiva} (féminin).
\item Modèle : \esp{El texto trata del trabajo y de la ayuda al desarrollo en África. La tesis del autor es que la ayuda es necesaria, pero que el verdadero desarrollo exige comercio justo y empleo duradero.}
\item « Le véritable développement doit créer un emploi durable pour les jeunes, et pas seulement recevoir des aides extérieures. »
\item \esp{El verdadero desarrollo debe crear empleo duradero para los jóvenes, no solo recibir ayudas exteriores.}
\end{enumerate}
\end{corrige}

\coursec{Comprensión del texto}

Dans la section précédente, nous avons lu et découvert le texte \esp{« El mundo laboral hoy »}, en dégagant son vocabulaire et sa situation d'énonciation. Il ne suffit pas de « comprendre dans l'ensemble » : au baccalauréat, l'épreuve de compréhension de l'écrit exige des réponses \cle{précises}, \cle{argumentées} et \cle{justifiées par le texte}. Cette section vous apprend à passer d'une lecture globale à une lecture méthodique, paragraphe par paragraphe.

\courssub{Relecture guidée du texte}

Avant de répondre aux questions, relisons le texte avec un objectif précis : repérer, dans chaque paragraphe, \textbf{une seule idée directrice}.

\begin{textebox}[El mundo laboral hoy (rappel du texte support)]
El mundo laboral ha cambiado mucho en las últimas décadas. En África, muchos jóvenes buscan un trabajo estable, pero el desempleo juvenil sigue siendo muy alto. Por eso, la ayuda al desarrollo sigue siendo necesaria para crear empleo y formar a los trabajadores.

Los países ricos conceden subvenciones y ayudas a los países del Sur. Estas ayudas pueden ser bilaterales, es decir, entre dos países, o multilaterales, cuando participan varios países u organismos internacionales como el Banco Mundial. Sin embargo, algunos economistas piensan que África no necesita solamente ayuda, sino también inversión y comercio justo.

Además, las ONG desempeñan un papel importante: construyen escuelas, forman a los campesinos y apoyan a las mujeres emprendedoras. En Camerún, por ejemplo, muchas cooperativas de productores de cacao reciben apoyo técnico y financiero de organizaciones internacionales.

En conclusión, el mundo laboral de hoy exige cooperación entre los países. La ayuda al desarrollo es útil, pero debe ir acompañada de inversiones, de educación y de acuerdos justos para que los jóvenes africanos encuentren un futuro digno en su propio continente.
\end{textebox}

\textbf{Glossaire :}
\begin{itemize}
\item \esp{el desempleo juvenil} : le chômage des jeunes
\item \esp{conceder} : accorder, octroyer
\item \esp{es decir} : c'est-à-dire
\item \esp{el comercio justo} : le commerce équitable
\item \esp{desempeñar un papel} : jouer un rôle
\item \esp{los campesinos} : les paysans
\item \esp{las mujeres emprendedoras} : les femmes entrepreneures
\item \esp{apoyo técnico y financiero} : appui technique et financier
\item \esp{ir acompañada de} : être accompagnée de
\item \esp{digno} : digne
\end{itemize}

\courssub{Comprensión global : thème, thèse, intention}

Avant les questions de détail, il faut répondre à trois questions générales qui structurent toute lecture.

\begin{definition}[Les trois questions de compréhension globale]
\begin{itemize}
\item \textbf{Le thème} (\esp{el tema}) : de quoi parle le texte ? Réponse en quelques mots : \esp{el mundo laboral y la ayuda al desarrollo en África}.
\item \textbf{La thèse} (\esp{la tesis}) : quelle position l'auteur défend-il ? Ici : l'aide au développement est nécessaire mais insuffisante ; il faut aussi l'investissement, l'éducation et des accords équitables.
\item \textbf{L'intention de l'auteur} (\esp{la intención del autor}) : pourquoi écrit-il ? Ici : \emph{informer} et \emph{sensibiliser} le lecteur aux enjeux du développement économique en Afrique.
\end{itemize}
\end{definition}

\begin{methode}[Répondre à une question de compréhension en trois pas]
\begin{enumerate}
\item \textbf{Répondre par une phrase complète}, en espagnol, en reprenant les mots de la question. Jamais un mot isolé.
\item \textbf{Citer le texte entre guillemets} (« ... ») pour justifier la réponse : choisir le segment exact qui prouve la réponse.
\item \textbf{Traduire la citation en français} si la consigne le demande (souvent : \esp{« Justifica tu respuesta citando el texto y traduce la cita »}).
\end{enumerate}
\end{methode}

\begin{exemplebox}[Modèle de réponse rédigée]
\textbf{Question :} \esp{¿Qué buscan muchos jóvenes africanos?}

\textbf{Réponse :} \esp{Muchos jóvenes africanos buscan un trabajo estable. El texto dice: « En África, muchos jóvenes buscan un trabajo estable ».}

\textbf{Traduction de la citation :} « En Afrique, beaucoup de jeunes cherchent un emploi stable. »
\end{exemplebox}

\begin{attention}[L'erreur classique : recopier tout le paragraphe]
Beaucoup d'élèves, par peur de se tromper, recopient \textbf{tout le paragraphe} comme justification. C'est une erreur sanctionnée : cela prouve que vous n'avez pas repéré l'information. La justification doit être \cle{un segment court et pertinent} : une phrase, parfois seulement un groupe de mots. Mauvaise justification : copier les quatre lignes du paragraphe 2. Bonne justification : \esp{« Estas ayudas pueden ser bilaterales, es decir, entre dos países, o multilaterales, cuando participan varios países »}.
\end{attention}

\courssub{Comprensión fine : questions par paragraphe}

Voici les questions type Bac, avec leur localisation dans le texte. Entraînez-vous à y répondre avant de consulter les corrigés.

\begin{exercice}[Questions de compréhension (type Bac)]
\begin{enumerate}
\item \esp{¿Qué buscan muchos jóvenes africanos?} (paragraphe 1)
\item \esp{¿Cuál es la diferencia entre un acuerdo bilateral y un acuerdo multilateral?} (paragraphe 2)
\item \esp{¿Qué hacen las ONG?} (paragraphe 3)
\item \esp{Según algunos economistas, ¿qué necesita África además de ayuda?} (paragraphe 2)
\item \esp{¿Cuál es la idea principal del texto?} (ensemble du texte)
\end{enumerate}
\end{exercice}

\begin{corrige}[Questions de compréhension]
\begin{enumerate}
\item \esp{Muchos jóvenes africanos buscan un trabajo estable.} Citation : \esp{« muchos jóvenes buscan un trabajo estable »} (« beaucoup de jeunes cherchent un emploi stable »).
\item \esp{Un acuerdo bilateral se firma entre dos países, mientras que un acuerdo multilateral reúne a varios países u organismos internacionales.} Citation : \esp{« bilaterales, es decir, entre dos países, o multilaterales, cuando participan varios países u organismos internacionales »} (« bilatérales, c'est-à-dire entre deux pays, ou multilatérales, quand participent plusieurs pays ou organismes internationaux »).
\item \esp{Las ONG construyen escuelas, forman a los campesinos y apoyan a las mujeres emprendedoras.} Citation : \esp{« construyen escuelas, forman a los campesinos y apoyan a las mujeres emprendedoras »} (« elles construisent des écoles, forment les paysans et soutiennent les femmes entrepreneures »).
\item \esp{Según algunos economistas, África necesita también inversión y comercio justo.} Citation : \esp{« África no necesita solamente ayuda, sino también inversión y comercio justo »} (« l'Afrique n'a pas seulement besoin d'aide, mais aussi d'investissement et de commerce équitable »).
\item \esp{La idea principal es que la ayuda al desarrollo es necesaria pero debe ir acompañada de inversiones, educación y acuerdos justos.} Citation : \esp{« La ayuda al desarrollo es útil, pero debe ir acompañada de inversiones, de educación y de acuerdos justos »} (« l'aide au développement est utile, mais elle doit s'accompagner d'investissements, d'éducation et d'accords équitables »).
\end{enumerate}
\end{corrige}

\courssub{Repérage des idées directrices}

Un texte argumentatif bien construit obéit à un principe simple : \cle{un paragraphe = une idée principale}. Repérer ces idées, c'est fabriquer le plan du texte, compétence indispensable pour le commentaire et la synthèse.

\begin{exemplebox}[Les idées directrices du texte]
\begin{itemize}
\item \textbf{Paragraphe 1 :} le chômage des jeunes rend l'aide au développement nécessaire.
\item \textbf{Paragraphe 2 :} les formes de l'aide (bilatérale, multilatérale) et ses limites.
\item \textbf{Paragraphe 3 :} le rôle concret des ONG sur le terrain (exemple du Cameroun).
\item \textbf{Paragraphe 4 :} conclusion : l'aide doit s'accompagner d'investissement et d'éducation.
\end{itemize}
\end{exemplebox}

\begin{popfigure}
\begin{tikzpicture}
\draw[fill=popblueL, draw=popblue, thick] (0,0) rectangle (3.2,6.4);
\draw[fill=popgreenL, draw=popgreen, thick] (3.6,0) rectangle (12.4,6.4);
\node[text=popdark, font=\bfseries] at (1.6,5.9) {Párrafo};
\node[text=popdark, font=\bfseries] at (8,5.9) {Idea principal};
\draw[popline] (0,5.5) -- (12.4,5.5);
\node[text=popdark, align=center] at (1.6,4.7) {1};
\node[text=popdark, align=left, text width=8.2cm] at (4.1,4.7) {\esp{El desempleo juvenil es alto, la ayuda al desarrollo es necesaria.}};
\draw[popline] (0,4.1) -- (12.4,4.1);
\node[text=popdark, align=center] at (1.6,3.3) {2};
\node[text=popdark, align=left, text width=8.2cm] at (4.1,3.3) {\esp{Ayudas bilaterales y multilaterales; África necesita también inversión y comercio justo.}};
\draw[popline] (0,2.7) -- (12.4,2.7);
\node[text=popdark, align=center] at (1.6,1.9) {3};
\node[text=popdark, align=left, text width=8.2cm] at (4.1,1.9) {\esp{Las ONG actúan sobre el terreno: escuelas, formación, apoyo a las mujeres.}};
\draw[popline] (0,1.3) -- (12.4,1.3);
\node[text=popdark, align=center] at (1.6,0.6) {4};
\node[text=popdark, align=left, text width=8.2cm] at (4.1,0.6) {\esp{Conclusión: la ayuda debe ir acompañada de inversión, educación y acuerdos justos.}};
\end{tikzpicture}
\legende{Tableau « Párrafo / Idea principal » : modèle complété. En classe, on distribue le tableau vide et les élèves le remplissent.}
\end{popfigure}

\courssub{Vrai ou faux : justifier par une citation}

L'exercice \esp{« verdadero o falso »} est un classique du Bac. La règle d'or : une réponse sans justification \textbf{ne rapporte aucun point}, même si elle est correcte.

\begin{exoresolu}[Verdadero o falso]
\textbf{Énoncé :} Indique si les affirmations suivantes sont \esp{verdaderas (V)} ou \esp{falsas (F)} et justifie chaque réponse par une citation du texte.
\begin{enumerate}
\item \esp{El desempleo juvenil ha desaparecido en África.}
\item \esp{Los acuerdos multilaterales solo conciernen a dos países.}
\item \esp{Las ONG apoyan a las mujeres emprendedoras.}
\item \esp{Según el texto, la ayuda al desarrollo ya no es necesaria.}
\end{enumerate}
\tcblower
\textbf{Solution détaillée :}
\begin{enumerate}
\item \textbf{Falso.} Citation : \esp{« el desempleo juvenil sigue siendo muy alto »} (« le chômage des jeunes reste très élevé »). Le texte dit le contraire de l'affirmation.
\item \textbf{Falso.} Citation : \esp{« multilaterales, cuando participan varios países u organismos internacionales »} (« multilatéraux, quand plusieurs pays ou organismes internationaux participent »). C'est l'accord \emph{bilatéral} qui concerne deux pays.
\item \textbf{Verdadero.} Citation : \esp{« apoyan a las mujeres emprendedoras »} (« elles soutiennent les femmes entrepreneures »).
\item \textbf{Falso.} Citation : \esp{« la ayuda al desarrollo sigue siendo necesaria »} (« l'aide au développement reste nécessaire »).
\end{enumerate}
\end{exoresolu}

\courssub{Les connecteurs logiques et leur fonction}

Les \cle{connecteurs} sont les panneaux de signalisation du texte : ils indiquent la relation logique entre les idées. Les repérer permet de suivre l'argumentation et de prévoir la suite.

\begin{propriete}[Les trois connecteurs du texte]
\begin{itemize}
\item \esp{por eso} : exprime la \textbf{conséquence} (= c'est pourquoi, donc). \esp{El desempleo es alto. Por eso, la ayuda al desarrollo sigue siendo necesaria.} « Le chômage est élevé. C'est pourquoi l'aide au développement reste nécessaire. »
\item \esp{sin embargo} : exprime l'\textbf{opposition} (= cependant, pourtant). \esp{Las ayudas existen. Sin embargo, algunos economistas piensan que no bastan.} « Les aides existent. Cependant, certains économistes pensent qu'elles ne suffisent pas. »
\item \esp{además} : exprime l'\textbf{addition} (= de plus, en outre). \esp{Además, las ONG desempeñan un papel importante.} « De plus, les ONG jouent un rôle important. »
\end{itemize}
\end{propriete}

\begin{center}
\begin{tabularx}{\linewidth}{|X|X|X|}
\hline
\rowcolor{popblueL} Conector & Función lógica & Équivalents français \\
\hline
\esp{por eso, por lo tanto, entonces} & consecuencia & c'est pourquoi, donc, par conséquent \\
\hline
\esp{sin embargo, pero, no obstante} & oposición & cependant, mais, néanmoins \\
\hline
\esp{además, también, asimismo} & adición & de plus, aussi, également \\
\hline
\esp{es decir, o sea} & explicación & c'est-à-dire \\
\hline
\esp{en conclusión, en resumen} & conclusión & en conclusion, en résumé \\
\hline
\esp{por ejemplo} & ejemplificación & par exemple \\
\hline
\end{tabularx}
\end{center}

\begin{exemplebox}[Exemples traduits]
\begin{itemize}
\item \esp{Por eso, la ayuda al desarrollo sigue siendo necesaria.} « C'est pourquoi l'aide au développement reste nécessaire. »
\item \esp{Sin embargo, África no necesita solamente ayuda.} « Cependant, l'Afrique n'a pas seulement besoin d'aide. »
\item \esp{Además, las ONG construyen escuelas.} « De plus, les ONG construisent des écoles. »
\item \esp{En conclusión, el mundo laboral exige cooperación.} « En conclusion, le monde du travail exige de la coopération. »
\end{itemize}
\end{exemplebox}

\begin{attention}[Faux amis et calques à éviter]
\begin{itemize}
\item \esp{sin embargo} ne se traduit jamais par « sans embargo » : c'est « cependant ».
\item \esp{además} signifie « de plus », pas « en plus de » tout seul ; « en plus de » se dit \esp{además de} : \esp{Además de ayuda, África necesita inversión.} « En plus de l'aide, l'Afrique a besoin d'investissement. »
\item Ne traduisez pas « c'est pourquoi » par \esp{es por qué} : la bonne forme est \esp{por eso} ou \esp{por lo tanto}.
\item \esp{actualmente} signifie « actuellement », pas « en fait » (faux ami de \emph{actually}) ; attention aussi à \esp{asistir} (« assister à », pas « aider » : \esp{ayudar}).
\end{itemize}
\end{attention}

\courssub{Reformulation : la paraphrase guidée}

La \cle{paraphrase} consiste à dire la même idée avec d'autres mots. Elle prouve au correcteur que vous avez réellement compris, et elle est indispensable dans le commentaire et la synthèse. Deux techniques simples : remplacer les mots par des \textbf{synonymes} et transformer la \textbf{structure} de la phrase (actif/passif, nom/verbe).

\begin{exemplebox}[Paraphrases guidées du texte]
\begin{itemize}
\item Phrase originale : \esp{El desempleo juvenil sigue siendo muy alto.}\\
Paraphrase : \esp{Muchos jóvenes todavía no encuentran trabajo.} (« Beaucoup de jeunes ne trouvent toujours pas de travail. »)
\item Phrase originale : \esp{Las ONG apoyan a las mujeres emprendedoras.}\\
Paraphrase : \esp{Las organizaciones no gubernamentales ayudan a las mujeres que crean empresas.} (« Les organisations non gouvernementales aident les femmes qui créent des entreprises. »)
\item Phrase originale : \esp{La ayuda debe ir acompañada de inversiones.}\\
Paraphrase : \esp{No basta con dar ayuda: también hace falta invertir.} (« Il ne suffit pas de donner de l'aide : il faut aussi investir. »)
\end{itemize}
\end{exemplebox}

\begin{methode}[Réussir une paraphrase en quatre pas]
\begin{enumerate}
\item \textbf{Comprendre} la phrase originale et la traduire mentalement en français.
\item \textbf{Remplacer} les mots-clés par des synonymes : \esp{buscar} $\rightarrow$ \esp{intentar encontrar} ; \esp{apoyar} $\rightarrow$ \esp{ayudar} ; \esp{necesitar} $\rightarrow$ \esp{hacer falta}.
\item \textbf{Transformer} la structure : nom $\rightarrow$ verbe (\esp{el desempleo} $\rightarrow$ \esp{estar desempleado}), actif $\rightarrow$ passif.
\item \textbf{Vérifier} que le sens n'a pas changé : une paraphrase ne doit ni ajouter ni retirer d'idée.
\end{enumerate}
\end{methode}

\begin{exercice}[À vous de jouer : paraphrases et connecteurs]
\begin{enumerate}
\item Paraphrasez : \esp{Los países ricos conceden subvenciones a los países del Sur.}
\item Paraphrasez : \esp{Muchas cooperativas de cacao reciben apoyo técnico y financiero.}
\item Complétez avec \esp{por eso}, \esp{sin embargo} ou \esp{además} : \esp{El desempleo juvenil es muy alto. \ldots{}, muchos jóvenes emigran. Las ayudas existen; \ldots{}, no siempre llegan a los más pobres. \ldots{}, las ONG forman a los campesinos.}
\end{enumerate}
\end{exercice}

\begin{corrige}[Paraphrases et connecteurs]
\begin{enumerate}
\item \esp{Las naciones desarrolladas dan ayudas económicas a los países pobres.} (synonymes : \esp{países ricos} $\rightarrow$ \esp{naciones desarrolladas} ; \esp{conceder} $\rightarrow$ \esp{dar}.)
\item \esp{Numerosas cooperativas de productores de cacao son ayudadas técnica y económicamente.} (transformation actif $\rightarrow$ passif.)
\item \esp{El desempleo juvenil es muy alto. Por eso, muchos jóvenes emigran. Las ayudas existen; sin embargo, no siempre llegan a los más pobres. Además, las ONG forman a los campesinos.}
\end{enumerate}
\end{corrige}

\begin{aretenir}[L'essentiel de la section]
\begin{itemize}
\item Répondre à une question de compréhension = \textbf{phrase complète + citation courte entre guillemets + traduction si demandée}.
\item Ne jamais recopier un paragraphe entier : choisir le \textbf{segment pertinent}.
\item Un paragraphe = \textbf{une idée directrice} : c'est le plan du texte.
\item Au vrai/faux, toute réponse non justifiée vaut zéro.
\item Les connecteurs structurent l'argumentation : \esp{por eso} (conséquence), \esp{sin embargo} (opposition), \esp{además} (addition).
\item La paraphrase dit la même idée autrement, grâce aux synonymes et aux transformations de structure.
\end{itemize}
\end{aretenir}

\begin{savaistu}[Cameroun, Guinée équatoriale et coopération internationale]
Le Cameroun et la Guinée équatoriale (seul pays africain hispanophone) sont membres de la \esp{CEMAC} (Communauté Économique et Monétaire de l'Afrique Centrale) et bénéficient d'accords de coopération bilatéraux (avec l'Espagne, la France, la Chine, les USA) et multilatéraux (Banque Mondiale \esp{Banco Mundial}, Banque Africaine de Développement \esp{Banco Africano de Desarrollo}, FMI \esp{FMI}, ONU \esp{ONU}). La Guinée équatoriale, pays pétrolier, investit dans les infrastructures ; le Cameroun, pays agricole (cacao, café, coton) et pétrolier, mise sur la transformation locale et l'industrialisation via le \esp{Plan National de Développement (PND)}. Les \esp{ONG} espagnoles (\esp{Manos Unidas, Cruz Roja, Ayuda en Acción}) et locales sont très actives dans la formation professionnelle rurale (ex. \esp{cooperativas de cacao} dans le Sud et le Centre).
\end{savaistu}

\begin{experience}[Débat oral : « ¿Ayuda o inversión? »]
\textbf{Consigne :} En groupes de 4, préparez un débat de 5 minutes. Deux élèves défendent la thèse « L'aide au développement est indispensable » (\esp{La ayuda es indispensable}), deux autres « L'investissement et le commerce équitable sont la vraie solution » (\esp{La inversión y el comercio justo son la solución}). Utilisez le vocabulaire de la leçon (\esp{subvención, bilateral, multilateral, ONG, cooperativa, comercio justo, desempleo juvenil}) et les connecteurs (\esp{por eso, sin embargo, además, en conclusión}). Un élève note les arguments clés pour une synthèse finale écrite.
\end{experience}

\begin{exercice}[Production guidée : rédaction d'un article d'opinion]
Rédigez en espagnol (120--150 mots) un article pour le journal du lycée intitulé \esp{« El futuro del trabajo en África : ¿ayuda o inversión? »}. Structurez votre texte :
\begin{enumerate}
\item \textbf{Introduction :} accroche + thèse (votre position).
\item \textbf{Développement (2 paragraphes) :} arguments pour l'aide (ex. ONG, formation) + arguments pour l'investissement/commerce juste (ex. entreprises, dignité). Utilisez \esp{por un lado / por otro lado, sin embargo, además}.
\item \textbf{Conclusion :} synthèse + ouverture (rôle de la jeunesse, éducation).
\end{enumerate}
\textbf{Mots-outils :} \esp{en mi opinión, es fundamental, no basta con, hace falta, por tanto, para que + subjonctif}.
\end{exercice}

\begin{corrige}[Production guidée : proposition de corrigé]
\textbf{El futuro del trabajo en África: ¿ayuda o inversión?}

En mi opinión, el desarrollo de África no depende solo de la ayuda exterior, sino de una inversión real y justa. Es cierto que la ayuda al desarrollo sigue siendo necesaria: por eso, las ONG construyen escuelas y forman a campesinos en Camerún. Además, los acuerdos multilaterales aportan fondos para la salud y la educación.

Sin embargo, la ayuda no basta. El desempleo juvenil es muy alto y los jóvenes buscan un trabajo digno, no solo asistencia. Por lo tanto, hace falta inversión en empresas locales y comercio justo para que los productos africanos (cacao, café, algodón) se vendan a buen precio. Por ejemplo, las cooperativas de cacao necesitan máquinas para transformar el grano, no solo dinero para comprar semillas.

En conclusión, la cooperación internacional debe ir acompañada de educación técnica y acuerdos comerciales equitativos. Solo así los jóvenes africanos construirán su futuro en su propio continente.
\end{corrige}

Vous savez désormais lire le texte avec méthode et précision. La section suivante approfondira le \textbf{lexique thématique} de l'aide au développement et du monde du travail, dont vous avez déjà rencontré les mots-clés dans ces questions.

\coursec{Lexique thématique : l'aide au développement et le monde du travail}

Dans la section précédente, nous avons lu et compris le texte \esp{El mundo laboral hoy} : nous avons identifié le thème, les idées principales et l'organisation du propos. Mais pour commenter un texte économique en espagnol — et surtout pour le traduire sans faute —, il ne suffit pas de le comprendre globalement : il faut maîtriser son \cle{vocabulaire spécialisé}. C'est l'objet de cette section : construire, champ par champ, le lexique de l'aide au développement et du monde du travail, en évitant les pièges que l'espagnol tend aux francophones.

\courssub{Texte d'observation : le lexique en contexte}

Avant toute liste de mots, observons ce vocabulaire \emph{en situation}. Lisons le texte suivant et soulignons mentalement tous les mots qui appartiennent au champ de l'aide au développement.

\begin{textebox}[Cooperación y empleo en África Central]
La cooperación internacional ocupa un lugar central en el desarrollo económico de África Central. El año pasado, el gobierno de Camerún firmó un acuerdo bilateral con España para financiar proyectos de formación profesional en Yaoundé y Douala. Gracias a este acuerdo, el Estado español concedió una subvención destinada a crear empleo duradero para los jóvenes.

Sin embargo, la ayuda al desarrollo no depende solo de los gobiernos. Numerosas ONG trabajan sobre el terreno para llegar a la población rural más pobre. Sus voluntarios luchan contra la pobreza y colaboran con las Naciones Unidas en programas de ayuda humanitaria. Además, la Unión Europea ha firmado acuerdos multilaterales con varios países africanos para promover el comercio justo: los productores de cacao y de café reciben así un precio estable por su trabajo.

El mundo laboral cambia rápidamente. Las empresas necesitan trabajadores cualificados, pero el desempleo sigue siendo alto entre los jóvenes diplomados. Por eso, muchos expertos piden más inversión extranjera y préstamos a bajo interés para las pequeñas empresas. Un donativo puntual puede ayudar en una emergencia, pero solo la inversión a largo plazo y la cooperación verdadera permiten un desarrollo económico sostenible.

En conclusión, la ayuda al desarrollo es eficaz cuando combina la financiación, los acuerdos entre Estados y el trabajo diario de las organizaciones sobre el terreno.
\end{textebox}

\textbf{Glossaire :} \esp{la formación profesional} = la formation professionnelle ; \esp{destinada a} = destinée à ; \esp{cualificados} = qualifiés ; \esp{a bajo interés} = à faible taux d'intérêt ; \esp{puntual} = ponctuel, occasionnel ; \esp{a largo plazo} = à long terme ; \esp{sostenible} = durable.

\textbf{Questions de repérage :}
\begin{enumerate}
\item Quels mots du texte désignent de l'argent donné ou prêté ?
\item Quels mots désignent des organisations ou des institutions ?
\item Relevez trois expressions verbales liées à la lutte contre la pauvreté.
\end{enumerate}

\courssub{Champ 1 — L'aide et le financement}

\begin{definition}[Les mots de l'argent]
\begin{itemize}
\item \esp{la ayuda al desarrollo} : l'aide au développement (ensemble des aides versées aux pays en développement).
\item \esp{la subvención} : la subvention (argent donné par un État ou une institution, sans remboursement).
\item \esp{el préstamo} : le prêt (argent prêté, à rembourser, souvent avec des intérêts).
\item \esp{la inversión} : l'investissement (argent placé dans une activité pour obtenir un résultat économique).
\item \esp{el donativo} : le don (contribution volontaire, souvent d'un particulier ou d'une association).
\item \esp{financiar} : financer ; \esp{conceder} : accorder (une subvention, un prêt).
\end{itemize}
\end{definition}

\begin{exemplebox}[Exemples traduits]
\esp{España concede una subvención a Camerún.} = L'Espagne accorde une subvention au Cameroun.

\esp{El banco concede préstamos a bajo interés a las pequeñas empresas.} = La banque accorde des prêts à faible taux d'intérêt aux petites entreprises.

\esp{La inversión extranjera crea empleo en Douala.} = L'investissement étranger crée de l'emploi à Douala.

\esp{Los voluntarios recogieron donativos para el hospital.} = Les volontaires ont collecté des dons pour l'hôpital.
\end{exemplebox}

\begin{attention}[Faux ami : \esp{inversión} $\neq$ invention]
\esp{Una inversión} signifie « un investissement », jamais « une invention » (qui se dit \esp{un invento}). Le verbe \esp{invertir} = investir est irrégulier (e $\rightarrow$ ie) : \esp{invierto, inviertes, invierte, invertimos, invertís, invierten}. Exemple : \esp{La empresa invierte mucho dinero en formación.} = L'entreprise investit beaucoup d'argent dans la formation.
\end{attention}

\courssub{Champ 2 — Les accords et la coopération}

\begin{definition}[Les mots de l'accord]
\begin{itemize}
\item \esp{el acuerdo bilateral} : l'accord bilatéral (entre deux États) ; \esp{el acuerdo multilateral} : l'accord multilatéral (entre plusieurs États).
\item \esp{la cooperación} : la coopération (action de travailler ensemble pour le développement).
\item \esp{el tratado} : le traité (accord formel et solennel entre États).
\item \esp{firmar un acuerdo} : signer un accord ; \esp{llegar a un acuerdo} : parvenir à un accord.
\item \esp{la ayuda humanitaria} : l'aide humanitaire (aide d'urgence : catastrophes, guerres, famines).
\end{itemize}
\end{definition}

\begin{exemplebox}[Exemples traduits]
\esp{Camerún y España firmaron un acuerdo bilateral de cooperación.} = Le Cameroun et l'Espagne ont signé un accord bilatéral de coopération.

\esp{Los países llegaron a un acuerdo multilateral sobre el comercio.} = Les pays sont parvenus à un accord multilatéral sur le commerce.

\esp{Tras las inundaciones, llegó la ayuda humanitaria internacional.} = Après les inondations, l'aide humanitaire internationale est arrivée.
\end{exemplebox}

\begin{propriete}[Construction : \esp{acuerdo} + complément]
On dit \esp{un acuerdo \textbf{entre} dos países} (un accord entre deux pays), \esp{un acuerdo \textbf{con} un país} (un accord avec un pays) et \esp{un acuerdo \textbf{sobre} un tema} (un accord sur un sujet). Le verbe \esp{firmar} se construit directement : \esp{firmar un acuerdo}, sans préposition.
\end{propriete}

\courssub{Champ 3 — Les acteurs du développement}

\begin{definition}[ONG et institutions]
\begin{itemize}
\item \esp{la ONG} = \esp{organización no gubernamental} : organisation non gouvernementale. On la prononce « o-ené-gé ». C'est un sigle toujours \textbf{féminin} : \esp{la ONG}, \esp{las ONG}.
\item \esp{el Estado} : l'État ; \esp{el gobierno} : le gouvernement.
\item \esp{la Unión Europea} : l'Union européenne ; \esp{las Naciones Unidas} : les Nations unies (toujours au pluriel).
\item \esp{el/la voluntario/a} : le/la volontaire, bénévole ; \esp{el voluntariado} : le bénévolat.
\end{itemize}
\end{definition}

\begin{exemplebox}[Exemples traduits]
\esp{Las ONG trabajan sobre el terreno.} = Les ONG travaillent sur le terrain.

\esp{Las Naciones Unidas coordinan la ayuda humanitaria.} = Les Nations unies coordonnent l'aide humanitaire.

\esp{Una voluntaria española enseña en una escuela de Yaoundé.} = Une volontaire espagnole enseigne dans une école de Yaoundé.
\end{exemplebox}

\begin{attention}[\esp{asistencia} : un faux ami dangereux]
\esp{La asistencia} signifie d'abord « la présence, l'assistance » (à une réunion, à un examen) : \esp{La asistencia a clase es obligatoria.} = La présence en classe est obligatoire. Pour « l'aide sociale », on dit \esp{la asistencia social}. Pour « aider », le verbe est \esp{ayudar} ; \esp{asistir a} signifie « assister à » (= être présent) : \esp{Asistió a la reunión.} = Il a assisté à la réunion.
\end{attention}

\courssub{Champ 4 — Le travail et l'économie}

\begin{definition}[Les mots du travail]
\begin{itemize}
\item \esp{el mundo laboral} : le monde du travail ; \esp{el mercado de trabajo} : le marché du travail.
\item \esp{el empleo} : l'emploi ; \esp{crear empleo} : créer de l'emploi.
\item \esp{el desempleo} / \esp{el paro} : le chômage ; \esp{estar en paro} : être au chômage.
\item \esp{la empresa} : l'entreprise ; \esp{el/la trabajador(a)} : le/la travailleur(se), salarié(e) ; \esp{el/la empresario/a} : le chef d'entreprise.
\item \esp{el comercio justo} : le commerce équitable ; \esp{el desarrollo económico} : le développement économique.
\end{itemize}
\end{definition}

\begin{attention}[\esp{el paro} ou \esp{el desempleo} ?]
En Espagne, le mot courant est \esp{el paro} : \esp{El paro juvenil es muy alto en España.} = Le chômage des jeunes est très élevé en Espagne. \esp{El desempleo} est le terme plus international, utilisé dans toute l'Amérique latine et dans les documents officiels. Dans un commentaire de texte, les deux sont corrects. Attention : \esp{un paro} peut aussi signifier « un arrêt de travail, une grève » en Amérique latine.
\end{attention}

\begin{savaistu}[Le commerce équitable et le cacao camerounais]
Le \esp{comercio justo} (commerce équitable) garantit un prix stable aux producteurs, même quand les cours mondiaux baissent. Le cacao et le café, secteurs clés de l'économie camerounaise, sont au cœur de ce système : grâce au commerce équitable, les planteurs de la région du Centre ou du Sud-Ouest peuvent financer l'école de leurs enfants et investir dans leurs plantations. En espagnol, on parle aussi de \esp{precio justo} (prix juste) et de \esp{comercio solidario}.
\end{savaistu}

\courssub{Expressions utiles pour le commentaire}

Ces expressions figées reviennent dans presque tous les textes économiques ; apprenez-les comme des blocs :

\begin{center}
\begin{tabularx}{\linewidth}{|X|X|}
\hline
\rowcolor{popblueL} Español & Français \\
\hline
\esp{trabajar sobre el terreno} & travailler sur le terrain \\
\hline
\esp{llegar a la población (rural)} & atteindre la population (rurale) \\
\hline
\esp{crear empleo duradero} & créer un emploi durable \\
\hline
\esp{luchar contra la pobreza} & lutter contre la pauvreté \\
\hline
\esp{promover el desarrollo económico} & promouvoir le développement économique \\
\hline
\esp{reducir el desempleo juvenil} & réduire le chômage des jeunes \\
\hline
\esp{destinar fondos a un proyecto} & affecter des fonds à un projet \\
\hline
\end{tabularx}
\end{center}

Exemple : \esp{Las ONG destinan fondos a proyectos que crean empleo duradero.} = Les ONG affectent des fonds à des projets qui créent un emploi durable.

\courssub{Carte mentale du lexique}

\begin{popfigure}
\begin{tikzpicture}[mindmap, grow cyclic, every node/.style={concept, text=white, align=center}, concept color=popblue,
level 1/.append style={level distance=3.4cm, sibling angle=90},
level 2/.append style={level distance=2.2cm, sibling angle=45}]
\node[font=\small\bfseries]{El mundo laboral y la ayuda}
child[concept color=poporange]{ node[font=\small]{Financiación}
  child{ node[font=\scriptsize]{subvención} }
  child{ node[font=\scriptsize]{préstamo} }
  child{ node[font=\scriptsize]{inversión} }
  child{ node[font=\scriptsize]{donativo} }
}
child[concept color=popgreen]{ node[font=\small]{Acuerdos}
  child{ node[font=\scriptsize]{bilateral} }
  child{ node[font=\scriptsize]{multilateral} }
  child{ node[font=\scriptsize]{tratado} }
  child{ node[font=\scriptsize]{cooperación} }
}
child[concept color=poppurple]{ node[font=\small]{Actores}
  child{ node[font=\scriptsize]{la ONG} }
  child{ node[font=\scriptsize]{el Estado} }
  child{ node[font=\scriptsize]{Naciones Unidas} }
  child{ node[font=\scriptsize]{voluntario} }
}
child[concept color=popgold]{ node[font=\small]{Trabajo}
  child{ node[font=\scriptsize]{empleo} }
  child{ node[font=\scriptsize]{desempleo} }
  child{ node[font=\scriptsize]{empresa} }
  child{ node[font=\scriptsize]{comercio justo} }
};
\end{tikzpicture}
\legende{Carte mentale des quatre champs lexicaux de la leçon.}
\end{popfigure}

\courssub{Tableau récapitulatif : pièges et correspondances}

\begin{center}
\begin{tabularx}{\linewidth}{|X|X|X|}
\hline
\rowcolor{popblueL} Español & Français & Remarque / piège \\
\hline
\esp{la inversión} & l'investissement & $\neq$ invention (\esp{invento}) \\
\hline
\esp{invertir (invierto)} & investir & e $\rightarrow$ ie au présent \\
\hline
\esp{el paro} & le chômage & surtout en Espagne \\
\hline
\esp{el desempleo} & le chômage & terme international \\
\hline
\esp{la asistencia} & la présence & « aider » = \esp{ayudar} \\
\hline
\esp{la asistencia social} & l'aide sociale & ne pas confondre \\
\hline
\esp{conceder} & accorder & \esp{conceder una subvención} \\
\hline
\esp{el préstamo} & le prêt & à rembourser \\
\hline
\esp{la subvención} & la subvention & non remboursable \\
\hline
\esp{la ONG} & l'ONG & toujours féminin \\
\hline
\esp{firmar un acuerdo} & signer un accord & sans préposition \\
\hline
\esp{el comercio justo} & le commerce équitable & \esp{precio justo} \\
\hline
\end{tabularx}
\end{center}

\courssub{Exercice résolu}

\begin{exoresolu}[Traduction thème : l'aide au développement]
Traduisez en espagnol : « Le gouvernement camerounais a signé un accord bilatéral avec l'Espagne. L'Espagne accorde une subvention pour financer la formation professionnelle des jeunes. Les ONG travaillent sur le terrain pour lutter contre la pauvreté et créer un emploi durable. »
\tcblower
\textbf{Solution détaillée.}

\esp{El gobierno camerunés ha firmado un acuerdo bilateral con España. España concede una subvención para financiar la formación profesional de los jóvenes. Las ONG trabajan sobre el terreno para luchar contra la pobreza y crear empleo duradero.}

Points de méthode :
\begin{enumerate}
\item « camerounais » (adjectif) = \esp{camerunés, camerunesa} ; le pays = \esp{Camerún} (avec accent).
\item « a signé » : passé composé espagnol \esp{ha firmado} (action récente et liée au présent).
\item « accorde une subvention » : \esp{conceder una subvención} — ne jamais utiliser \esp{acordar} (qui signifie « convenir, décider »).
\item « sur le terrain » : expression figée \esp{sobre el terreno} (et non \esp{en el terreno}).
\item « pour + infinitif » de but = \esp{para + infinitivo} : \esp{para financiar}, \esp{para luchar}.
\end{enumerate}
\end{exoresolu}

\begin{exercice}[À vous de jouer]
\begin{enumerate}
\item Complétez avec le mot qui convient (\esp{subvención, préstamo, inversión, donativo}) : \esp{El banco concedió un \ldots{} a diez años.} ; \esp{El Estado concedió una \ldots{} a la escuela.} ; \esp{La \ldots{} extranjera crea empleo.}
\item Traduisez : « Les Nations unies coordonnent l'aide humanitaire après les catastrophes. »
\item Vrai ou faux : \esp{la asistencia} signifie « l'aide ». Justifiez.
\end{enumerate}
\end{exercice}

\begin{corrige}[Exercice « À vous de jouer »]
\begin{enumerate}
\item \esp{préstamo} (à rembourser sur dix ans) ; \esp{subvención} (argent donné par l'État) ; \esp{inversión}.
\item \esp{Las Naciones Unidas coordinan la ayuda humanitaria después de las catástrofes.}
\item Faux : \esp{la asistencia} signifie d'abord « la présence » ; « l'aide » se dit \esp{la ayuda} et « l'aide sociale » \esp{la asistencia social}.
\end{enumerate}
\end{corrige}

\begin{aretenir}[L'essentiel du lexique]
\begin{itemize}
\item Quatre champs : \cle{financiación} (\esp{subvención, préstamo, inversión, donativo}), \cle{acuerdos} (\esp{bilateral, multilateral, tratado, cooperación}), \cle{acteurs} (\esp{la ONG, el Estado, las Naciones Unidas, el voluntario}), \cle{trabajo} (\esp{empleo, desempleo/paro, empresa, comercio justo}).
\item Expressions figées : \esp{trabajar sobre el terreno}, \esp{luchar contra la pobreza}, \esp{crear empleo duradero}, \esp{firmar un acuerdo}.
\item Pièges majeurs : \esp{inversión} $\neq$ invention ; \esp{asistencia} $\neq$ aide ; \esp{el paro} = chômage (Espagne) ; \esp{la ONG} est féminin ; \esp{invertir} $\rightarrow$ \esp{invierto}.
\end{itemize}
\end{aretenir}

Ce lexique maîtrisé, nous sommes prêts pour l'étape suivante : apprendre la \emph{méthode du commentaire de texte}, qui nous permettra d'organiser ces connaissances en une analyse structurée et rédigée.

\coursec{Méthode du commentaire de texte}

\courssub{Transition depuis le lexique thématique}

Après avoir maîtrisé le vocabulaire spécifique de l'aide au développement et du monde du travail (\emph{subvención, acuerdo bilateral, ONG, inversión, cooperación}), vous disposez des outils lexicaux nécessaires pour aborder sereinement l'exercice roi du baccalauréat : le commentaire de texte. Ce n'est pas une simple paraphrase, ni une dissertation hors sujet. C'est une \textbf{analyse structurée} qui prouve que vous avez compris le texte, que vous savez l'organiser et que vous maîtrisez la langue espagnole (grammaire, connecteurs, registres). La méthode ci-dessous, rodée pour l'épreuve du Bac camerounais (LV2), vous guide pas à pas.

\begin{textebox}{Texte d'entraînement : « Retos y oportunidades del empleo juvenil en África »}
África es el continente más joven del mundo: el 60\% de su población tiene menos de 25 años. Este « \textbf{dividendo demográfico} » representa una oportunidad histórica, pero también un desafío urgente. Cada año, más de diez millones de jóvenes entran en un mercado laboral que no genera suficientes empleos formales. La economía informal absorbe al 80\% de ellos, sin contratos, sin seguridad social y con salarios de miseria.

Ante esta realidad, la cooperación internacional juega un papel clave. Los acuerdos bilaterales entre la Unión Europea y la Unión Africana, así como las subvenciones de organismos como el BID o el Banco Mundial, financian programas de formación profesional y de apoyo al emprendimiento. Sin embargo, los expertos advierten: la ayuda al desarrollo no basta si no va acompañada de reformas estructurales. La corrupción, la inestabilidad política y la falta de infraestructuras frenan la inversión productiva.

En Camerún, por ejemplo, el Plan Nacional de Desarrollo 2020-2030 prioriza la industrialización y la transformación del sector primario. Las ONG locales, apoyadas por fondos multilaterales, crean incubadoras de start-ups en Yaoundé y Douala. Pero los jóvenes reclaman más: quieren una educación que responda a las necesidades reales del mercado, y exigen transparencia en la gestión de los fondos públicos.

En conclusión, el futuro laboral de África no depende solo de la ayuda externa, sino de la capacidad de sus gobiernos para crear un entorno favorable a la empresa y a la innovación. El talento africano existe; falta el marco que le permita florecer.
\end{textebox}

\begin{definition}[Glossaire du texte d'entraînement]
\textbf{Dividendo demográfico} : dividende démographique (structure d'âge favorable). \\
\textbf{Empleo formal / informal} : emploi formel / informel. \\
\textbf{Seguridad social} : sécurité sociale. \\
\textbf{Salario de miseria} : salaire de misère. \\
\textbf{BID} : Banco Interamericano de Desarrollo (Banque interaméricaine de développement). \\
\textbf{Emprendimiento} : entrepreneuriat. \\
\textbf{Reformas estructurales} : réformes structurelles. \\
\textbf{Incubadora de start-ups} : incubateur de start-ups. \\
\textbf{Marco favorable} : cadre favorable.
\end{definition}

\courssub{Étape 1 — Lecture active et repérage des mots-clés}

\begin{methode}[Lecture active : deux lectures obligatoires]
\begin{itemize}
\item \textbf{Première lecture (globale) :} Lisez le texte d'une traite, sans dictionnaire. Identifiez : la nature du texte (article de presse, essai, rapport), le thème général (le monde du travail, l'aide au développement), le ton (alarmiste, optimiste, analytique) et l'intention de l'auteur (informer, alerter, convaincre).
\item \textbf{Deuxième lecture (détaillée, crayon en main) :} Soulignez ou surlignez :
\begin{itemize}
    \item Les \cle{mots-clés} récurrents (empleo, juventud, ayuda, inversión, África, gobierno).
    \item Les \cle{connecteurs logiques} (sin embargo, por tanto, en conclusión, ante esta realidad) qui révèlent la structure interne.
    \item Les \cle{chiffres} et \cle{exemples concrets} (60\%, 10 millions, Camerún, Yaoundé, Douala).
    \item Les \cle{opinions / jugements} de l'auteur (expertos advierten, los jóvenes reclaman, falta el marco).
\end{itemize}
\item \textbf{Troisième lecture (de vérification) :} Relisez en vous posant la question : « Quelle est la thèse défendue ? » (Réponse attendue : L'aide extérieure est nécessaire mais insuffisante sans bonnes gouvernances et réformes internes).
\end{itemize}
\end{methode}

\begin{attention}[Piège fréquent : la lecture passive]
Ne vous contentez pas de « comprendre le sens global ». Au Bac, on vous demande d'\textbf{analyser}, pas de résumer. Si vous ne soulignez pas les connecteurs (\emph{sin embargo, por tanto}), vous ne verrez pas la structure argumentative (thèse / antithèse / synthèse) et votre développement sera plat.
\end{attention}

\courssub{Étape 2 — Définir le thème et la thèse}

\begin{definition}[Thème vs Thèse]
\begin{itemize}
    \item \textbf{Le thème (el tema)} : Sujet général du texte, en quelques mots. \emph{Exemple :} « El empleo juvenil en África y la cooperación internacional. »
    \item \textbf{La thèse (la tesis)} : Idée directrice, point de vue de l'auteur sur le thème. C'est une phrase affirmative. \emph{Exemple :} « La ayuda al desarrollo es indispensable pero insuficiente sin reformas políticas y económicas internas. »
\end{itemize}
\end{definition}

\begin{exemplebox}{Application sur le texte d'entraînement}
\textbf{Thème :} Los retos del empleo juvenil en África y el papel de la ayuda internacional. \\
\textbf{Thèse :} El autor sostiene que la ayuda exterior (bilateral, multilateral, ONG) es necesaria pero no basta; se requieren reformas estructurales, buena gobernanza y educación adaptada para transformar el dividendo demográfico en desarrollo real.
\end{exemplebox}

\courssub{Étape 3 — Rédiger l'introduction : le « contrat » avec le correcteur}

L'introduction suit un schéma immuable en 4 temps (méthode « entonnoir ») :

\begin{methode}[Structure type de l'introduction (4 étapes)]
\begin{itemize}
\item \textbf{Accroche / Présentation du document :} Nature, source (si connue), date, auteur (si connu), thème général.
\item \textbf{Problématique / Thèse :} Reformulez la thèse de l'auteur en une phrase claire.
\item \textbf{Annonce du plan :} Présentez les 2 ou 3 grandes parties du développement en utilisant les connecteurs de structure.
\item \textbf{Phrase de transition :} Lance le développement.
\end{itemize}
\end{methode}

\begin{exemplebox}{Modèle d'introduction (à adapter)}
\esp{Este texto es un artículo de prensa que trata del mundo laboral y de la ayuda al desarrollo en África. El autor explica que, aunque el continente tiene una población joven muy numerosa, el mercado laboral no crea suficientes empleos formales. La tesis principal es que la cooperación internacional y las subvenciones son necesarias, pero no son suficientes sin reformas estructurales y buena gobernanza.}

\esp{Vamos a analizar primero los desafíos del empleo juvenil y la economía informal. Luego, estudiaremos el papel de la ayuda internacional y sus límites. Finalmente, veremos las soluciones propuestas: educación, emprendimiento y transparencia.}
\end{exemplebox}

\begin{propriete}[Formules types pour l'introduction]
\begin{center}
\begin{tabularx}{\linewidth}{|X|X|}\hline
\rowcolor{popblueL} \textbf{Fonction} & \textbf{Formule espagnole} \\ \hline
Présenter le texte & \esp{Este texto es un artículo / informe / ensayo que trata de...} \\ \hline
Thème & \esp{El tema central es... / El texto aborda la problemática de...} \\ \hline
Thèse & \esp{El autor sostiene que... / La tesis defendida es que... / El autor denuncia / destaca / propone...} \\ \hline
Annonce du plan (2 parties) & \esp{En primer lugar, analizaremos... ; en segundo lugar, examinaremos...} \\ \hline
Annonce du plan (3 parties) & \esp{Vamos a estudiar primero..., luego..., y finalmente...} \\ \hline
Transition & \esp{Comencemos por... / Para empezar...} \\ \hline
\end{tabularx}
\end{center}
\end{propriete}

\courssub{Étape 4 — Le développement : structure, citation, explication}

C'est le cœur de la note. Respectez la règle d'or : \textbf{Une partie = Une idée directrice = Un paragraphe (ou deux) structuré}.

\begin{methode}[Architecture d'une partie de développement (méthode ICE)]
\begin{itemize}
\item \textbf{I — Idée directrice (phrase sujet) :} Annoncez l'idée en vos mots (espagnol).
\item \textbf{C — Citation / Référence précise :} Citez un segment court du texte (entre guillemets \esp{« ... »}) ou référez-vous à un exemple précis (\esp{el ejemplo de Camerún, la cifra del 60\%}). \textbf{Jamais de citation longue.}
\item \textbf{E — Explication / Analyse :} Expliquez \emph{pourquoi} cette citation prouve l'idée, faites le lien avec le thème, utilisez votre connaissance du contexte (Cameroun, Guinée équatoriale, ODD). Utilisez des connecteurs logiques.
\end{itemize}
\end{methode}

\begin{exemplebox}{Développement partie 1 (Idée 1 : Le défi démographique et l'informel)}
\textbf{I :} \esp{En primer lugar, el texto destaca \textbf{la} paradoja del dividendo demográfico africano.} \\
\textbf{C :} \esp{El autor señala que « el 60\% de su población tiene menos de 25 años » y que « la economía informal absorbe al 80\% » de los jóvenes.} \\
\textbf{E :} \esp{Esto significa que la juventud, que debería ser un motor de crecimiento, se convierte en una carga social por falta de empleo formal. En Camerún, como en muchos países de la CEMAC, los jóvenes titulados conducen motos-taxis (bendskins) por falta de alternativas. La ausencia de seguridad social y de salarios dignos perpetúa la pobreza.}
\end{exemplebox}

\begin{exemplebox}{Développement partie 2 (Idée 2 : L'aide internationale, nécessaire mais limitée)}
\textbf{I :} \esp{En segundo lugar, se analiza el papel de la cooperación internacional.} \\
\textbf{C :} \esp{El texto menciona « acuerdos bilaterales », « subvenciones del BID o el Banco Mundial » y « ONG locales apoyadas por fondos multilaterales ».} \\
\textbf{E :} \esp{Estas ayudas financian formación e incubadoras (ej. Yaoundé, Douala). Sin embargo, el autor introduce una concesión adversativa clave: « \textbf{sin embargo}, los expertos advierten: la ayuda al desarrollo no basta si no va acompañada de reformas estructurales ». La corrupción y la inestabilidad son « frenos » para la inversión productiva.}
\end{exemplebox}

\begin{exemplebox}{Développement partie 3 (Idée 3 : Solutions endogènes et gouvernance)}
\textbf{I :} \esp{Finalmente, el autor apuesta por soluciones internas y participación ciudadana.} \\
\textbf{C :} \esp{« Los jóvenes reclaman más: quieren una educación que responda a las necesidades reales del mercado » y « exigen transparencia ».} \\
\textbf{E :} \esp{El Plan Nacional de Desarrollo 2020-2030 de Camerún ilustra esta vía: industrialización y transformación del sector primario. La tesis final es clara: « el futuro laboral de África no depende solo de la ayuda externa, sino de la capacidad de sus gobiernos ». El talento existe; falta el marco institucional.}
\end{exemplebox}

\begin{propriete}[Connecteurs indispensables pour le développement]
\begin{center}
\begin{tabularx}{\linewidth}{|X|X|X|}\hline
\rowcolor{popblueL} \textbf{Fonction} & \textbf{Connecteurs (Español)} & \textbf{Français} \\ \hline
Ordre / Énumération & \esp{En primer lugar, en segundo lugar, finalmente, para empezar, luego} & En premier lieu, en second lieu, enfin, pour commencer, puis \\ \hline
Addition & \esp{Además, también, asimismo, de igual manera} & De plus, aussi, de même \\ \hline
Opposition / Concession & \esp{Sin embargo, no obstante, por un lado / por otro lado, aunque} & Cependant, néanmoins, d'un côté / de l'autre, bien que \\ \hline
Conséquence & \esp{Por tanto, por consiguiente, como resultado, así que} & Par conséquent, donc, si bien que \\ \hline
Cause & \esp{Puesto que, ya que, debido a, a causa de} & Puisque, car, à cause de \\ \hline
Exemple / Illustration & \esp{Por ejemplo, como muestra el caso de, en concreto} & Par exemple, comme le montre le cas de, concrètement \\ \hline
Conclusion partielle & \esp{En definitiva, en resumen, en suma} & En définitive, en résumé \\ \hline
\end{tabularx}
\end{center}
\end{propriete}

\courssub{Étape 5 — La conclusion : bilan et ouverture}

La conclusion ne doit pas être une simple répétition. Elle suit un schéma en 3 temps.

\begin{methode}[Structure de la conclusion (3 étapes)]
\begin{itemize}
\item \textbf{Bilan (synthèse) :} Reprenez la thèse en une phrase, résumez les 2-3 grandes idées du développement sans ajouter d'information nouvelle.
\item \textbf{Opinion personnelle argumentée :} Donnez votre avis \textbf{en espagnol}, en vous appuyant sur vos connaissances (contexte camerounais, ODD, actualité). Utilisez \esp{En mi opinión}, \esp{A mi juicio}, \esp{Considero que}.
\item \textbf{Ouverture (perspective) :} Posez une question, élargissez à un autre domaine (éducation, genre, changement climatique), ou proposez une piste de solution future.
\end{itemize}
\end{methode}

\begin{exemplebox}{Modèle de conclusion}
\esp{En conclusión, el texto demuestra que el empleo juvenil en África es un reto estructural: la ayuda internacional (subvenciones, acuerdos bilaterales, ONG) alivia la situación pero no la resuelve sin voluntad política interna. Hemos visto que la economía informal domina, que la cooperación tiene límites y que la clave está en la educación, el emprendimiento y la transparencia.}

\esp{\textbf{En mi opinión}, la prioridad absoluta debe ser la reforma del sistema educativo para alinear la formación con el mercado laboral real, como pide el texto. Además, la lucha contra la corrupción no es un eslogan, es una condición \emph{sine qua non} para que la inversión (nacional o extranjera) cree empleo digno.}

\esp{\textbf{Para terminar}, cabe preguntarse si la Zona de Libre Comercio Continental Africana (ZLECAf / AfCFTA) logrará crear el « marco favorable » que reclama el autor, o si los jóvenes africanos seguirán esperando un empleo que tarde en llegar.}
\end{exemplebox}

\begin{attention}[Erreurs fatales au Bac : Calques et fautes de grammaire]
\begin{itemize}
    \item \textbf{« Según a el autor »} $\rightarrow$ \cle{INCORRECT}. On dit : \esp{Según el autor} (pas de \emph{a} devant \emph{el}, pas d'article contracté \emph{al} après préposition \emph{según}).
    \item \textbf{« En mi opinion »} $\rightarrow$ \cle{INCORRECT} (accent manquant). On écrit : \esp{En mi opinión}.
    \item \textbf{« Pienso que la ayuda es... »} vs \textbf{« No pienso que la ayuda sea... »} $\rightarrow$ \textbf{Règle d'or :} Après \esp{pienso que / creo que / estoy seguro que} + \textbf{INDICATIF} (certitude). Après \esp{no pienso que / no creo que / dudo que / es posible que} + \textbf{SUBJONCTIF} (doute/négation).
    \item \textbf{« La ayuda es importante para desarrollar el país »} $\rightarrow$ \esp{Para desarrollar} (but) vs \esp{Por desarrollar} (cause). \cle{Por / Para} : voir fiche spécifique Module 4.
    \item \textbf{Traduction littérale de « Il faut que »} : \esp{Hay que + infinitif} (impersonnel) ou \esp{Es necesario que + subjonctif}. Pas \esp{Il faut que} $\rightarrow$ \esp{*Il faut que}.
\end{itemize}
\end{attention}

\courssub{Critères d'évaluation au Baccalauréat (Grille officielle MINESEC)}

\begin{center}
\begin{tabularx}{\linewidth}{|X|X|}\hline
\rowcolor{popblueL} \textbf{Critère} & \textbf{Indicateurs de réussite (Note sur 20)} \\ \hline
\cle{Fidélité au texte} (5 pts) & Respect du sens, pas de contresens, citations exactes et pertinentes, pas d'invention. \\ \hline
\cle{Structure} (5 pts) & Introduction canonique (4 temps), développement équilibré (2-3 parties ICE), conclusion (3 temps), connecteurs variés et justes. \\ \hline
\cle{Correction linguistique} (5 pts) & Accords (GN), conjugaisons (indicatif/subjonctif si/que, valeurs des temps), syntaxe (place de l'adjectif, pronoms), orthographe (accents, ñ, ¿ ¡). \\ \hline
\cle{Richesse lexicale} (5 pts) & Vocabulaire précis du module (subvención, acuerdo bilateral, brecha, emprendimiento, gobernanza), pas de répétitions, registre soutenu, pas de français (ex. \emph{actualmente} au lieu de \emph{en la actualidad} est ok, mais \emph{problema} sans accent non). \\ \hline
\end{tabularx}
\end{center}

\begin{savaistu}[Le commentaire au Cameroun et en Guinée équatoriale]
Au Cameroun, l'espagnol est LV2. L'épreuve écrite (coefficient variable selon séries) comporte souvent un commentaire de texte sur un thème du programme (Module 4 : Monde du travail / Développement). En Guinée équatoriale (pays hispanophone voisin), le commentaire de texte (\emph{comentario de texto}) est un exercice universitaire fondamental (Selectividad / EBAU). La méthode y est identique : \emph{Introducción, Tema, Tesis, Estructura, Análisis lingüístico, Conclusión}. Maîtriser cette méthode au lycée vous donne un avantage majeur si vous poursuivez des études en zone hispanophone (Espagne, Mexique, Argentine, Guinée équatoriale).
\end{savaistu}

\courssub{Organigramme récapitulatif de la structure du commentaire}

\begin{popfigure}
\begin{tikzpicture}[
    node distance=1.2cm and 2cm,
    every node/.style={align=center, font=\small},
    box/.style={draw=popblue, thick, rounded corners, fill=popblueL, minimum width=3.5cm, minimum height=1.2cm},
    arrow/.style={-Stealth, thick, popdark}
]
% Nodes
\node[box, align=center] (intro){\textbf{INTRODUCTION}\\(4 étapes)\\ature, source, thème\\thèse de l'auteur\\annonce du plan (2-3 parties)\\transition};
\node[box, below=of intro, align=center] (dev1){\textbf{DÉVELOPPEMENT}\\\textbf{Partie 1 : Idée 1}\\I : Phrase sujet\\C : Citation courte\\E : Explication + connecteurs};
\node[box, below=of dev1, align=center] (dev2){\textbf{Partie 2 : Idée 2}\\I / C / E\\\emph{(Opposition / Nuance)}\\\esp{Sin embargo, no obstante}};
\node[box, below=of dev2, align=center] (dev3){\textbf{Partie 3 : Idée 3}\\I / C / E\\\emph{(Solution / Avenir)}\\\esp{Finalmente, en definitiva}};
\node[box, below=of dev3, align=center] (concl){\textbf{CONCLUSION}\\(3 étapes)\\Bilan (thèse + synthèse)\\Opinion \esp{En mi opinión / A mi juicio}\\Ouverture (question / perspective)};

% Arrows
\draw[arrow] (intro.south) -- (dev1.north) node[midway, right, xshift=0.5cm, font=\tiny, text=poporange] {\esp{Comencemos por...}};
\draw[arrow] (dev1.south) -- (dev2.north) node[midway, right, xshift=0.5cm, font=\tiny, text=poporange] {\esp{Luego / Por otro lado}};
\draw[arrow] (dev2.south) -- (dev3.north) node[midway, right, xshift=0.5cm, font=\tiny, text=poporange] {\esp{Finalmente}};
\draw[arrow] (dev3.south) -- (concl.north) node[midway, right, xshift=0.5cm, font=\tiny, text=poporange] {\esp{En conclusión}};
\end{tikzpicture}
\legende{Structure type du commentaire de texte (Bac LV2) : Introduction $\rightarrow$ Développement (3 parties ICE) $\rightarrow$ Conclusion. Les flèches indiquent les connecteurs de transition obligatoires.}
\end{popfigure}

\courssub{Exercice résolu : De l'analyse à l'introduction}

\begin{exoresolu}[Rédiger l'introduction à partir d'un court extrait]
\textbf{Énoncé :} À partir de l'extrait suivant, rédigez une introduction complète (4 étapes) en espagnol. Respectez le modèle.

\esp{« La brecha digital en el mundo laboral latinoamericano se ha ampliado tras la pandemia. Según la OIT, el 70\% de los nuevos empleos requieren competencias digitales básicas, pero solo el 30\% de la fuerza laboral las posee. Los gobiernos deben invertir en formación continua y en infraestructura tecnológica rural. Sin la participación del sector privado, las subvenciones públicas no llegarán a los más vulnerables. »}

\tcblower
\textbf{Solution détaillée :}

\textbf{1. Analyse rapide :}
\begin{itemize}
    \item \textbf{Nature :} Article / Rapport (style analytique, chiffre OIT).
    \item \textbf{Thème :} Brecha digital / empleo / formación en Latinoamérica.
    \item \textbf{Thèse :} La brecha digital s'agrandit ; l'État doit investir (formation, infra) mais le privé est indispensable.
    \item \textbf{Plan (2 parties) :} 1. Constat : brecha digital post-pandemia y falta de competencias. 2. Soluciones : inversión pública + alianza público-privada.
\end{itemize}

\textbf{2. Introduction rédigée :}
\esp{Este texto es un artículo de análisis económico que trata de la brecha digital en el mercado laboral latinoamericano tras la pandemia. El autor denuncia que el 70\% de los nuevos empleos exigen competencias digitales que solo el 30\% de los trabajadores \textbf{poseen}. La tesis central es que los gobiernos deben invertir masivamente en formación e infraestructura rural, pero advierte que sin la colaboración del sector privado las subvenciones no llegarán a los más vulnerables.}

\esp{Vamos a analizar primero el diagnóstico alarmante de la brecha digital y el déficit de competencias. Luego, examinaremos las soluciones propuestas: la inversión pública en formación y la necesaria alianza con el sector privado. Comencemos por el diagnóstico.}
\end{exoresolu}

\begin{exercice}[Entraînement : À vous de jouer !]
\textbf{Consigne :} Reprenez le texte d'entraînement de cette leçon (\emph{« Retos y oportunidades del empleo juvenil en África »}).
\begin{enumerate}
    \item Identifiez les 3 idées directrices pour le développement (une par paragraphe du texte).
    \item Rédigez la \textbf{Partie 2 du développement} (L'aide internationale : rôle et limites) en appliquant strictement la méthode \textbf{ICE} (Idée - Citation - Explication). Utilisez au moins 3 connecteurs du tableau.
    \item Rédigez la \textbf{conclusion complète} (Bilan + Opinion + Ouverture) en vous inspirant du modèle mais en utilisant vos propres arguments (contexte camerounais : PND 2030, ZLECAf, formation professionnelle).
\end{enumerate}
\textit{Écrivez tout en espagnol. Vérifiez : subjonctif après opinion négative, accents, connecteurs, vocabulaire du module.}
\end{exercice}

\begin{corrige}[Exercice Entraînement]
\textbf{1. Idées directrices :}
\begin{itemize}
    \item Idée 1 (§1) : Le paradoxe du dividende démographique (jeunesse nombreuse mais emploi informel précaire).
    \item Idée 2 (§2) : L'aide internationale (bilatérale, multilatérale, ONG) nécessaire mais insuffisante sans réformes (corruption, instabilité).
    \item Idée 3 (§3-4) : Solutions endogènes : éducation adaptée, entrepreneuriat, transparence, bonne gouvernance (ex. Cameroun).
\end{itemize}

\textbf{2. Partie 2 (Modèle ICE) :}
\esp{En segundo lugar, el texto examina el papel de la cooperación internacional como respuesta a la crisis. \textbf{Por un lado}, reconoce su utilidad: « los acuerdos bilaterales entre la Unión Europea y la Unión Africana, así como las subvenciones de organismos como el BID o el Banco Mundial, financian programas de formación profesional ». \textbf{Además}, cita el caso concreto de « ONG locales, apoyadas por fondos multilaterales, crean incubadoras de start-ups en Yaoundé y Douala ». \textbf{Sin embargo}, \textbf{por otro lado}, el autor introduce una limitación fundamental mediante una concesión adversativa: « los expertos advierten: la ayuda al desarrollo no basta si no va acompañada de reformas estructurales ». \textbf{Por consiguiente}, la corrupción y la inestabilidad actúan como frenos que impiden que la inversión sea productiva. \textbf{En definitiva}, la ayuda es una condición necesaria pero no suficiente.}

\textbf{3. Conclusion (Modèle) :}
\esp{En conclusión, el texto demuestra que el futuro laboral de la juventud africana pende de un hilo entre la oportunidad demográfica y la realidad del empleo informal. Hemos visto que la economía informal domina, que la cooperación tiene límites y que la clave está en la educación pertinente, el emprendimiento apoyado y, sobre todo, buena gobernanza.}

\esp{\textbf{En mi opinión}, el Plan Nacional de Desarrollo de Camerún 2020-2030 va en la buena dirección al priorizar la industrialización, pero la clave está en la ejecución presupuestaria transparente. La ZLECAf ofrece un mercado único enorme; si los jóvenes cameruneses acceden a una formación técnica de calidad (ej. INP, IUT, \textbf{centros de formación profesional}), podrán exportar servicios y no solo materias primas.}

\esp{\textbf{Para terminar}, la pregunta sigue abierta: ¿lograrán los gobiernos africanos transformar la « \textbf{youth bulge} » (bonanza demográfica) en motor de crecimiento antes de que la frustración social derive en inestabilidad? La respuesta depende menos de la ayuda que de la voluntad política.}
\end{corrige}

\begin{aretenir}[Check-list finale avant l'épreuve]
\begin{itemize}
    \item [ ] J'ai lu le texte 3 fois (globale, détaillée, vérification).
    \item [ ] J'ai souligné : mots-clés, connecteurs, chiffres, opinions.
    \item [ ] J'ai formulé le \textbf{Thème} (nominal) et la \textbf{Thèse} (phrase).
    \item [ ] Mon \textbf{Introduction} a les 4 étapes (Nature/Source/Thème $\rightarrow$ Thèse $\rightarrow$ Plan $\rightarrow$ Transition).
    \item [ ] Mon \textbf{Développement} a 2 ou 3 parties \textbf{ICE} (Idée + Citation courte + Explication).
    \item [ ] J'ai utilisé des \textbf{connecteurs variés} (orden, oposición, consecuencia, ejemplo).
    \item [ ] Ma \textbf{Conclusion} a 3 temps (Bilan $\rightarrow$ Opinion \esp{En mi opinión} $\rightarrow$ Ouverture).
    \item [ ] \textbf{Relue :} Accords, Subjonctif (après \emph{no creo que, es necesario que}), Accents (\emph{opinión, también, además}), \emph{Por/Para}, Ponctuation \esp{¿ ¡}.
\end{itemize}
\end{aretenir}

\coursec{Commentaire modèle guidé}

Dans la section précédente, nous avons construit pas à pas la \cle{méthode} du commentaire de texte : lire deux fois, dégager le thème, bâtir un plan en trois parties, citer, conclure avec une opinion personnelle. Il est temps maintenant de voir cette méthode \emph{en action} : nous allons lire un \cle{comentario modelo} rédigé sur notre texte support \esp{El mundo laboral hoy}, le disséquer ensemble, puis le comparer à un commentaire faible pour en tirer les critères de réussite à l'examen.

\courssub{Lecture du commentaire modèle}

\begin{textebox}[Comentario modelo — \esp{El mundo laboral hoy}]
Este texto es un artículo de prensa que trata del mundo laboral actual y de la ayuda al desarrollo en África. En primer lugar, el autor describe la situación de los jóvenes africanos que « buscan un empleo digno ». Luego, explica los diferentes tipos de ayuda: las subvenciones, las inversiones y los acuerdos bilaterales y multilaterales, como el acuerdo entre España y Camerún. Además, subraya el papel de las ONG que « trabajan sobre el terreno ». Sin embargo, según algunos economistas, África necesita también « un comercio justo y más inversiones productivas ». En conclusión, este texto muestra que el verdadero desarrollo debe crear empleo duradero. En mi opinión, la cooperación es útil, pero los países africanos deben también desarrollar sus propias empresas.
\end{textebox}

\textbf{Traduction.} « Ce texte est un article de presse qui traite du monde du travail actuel et de l'aide au développement en Afrique. D'abord, l'auteur décrit la situation des jeunes Africains qui "cherchent un emploi digne". Ensuite, il explique les différents types d'aide : les subventions, les investissements et les accords bilatéraux et multilatéraux, comme l'accord entre l'Espagne et le Cameroun. De plus, il souligne le rôle des ONG qui "travaillent sur le terrain". Cependant, selon certains économistes, l'Afrique a aussi besoin d'"un commerce juste et de davantage d'investissements productifs". En conclusion, ce texte montre que le véritable développement doit créer un emploi durable. À mon avis, la coopération est utile, mais les pays africains doivent aussi développer leurs propres entreprises. »

\textbf{Glossaire.}

\begin{center}
\begin{tabularx}{\linewidth}{|X|X|}\hline
\rowcolor{popblueL} \textbf{Español} & \textbf{Français} \\ \hline
tratar (sobre / de) / versar sobre & traiter de, porter sur \\ \hline
el empleo digno & l'emploi digne \\ \hline
la subvención & la subvention \\ \hline
el acuerdo bilateral / multilateral & l'accord bilatéral / multilatéral \\ \hline
subrayar & souligner \\ \hline
sobre el terreno & sur le terrain \\ \hline
el comercio justo & le commerce équitable \\ \hline
duradero, -a & durable \\ \hline
desarrollar & développer \\ \hline
las propias empresas & leurs propres entreprises \\ \hline
\end{tabularx}
\end{center}

\begin{attention}[Point de vigilance : \esp{tratar de}]
Le verbe \esp{tratar} suivi de la préposition \esp{de} + nom (\esp{trata del mundo laboral}) signifie « traiter de, porter sur ». Mais \esp{tratar de} + infinitif signifie « essayer de » (\esp{trato de estudiar} = j'essaie d'étudier). Ne confondez pas les deux constructions.
\end{attention}

\textbf{Questions d'observation collective.}

\begin{enumerate}
\item Combien de parties distinguez-vous dans ce commentaire ? Soulignez les mots qui les annoncent.
\item Combien de citations du texte support sont utilisées ? Entre quels signes sont-elles placées ?
\item Où se trouve l'opinion personnelle du rédacteur ? Quelle expression l'introduit ?
\end{enumerate}

\courssub{Identification des trois parties et des citations}

\begin{propriete}[La structure du comentario modelo]
Le commentaire modèle respecte le plan en \cle{trois parties} vu dans la méthode :
\begin{enumerate}
\item \textbf{Presentación} : \esp{Este texto es un artículo de prensa que trata del mundo laboral actual y de la ayuda al desarrollo en África.} — on présente la \emph{nature} du texte (article de presse) et son \emph{thème}.
\item \textbf{Análisis de las ideas} : le corps du commentaire, rythmé par \esp{En primer lugar} … \esp{Luego} … \esp{Además} … \esp{Sin embargo} … — chaque idée du texte est reprise et illustrée par une citation.
\item \textbf{Conclusión y opinión personal} : \esp{En conclusión} … puis \esp{En mi opinión} … — on synthétise, puis on donne son avis.
\end{enumerate}
\end{propriete}

\begin{exemplebox}[Les citations du commentaire modèle]
Trois citations appuient l'analyse :
\begin{itemize}
\item \esp{« buscan un empleo digno »} → illustre la situation des jeunes Africains ;
\item \esp{« trabajan sobre el terreno »} → illustre le rôle concret des ONG ;
\item \esp{« un comercio justo y más inversiones productivas »} → illustre le point de vue des économistes.
\end{itemize}
Remarquez que chaque citation est \emph{courte} (quelques mots), placée entre \esp{« »}, et toujours \emph{expliquée} par la phrase qui l'entoure.
\end{exemplebox}

\begin{methode}[Comment citer un texte dans un commentaire]
Pour introduire une citation, utilisez des \cle{verbes de parole} et des formules variées :
\begin{enumerate}
\item \esp{El autor dice que...} — « L'auteur dit que... »
\item \esp{Según el texto, ...} — « Selon le texte, ... »
\item \esp{Como escribe el autor: « ... »} — « Comme l'auteur l'écrit : "..." »
\item \esp{El texto afirma / subraya / explica que...} — « Le texte affirme / souligne / explique que... »
\end{enumerate}
Après la citation, ajoutez toujours une phrase de commentaire : \esp{Estas palabras muestran que...} (« Ces mots montrent que... »), \esp{Con esta expresión, el autor quiere decir que...} (« Par cette expression, l'auteur veut dire que... »).
\end{methode}

\begin{attention}[Piège : la citation « orpheline »]
Une \cle{citation} ne doit \textbf{jamais} être laissée seule, sans introduction ni explication. Écrire simplement : \esp{Los jóvenes africanos. « Buscan un empleo digno ». Las ONG.} n'a aucune valeur : le correcteur ne sait pas si vous avez compris. Règle d'or : \textbf{avant} la citation, une formule d'introduction ; \textbf{après} la citation, une phrase qui l'explique avec vos propres mots. Autre piège de francophone : ne recopiez pas de longues phrases entières du texte — citez 3 à 8 mots, pas trois lignes.
\end{attention}

\courssub{Analyse des connecteurs et de la progression logique}

Relisons le modèle en soulignant ses \cle{conectores} : \esp{En primer lugar} → \esp{Luego} → \esp{Además} → \esp{Sin embargo} → \esp{En conclusión} → \esp{En mi opinión}. Ces petits mots sont la « charpente » du commentaire : ils montrent au lecteur où l'on va.

\begin{center}
\begin{tabularx}{\linewidth}{|X|X|X|}\hline
\rowcolor{popblueL} \textbf{Conector} & \textbf{Français} & \textbf{Fonction dans le commentaire} \\ \hline
\esp{En primer lugar} & D'abord, en premier lieu & annoncer la première idée \\ \hline
\esp{Luego / Después} & Ensuite & enchaîner l'idée suivante \\ \hline
\esp{Además / También} & De plus, également & ajouter un argument \\ \hline
\esp{Sin embargo / Pero} & Cependant, mais & introduire une nuance ou une opposition \\ \hline
\esp{Según (+ nom)} & Selon & attribuer une opinion à quelqu'un \\ \hline
\esp{En conclusión / En resumen} & En conclusion, en résumé & ouvrir la synthèse finale \\ \hline
\esp{En mi opinión / Pienso que / Creo que} & À mon avis / Je pense que & exprimer l'opinion personnelle \\ \hline
\end{tabularx}
\end{center}

\begin{aretenir}[La progression logique du modèle]
Le commentaire modèle suit une progression en \emph{entonnoir} : du général (présentation du texte) vers le particulier (les idées une à une), puis retour au général (conclusion), et enfin ouverture personnelle. L'opposition \esp{Sin embargo} est précieuse : elle montre que vous savez relever les \emph{nuances} du texte, ce qui distingue un bon commentaire d'un simple résumé.
\end{aretenir}

\courssub{Repérage de l'opinion personnelle dans la conclusion}

Dans le modèle, l'opinion personnelle occupe la toute dernière phrase : \esp{En mi opinión, la cooperación es útil, pero los países africanos deben también desarrollar sus propias empresas.} Observez sa construction : un marqueur d'opinion (\esp{En mi opinión}), puis un jugement \emph{nuancé} (\esp{es útil, pero...}) : le rédacteur reconnaît la valeur de la coopération tout en ajoutant une condition. C'est exactement ce qu'attend l'examinateur : un avis \emph{personnel}, \emph{argumenté} et \emph{équilibré}, pas un slogan.

\begin{exemplebox}[Autres formules d'opinion utiles]
\begin{itemize}
\item \esp{Pienso que la ayuda al desarrollo es necesaria, pero no suficiente.} « Je pense que l'aide au développement est nécessaire, mais pas suffisante. »
\item \esp{Me parece que los acuerdos bilaterales benefician a los dos países.} « Il me semble que les accords bilatéraux profitent aux deux pays. »
\item \esp{Creo que el texto tiene razón: sin empleo no hay desarrollo.} « Je crois que le texte a raison : sans emploi, pas de développement. »
\end{itemize}
\end{exemplebox}

\begin{attention}[Faux amis et calques dans l'opinion]
Ne dites pas \esp{en mi opinión personal} (pléonasme calqué du français « mon opinion personnelle ») : \esp{en mi opinión} suffit. Attention aussi à \esp{actualmente}, qui signifie « actuellement » et non « en fait » ; et à \esp{estoy de acuerdo con el autor} (« je suis d'accord avec l'auteur »), où le français « d'accord » ne se traduit jamais par \esp{de acorde}.
\end{attention}

\courssub{Comparaison avec un commentaire faible : les critères de réussite}

Voici maintenant, à titre de contre-exemple, un \cle{comentario débil} rédigé sur le même texte :

\begin{textebox}[Comentario débil (à ne pas imiter)]
El texto habla del trabajo en África. Los jóvenes no tienen empleo. Hay subvenciones y acuerdos. Las ONG ayudan. El texto es interesante y me gusta mucho.
\end{textebox}

Traduction : « Le texte parle du travail en Afrique. Les jeunes n'ont pas d'emploi. Il y a des subventions et des accords. Les ONG aident. Le texte est intéressant et il me plaît beaucoup. »

Ce texte n'est pas \emph{faux} sur le plan de la langue, mais il est \emph{pauvre} comme commentaire : pas de présentation du type de texte, pas de plan visible, aucune citation, aucun connecteur, et une « opinion » vide (\esp{me gusta mucho}) qui ne dit rien. Comparons les deux productions :

\begin{popfigure}
\centering
\begin{tikzpicture}[
  box/.style={draw=popline, rounded corners=4pt, align=left, text width=6.4cm, inner sep=6pt, font=\small},
  head/.style={font=\bfseries\small, align=center},
  lab/.style={font=\bfseries\small, text=popdark, align=center}
]
\node[head, text=popgreen] (h1) at (0,0) {Comentario fuerte};
\node[head, text=poppink] (h2) at (7.6,0) {Comentario débil};
\node[box, fill=popgreenL, anchor=north west] at (-3.4,-0.5)
  {\textbf{Plan :} tres partes claras (presentación, análisis, conclusión).\\[2pt]
   \textbf{Citas :} tres citas cortas entre « », siempre explicadas.\\[2pt]
   \textbf{Conectores :} \emph{En primer lugar, Luego, Además, Sin embargo, En conclusión}.\\[2pt]
   \textbf{Opinión :} personal y matizada (\emph{En mi opinión... pero...}).};
\node[box, fill=poppinkL, anchor=north west] at (4.2,-0.5)
  {\textbf{Plan :} ninguno: frases sueltas sin orden.\\[2pt]
   \textbf{Citas :} ninguna: no se ve el texto.\\[2pt]
   \textbf{Conectores :} ninguno: solo puntos.\\[2pt]
   \textbf{Opinión :} vacía (\emph{me gusta mucho}).};
\draw[-{Stealth}, very thick, poporange] (3.1,-3.4) -- (4.1,-3.4)
  node[midway, above, font=\small\bfseries, text=poporange, align=center]{¡Mejora!};
\end{tikzpicture}
\legende{Tableau comparatif : les quatre critères qui distinguent un commentaire fort d'un commentaire faible.}
\end{popfigure}

\begin{aretenir}[Les quatre critères de réussite d'un commentaire]
\begin{enumerate}
\item Un \cle{plan} en trois parties, annoncé par des connecteurs ;
\item Des \cle{citations} courtes, introduites et expliquées ;
\item Une \cle{progression logique} avec au moins une nuance (\esp{Sin embargo}) ;
\item Une \cle{opinión personal} argumentée et nuancée en conclusion.
\end{enumerate}
\end{aretenir}

\begin{savaistu}[Le commentaire au Baccalauréat A]
Au Baccalauréat A, le commentaire de texte est souvent noté autant sur la \emph{pertinence des idées} que sur la langue : un plan clair avec des phrases simples et correctes vaut toujours mieux qu'un style compliqué plein de fautes. Le correcteur cherche d'abord : avez-vous compris le texte ? Citez-vous le texte ? Donnez-vous un avis ? La langue vient ensuite. Entraînez-vous donc à rédiger des commentaires \emph{courts mais structurés} plutôt que longs et confus.
\end{savaistu}

\begin{exoresolu}[Transformer un commentaire faible en commentaire fort]
Énoncé. Reprenez la phrase du commentaire débil : \esp{Las ONG ayudan.} Transformez-la en un véritable paragraphe d'analyse (3 phrases) avec un connecteur, une citation et une explication.
\tcblower
Solution détaillée. On applique la règle « connecteur + idée + citation + explication » :

\esp{Además, el autor subraya el papel de las ONG que « trabajan sobre el terreno ». Estas palabras muestran que las ONG conocen los problemas reales de la población, por ejemplo en los barrios de Yaoundé o de Douala. Por eso, su ayuda es muy concreta.}

Analyse : le connecteur \esp{Además} relie le paragraphe au précédent ; la formule \esp{el autor subraya} introduit la citation ; la citation \esp{« trabajan sobre el terreno »} est courte ; la phrase \esp{Estas palabras muestran que...} l'explique avec un exemple camerounais ; la dernière phrase conclut le paragraphe. Trois phrases suffisent pour obtenir tous les points du critère « analyse ».
\end{exoresolu}

\begin{exercice}[À vous de jouer]
Voici trois phrases extraites d'un commentaire faible sur \esp{El mundo laboral hoy} : \esp{Los jóvenes buscan trabajo. Hay acuerdos entre España y Camerún. El texto es bueno.}
\begin{enumerate}
\item Ajoutez un connecteur adapté au début de chaque phrase.
\item Insérez dans la première phrase une citation du texte support, introduite par \esp{según el texto}.
\item Remplacez la dernière phrase par une véritable opinion personnelle nuancée, introduite par \esp{En mi opinión}.
\end{enumerate}
\end{exercice}

\begin{corrige}[Exercice — Commentaire modèle guidé]
Proposition de correction (d'autres réponses sont possibles) :
\begin{enumerate}
\item \textbf{Étape 1 (ajout des connecteurs) :} \esp{En primer lugar, los jóvenes buscan trabajo. Además, hay acuerdos entre España y Camerún. En conclusión, el texto es bueno.}
\item \textbf{Étape 2 (insertion de la citation) :} \esp{En primer lugar, según el texto, los jóvenes africanos « buscan un empleo digno ».}
\item \textbf{Étape 3 (opinion personnelle nuancée — phrase finale du commentaire) :} \esp{En mi opinión, este texto es interesante porque muestra que la ayuda es útil, pero que el empleo duradero es más importante.}
\end{enumerate}
\textbf{Paragraphe final reconstitué :}
\esp{En primer lugar, según el texto, los jóvenes africanos « buscan un empleo digno ». Además, hay acuerdos entre España y Camerún. En mi opinión, este texto es interesante porque muestra que la ayuda es útil, pero que el empleo duradero es más importante.}

Barème indicatif : 1 point par connecteur correct, 2 points pour la citation bien introduite, 2 points pour une opinion nuancée avec \esp{pero} ou \esp{sin embargo}.
\end{corrige}

\begin{experience}[Atelier collectif : la chasse aux critères]
En groupes de quatre, relisez à voix haute le \esp{Comentario modelo}. Chaque groupe reçoit une mission : le groupe 1 surligne les \emph{connecteurs}, le groupe 2 encadre les \emph{citations}, le groupe 3 souligne les \emph{verbes de parole} (\esp{dice, explica, subraya, muestra}), le groupe 4 repère l'\emph{opinion personnelle}. Mise en commun au tableau : la classe reconstitue la « radiographie » complète du commentaire, puis chaque élève rédige une phrase d'analyse supplémentaire sur le modèle « connecteur + idée + citation + explication ».
\end{experience}

\coursec{Production guidée}

Dans la section précédente, nous avons construit ensemble un commentaire modèle guidé du texte \esp{El mundo laboral hoy} : introduction, plan, analyse des idées directrices. Il est temps maintenant de passer du modèle du professeur à \textbf{votre propre production}. Cette dernière étape suit une progression en trois marches : le \cle{résumé} (reformuler), l'\cle{opinion} (juger), le \cle{débat} (argumenter à l'oral). Chaque tâche réutilise le lexique de l'aide au développement et les outils grammaticaux du module.

\begin{popfigure}
\begin{tikzpicture}[
  marche/.style={rectangle, draw=popblue, fill=popblueL, rounded corners=4pt, minimum width=4.6cm, minimum height=1.1cm, align=center, font=\small},
  fleche/.style={-{Stealth[length=3mm]}, thick, poporange}
]
\node[marche, align=center] (m1) at (0,0){\textbf{1. Resumen}\\ reformuler en 5--6 phrases};
\node[marche, fill=poporangeL, draw=poporange, align=center] (m2) at (3.4,1.3){\textbf{2. Opinión}\\ conclusion + subjonctif};
\node[marche, fill=popgreenL, draw=popgreen, align=center] (m3) at (6.8,2.6){\textbf{3. Debate}\\ argumenter à l'oral};
\draw[fleche] (m1.east) -- (m2.west);
\draw[fleche] (m2.east) -- (m3.west);
\node[font=\small\itshape, popmuted, align=center] at (3.4,-1.0) {Difficulté croissante : de la reformulation à la prise de position orale};
\end{tikzpicture}
\legende{L'escalier de la production : résumé, opinion, débat}
\end{popfigure}

\courssub{Tâche 1 (écrite) : rédiger un résumé du texte}

\textbf{Consigne.} \esp{Resume el texto « El mundo laboral hoy » en 5 o 6 frases, utilizando al menos cinco palabras del léxico estudiado (subvención, ayuda al desarrollo, acuerdo bilateral, inversión, cooperación, ONG).}

\begin{methode}[Réussir un résumé en espagnol]
\begin{enumerate}
\item \textbf{Relire} le texte en soulignant les idées principales (une par paragraphe).
\item \textbf{Supprimer} les exemples, les chiffres, les répétitions et les détails secondaires.
\item \textbf{Reformuler} avec ses propres mots : interdiction de recopier des phrases entières du texte.
\item \textbf{Relier} les idées avec des connecteurs : \esp{primero}, \esp{además}, \esp{sin embargo}, \esp{por eso}, \esp{en conclusión}.
\item \textbf{Vérifier} : le résumé doit être compréhensible par quelqu'un qui n'a pas lu le texte, et contenir au moins cinq mots du lexique.
\end{enumerate}
\end{methode}

\begin{attention}[Pièges du résumé pour les francophones]
\begin{itemize}
\item Ne traduisez pas « résumer » par \esp{resumir} au sens de « conclure » : \esp{resumir} signifie bien « résumer », mais le nom est \esp{el resumen} (et non \esp{el resúmen} : pas d'accent !).
\item Attention au faux ami \esp{actualmente} : il signifie « actuellement » (en ce moment), \textbf{pas} « en fait » (pour « en fait », on utilise \esp{de hecho} ou \esp{en realidad}).
\item Évitez le calque « selon le texte » $\rightarrow$ \esp{según el texto} (correct), mais jamais \esp{acuerdo al texto}.
\end{itemize}
\end{attention}

\begin{exoresolu}[Résumé modèle]
Énoncé : rédigez le résumé du texte étudié en 5--6 phrases avec cinq mots du lexique.
\tcblower
Solution détaillée (modèle de résumé) :

\esp{El texto presenta la situación del mundo laboral actual, marcada por el desempleo juvenil y la precariedad. Primero, explica que la cooperación internacional y la ayuda al desarrollo intentan mejorar las condiciones de vida en África. Además, menciona el papel de las ONG y de las subvenciones públicas en la creación de empleo. Sin embargo, subraya que los acuerdos bilaterales no siempre benefician a los países africanos. Por eso, el autor defiende una inversión más justa y duradera. En conclusión, el texto invita a reflexionar sobre un desarrollo económico más equilibrado.}

Vérification : 6 phrases ; mots du lexique employés : \esp{cooperación}, \esp{ayuda al desarrollo}, \esp{ONG}, \esp{subvenciones}, \esp{acuerdos bilaterales}, \esp{inversión} (6 mots $\geq$ 5) ; connecteurs : \esp{primero}, \esp{además}, \esp{sin embargo}, \esp{por eso}, \esp{en conclusión}.
\end{exoresolu}

\courssub{Tâche 2 (écrite) : rédiger la conclusion d'un commentaire}

\textbf{Consigne.} \esp{Escribe la conclusión de un comentario del texto, dando tu opinión sobre la ayuda al desarrollo (8--10 líneas).}

La conclusion d'un commentaire de texte, en Terminale A, comporte trois éléments : le \textbf{bilan} (rappel de l'idée principale), l'\textbf{ouverture} (élargissement du sujet) et l'\textbf{opinion personnelle}. C'est l'opinion qui exige le plus de soin grammatical : après les verbes d'opinion à la forme négative et les expressions de doute, l'espagnol impose le \cle{subjonctif}.

\begin{propriete}[Exprimer et nuancer son opinion]
\begin{itemize}
\item \textbf{Opinion affirmative} (indicatif) : \esp{Creo que la cooperación es necesaria.} « Je pense que la coopération est nécessaire. »
\item \textbf{Opinion négative ou doute} (subjonctif) : \esp{No creo que la ayuda sea suficiente.} « Je ne pense pas que l'aide soit suffisante. »
\item \textbf{Nuance} : \esp{Es verdad que..., pero...} « Il est vrai que..., mais... » ; \esp{Aunque la ayuda es necesaria, no es suficiente.} « Bien que l'aide soit nécessaire, elle n'est pas suffisante. » (avec \esp{aunque} + indicatif quand le fait est présenté comme réel).
\end{itemize}
\end{propriete}

\begin{exemplebox}[Formules d'opinion traduites]
\begin{itemize}
\item \esp{A mi juicio, la cooperación internacional es indispensable.} « À mon avis, la coopération internationale est indispensable. »
\item \esp{En mi opinión, los acuerdos multilaterales deben respetar los intereses africanos.} « Selon moi, les accords multilatéraux doivent respecter les intérêts africains. »
\item \esp{Me parece que las ONG hacen un trabajo admirable.} « Il me semble que les ONG font un travail admirable. »
\item \esp{Dudo que las subvenciones resuelvan todos los problemas.} « Je doute que les subventions résolvent tous les problèmes. » (subjonctif après \esp{dudar que})
\end{itemize}
\end{exemplebox}

\begin{exemplebox}[Conclusion modèle]
\esp{En conclusión, el texto muestra que el mundo laboral actual es un desafío para los jóvenes africanos. A mi juicio, la ayuda al desarrollo es útil, pero no creo que sea suficiente sin un comercio más justo. Es verdad que las subvenciones crean empleos, pero la inversión extranjera debe beneficiar también a la población local. Este tema nos invita a pensar en el futuro de la cooperación entre África y Europa.}
\end{exemplebox}

\courssub{Tâche 3 (orale) : mini-débat en binômes}

\begin{experience}[¿África necesita ayuda o comercio justo?]
Par binômes, l'élève A défend la thèse de l'aide au développement, l'élève B défend celle du commerce équitable. Chaque intervention (1 minute) doit obligatoirement contenir les quatre mots-clés : \esp{subvención}, \esp{inversión}, \esp{ONG}, \esp{empleo}. Déroulement :
\begin{enumerate}
\item Préparation individuelle (5 min) : noter deux arguments et un exemple (Cameroun, Guinée équatoriale, Espagne...).
\item Débat (4 min) : alternance des prises de parole, réaction obligatoire à l'argument du partenaire.
\item Synthèse commune (1 min) : les deux élèves formulent ensemble une position nuancée.
\end{enumerate}
\end{experience}

\begin{attention}[Dire son accord : la bonne construction]
En français on dit « je suis d'accord » avec le verbe « être » + « d'accord ». En espagnol, la locution est \esp{estar de acuerdo} : la préposition \esp{de} est \textbf{obligatoire}.
\begin{itemize}
\item Correct : \esp{Estoy de acuerdo contigo.} « Je suis d'accord avec toi. »
\item Correct : \esp{No estoy de acuerdo con esa idea.} « Je ne suis pas d'accord avec cette idée. »
\item Faux : \esp{Estoy acuerdo.} (calque du français, à bannir absolument)
\item Autre formule utile : \esp{Tienes razón.} « Tu as raison. » / \esp{No tienes razón.} « Tu n'as pas raison. »
\end{itemize}
\end{attention}

\begin{center}
\begin{tabularx}{\linewidth}{|X|X|X|}
\hline
\rowcolor{popblueL} Français & Español & Remarque \\
\hline
Je suis d'accord & Estoy de acuerdo & \esp{de} obligatoire \\
\hline
Je ne suis pas d'accord & No estoy de acuerdo & négation devant le verbe \\
\hline
À mon avis & A mi juicio / En mi opinión & pour ouvrir une prise de parole \\
\hline
Il est vrai que..., mais... & Es verdad que..., pero... & nuance + indicatif \\
\hline
Bien que l'aide soit nécessaire & Aunque la ayuda es necesaria & fait réel $\rightarrow$ indicatif \\
\hline
Je ne pense pas que ce soit juste & No creo que sea justo & subjonctif obligatoire \\
\hline
D'un côté / de l'autre & Por un lado / por otro lado & connecteurs de débat \\
\hline
\end{tabularx}
\end{center}

\courssub{Texte d'entraînement à l'argumentation}

Pour alimenter le débat, lisez ce court texte d'opinion original ; il servira de réservoir d'arguments pour les deux camps.

\begin{textebox}[¿Ayuda o comercio? Un debate necesario]
En muchos países africanos, el debate sobre el desarrollo económico está más vivo que nunca. Por un lado, los defensores de la ayuda internacional recuerdan que las subvenciones financian escuelas, hospitales y carreteras. Sin la cooperación de los países ricos, dicen, muchas ONG no podrían actuar en las zonas rurales, y miles de jóvenes quedarían sin formación profesional ni empleo.

Por otro lado, cada vez más voces afirman que África no necesita caridad, sino justicia comercial. Según ellos, los acuerdos bilaterales y multilaterales actuales benefician sobre todo a las empresas extranjeras: el cacao camerunés o el petróleo de Guinea Ecuatorial se venden baratos, mientras los productos transformados vuelven caros. Lo que hace falta, repiten, es una inversión que cree fábricas y empleos en el continente, y no solamente donaciones.

Entre las dos posiciones, algunos economistas proponen un camino intermedio: mantener la ayuda al desarrollo para las urgencias sociales, pero transformar poco a poco las relaciones económicas hacia un comercio más equilibrado. Aunque la ayuda es necesaria hoy, no es suficiente para construir el futuro. El verdadero objetivo, concluyen, es que los jóvenes africanos encuentren un empleo digno en su propio país, sin tener que emigrar.
\end{textebox}

\textbf{Glossaire.} \esp{vivo} : vivant, actuel ; \esp{carreteras} : routes ; \esp{formación profesional} : formation professionnelle ; \esp{caridad} : charité ; \esp{baratos / caros} : bon marché / chers ; \esp{donaciones} : dons ; \esp{equilibrado} : équilibré ; \esp{digno} : digne ; \esp{emigrar} : émigrer.

\textbf{Questions de préparation au débat.}
\begin{enumerate}
\item \esp{¿Qué argumentos presentan los defensores de la ayuda internacional?}
\item \esp{¿Por qué critican algunos los acuerdos bilaterales actuales?}
\item \esp{¿Cuál es el « camino intermedio » propuesto por los economistas?}
\item \esp{¿Y tú, estás de acuerdo con esta posición intermedia? Justifica tu respuesta.}
\end{enumerate}

\courssub{Grille d'auto-évaluation}

Après chaque production (écrite ou orale), évaluez-vous honnêtement à l'aide de cette grille. Cochez ce qui est réellement présent dans votre texte ou votre intervention.

\begin{center}
\begin{tabularx}{\linewidth}{|X|X|X|}
\hline
\rowcolor{popgreenL} Critère & Sí & Todavía no \\
\hline
He utilizado 5 palabras del léxico (subvención, inversión, ONG, acuerdo, cooperación, empleo) & & \\
\hline
He utilizado al menos 2 conectores (sin embargo, además, por eso, en conclusión...) & & \\
\hline
He expresado 1 opinión con el subjuntivo (no creo que sea..., dudo que...) & & \\
\hline
He dicho \esp{estoy de acuerdo} con la preposición \esp{de} & & \\
\hline
Mi resumen no copia ninguna frase del texto & & \\
\hline
\end{tabularx}
\end{center}

\begin{aretenir}[Les trois clés de la production]
\begin{itemize}
\item \textbf{Résumé} : idées principales + mots personnels + connecteurs.
\item \textbf{Opinion} : \esp{a mi juicio} + indicatif ; \esp{no creo que} + subjonctif.
\item \textbf{Débat} : \esp{estar de acuerdo}, arguments appuyés sur le lexique économique.
\end{itemize}
\end{aretenir}

\courssub{Correction collective d'une production d'élève}

Voici la conclusion rédigée par un élève de Terminale A, reproduite avec ses erreurs. Corrigeons-la ensemble au tableau, erreur par erreur.

\begin{exoresolu}[Texte d'élève à corriger]
Énoncé : relevez et corrigez les cinq erreurs de cette production.

\esp{En conclusión, yo estoy acuerdo con el autor: la ayuda al desarrollo es muy importante. No creo que las subvenciones son suficientes para crear empleo. Aunque las ONG trabajan mucho, el desempleo juvenil sigue siendo un problema grave. Según mi opinión, los acuerdos bilaterales deben cambiar. Actualmente, la inversión extranjera no beneficia a los jóvenes africanos.}
\tcblower
Solution détaillée :
\begin{enumerate}
\item \esp{yo estoy acuerdo} $\rightarrow$ \esp{estoy de acuerdo} : préposition \esp{de} obligatoire ; le pronom \esp{yo} est inutile ici (la désinence verbale suffit).
\item \esp{no creo que las subvenciones son} $\rightarrow$ \esp{no creo que las subvenciones sean} : après \esp{no creer que}, le subjonctif est obligatoire.
\item \esp{Según mi opinión} $\rightarrow$ \esp{En mi opinión} ou \esp{A mi juicio} : \esp{según} ne se combine pas avec \esp{mi opinión} en espagnol (calque du français « selon mon opinion »).
\item \esp{Actualmente} : l'élève a écrit \esp{Actualmente} (actuellement). Si le sens est temporel, c'est correct. Mais dans le contexte d'une correction/précision (« En fait, ... »), c'est un faux ami ; il faut \esp{De hecho} ou \esp{En realidad}. La version corrigée utilise \esp{De hecho}.
\item Style : répétition de \esp{problema grave} évitable ; on peut varier avec \esp{un desafío} et ajouter un connecteur de conclusion : \esp{Por eso, los acuerdos bilaterales deben cambiar.}
\end{enumerate}
Version corrigée :

\esp{En conclusión, estoy de acuerdo con el autor: la ayuda al desarrollo es muy importante. No creo que las subvenciones sean suficientes para crear empleo. Aunque las ONG trabajan mucho, el desempleo juvenil sigue siendo un desafío grave. En mi opinión, los acuerdos bilaterales deben cambiar. De hecho, la inversión extranjera no beneficia bastante a los jóvenes africanos.}
\end{exoresolu}

\begin{exercice}[Pour aller plus loin]
Rédigez à la maison votre propre conclusion de commentaire (8--10 lignes) sur la question : \esp{¿Debe el Camerún firmar más acuerdos bilaterales con España?} Exigences : cinq mots du lexique, deux connecteurs, une opinion au subjonctif, une nuance avec \esp{aunque}. Vous l'auto-évaluerez avec la grille avant de la rendre.
\end{exercice}

\newpage

\coursec{Activité d'intégration}

\coursec{Leçon E16 : Lectura y comentario : « El mundo laboral hoy » ; ayudas, subvenciones, acuerdos bilaterales y multilaterales}

\courssub{Objectifs de l'activité}

Cette activité d'intégration clôt la leçon E16 en faisant travailler ensemble tous les savoir-faire étudiés. À l'issue de la séance, l'élève sera capable de :

\begin{itemize}
\item \cle{comprendre} un document authentique (discours, article de presse) portant sur la coopération économique entre l'Espagne et le Cameroun ;
\item \cle{réemployer} à l'oral et à l'écrit le lexique de l'aide au développement : \esp{subvención}, \esp{ayuda al desarrollo}, \esp{acuerdo bilateral/multilateral}, \esp{inversión}, \esp{cooperación}, \esp{ONG} ;
\item \cle{argumenter} une position en espagnol à l'aide des connecteurs logiques et des structures d'opinion étudiées (\esp{en mi opinión}, \esp{estoy convencido de que} + indicatif, \esp{es necesario que} + subjonctif) ;
\item \cle{transférer} la méthode du commentaire de texte (repérage des idées directrices, synthèse) dans une prise de parole structurée ;
\item \cle{rédiger} un compte rendu de débat en huit lignes, en passant correctement du discours direct au discours indirect.
\end{itemize}

\courssub{Situation}

La ville de Yaoundé accueille cette semaine le \emph{Foro Económico Hispano-Camerunés}, organisé au Palais des Congrès. Des délégations d'entreprises espagnoles, des représentants du gouvernement camerounais, des diplomates et des organisations non gouvernementales se réunissent pour discuter d'une question qui concerne directement les jeunes : la création d'emplois. Votre classe de Terminale A a été invitée à simuler la séance plénière de clôture. Le journal \emph{La Tribuna del Foro} a publié ce matin l'éditorial suivant, qui servira de document déclencheur au débat.

\begin{textebox}[La Tribuna del Foro — editorial del día]
¿Empleo joven o dependencia eterna?

Yaoundé recibe hoy a la delegación española con una pregunta en los labios de todos: ¿cómo puede la cooperación crear empleo para los jóvenes cameruneses? Las cifras del desempleo juvenil preocupan, pero también existen razones para la esperanza. Los acuerdos bilaterales firmados en los últimos años han financiado carreteras, escuelas técnicas y proyectos agrícolas en la región del Centro. Sin embargo, muchos jóvenes se preguntan si las subvenciones llegan realmente hasta los barrios de Douala o de Bamenda.

Los empresarios españoles hablan de inversión y de formación profesional; las ONG recuerdan que la ayuda al desarrollo debe respetar las prioridades locales; los jóvenes emprendedores piden créditos y menos burocracia. El debate está abierto. La cooperación no es caridad: es, o debería ser, una alianza entre socios que se respetan. El futuro de miles de jóvenes cameruneses depende de las decisiones que se tomen hoy en este foro.
\end{textebox}

\textbf{Glossaire du document :} \esp{en los labios de todos} : sur toutes les lèvres ; \esp{las cifras} : les chiffres ; \esp{carreteras} : routes ; \esp{créditos} : prêts ; \esp{burocracia} : bureaucratie ; \esp{caridad} : charité ; \esp{alianza} : alliance ; \esp{socios} : partenaires.

\courssub{Consignes}

Travail en groupes de quatre. Chaque groupe forme une délégation complète du forum. Suivez les étapes dans l'ordre.

\begin{enumerate}
\item \textbf{Lecture active (10 min).} Relisez l'éditorial en silence. Soulignez les six mots du lexique de la leçon qui y apparaissent et entourez la thèse principale du journaliste. Vérifiez en groupe que tout le monde a compris le glossaire.

\item \textbf{Répartition des rôles (5 min).} Dans chaque groupe, attribuez les quatre rôles : \esp{el representante del gobierno camerunés}, \esp{el diplomático español}, \esp{el responsable de una ONG}, \esp{el joven emprendedor}. Chacun défend un point de vue différent sur la question du forum : \esp{¿Cómo puede la cooperación crear empleo para los jóvenes cameruneses?}

\item \textbf{Préparation du discours (15 min).} Rédigez individuellement un discours de \textbf{cinq phrases minimum}. Contraintes obligatoires : employer \textbf{au moins quatre mots du lexique} de la leçon (\esp{subvención}, \esp{acuerdo bilateral}, \esp{inversión}, \esp{cooperación}, \esp{ONG}, \esp{ayuda al desarrollo}) ; utiliser \textbf{au moins une structure d'opinion} et \textbf{un verbe au subjonctif} après \esp{es necesario que} ou \esp{es importante que}.

\item \textbf{Le forum (20 min).} Chaque délégué présente son discours devant la classe, debout, sans lire intégralement ses notes. Les autres groupes écoutent et prennent des notes : ils devront citer au moins un argument entendu dans leur compte rendu.

\item \textbf{Le vote (5 min).} Après les discours, la classe vote à main levée pour désigner la proposition la plus convaincante. Justifiez votre vote en une phrase en espagnol : \esp{Voto por... porque...}.

\item \textbf{Trace écrite (15 min).} Rédigez individuellement un \textbf{compte rendu du débat en huit lignes} : présentez le contexte, rapportez au discours indirect au moins deux arguments entendus, et concluez par votre opinion personnelle argumentée.
\end{enumerate}

\begin{attention}[Pièges à éviter pendant l'activité]
Attention au faux ami \esp{eventualmente} (qui signifie « éventuellement, peut-être ») : ne l'utilisez pas pour dire « actuellement » — le mot correct est \esp{actualmente}. De même, \esp{asistir al foro} signifie « assister au forum » (être présent), et non « aider » : pour aider, utilisez \esp{ayudar a}. Enfin, n'oubliez pas le subjonctif après \esp{es necesario que} : \esp{Es necesario que los jóvenes \textbf{reciban} formación}, et non \esp{reciben}.
\end{attention}

\courssub{Ressources à mobiliser}

\begin{aretenir}[Boîte à outils linguistique]
\begin{itemize}
\item \textbf{Lexique de base :} \esp{la subvención} (la subvention), \esp{la ayuda al desarrollo} (l'aide au développement), \esp{el acuerdo bilateral / multilateral} (l'accord bilatéral / multilatéral), \esp{la inversión} (l'investissement), \esp{la cooperación} (la coopération), \esp{la ONG} (l'ONG).
\item \textbf{Lexique étendu du texte :} \esp{el desempleo juvenil} (le chômage des jeunes), \esp{la formación profesional} (la formation professionnelle), \esp{el emprendedor / la emprendedora} (l'entrepreneur), \esp{la burocracia} (la bureaucratie), \esp{el crédito} (le prêt / le crédit).
\item \textbf{Structures d'opinion :} \esp{En mi opinión...}, \esp{Estoy convencido/a de que} + indicatif, \esp{Me parece que} + indicatif, \esp{Es necesario / importante / fundamental que} + \textbf{subjonctif}.
\item \textbf{Connecteurs logiques :} \esp{en primer lugar} (d'abord), \esp{además} (de plus), \esp{sin embargo} (cependant), \esp{por lo tanto} (par conséquent), \esp{en conclusión} (en conclusion).
\item \textbf{Discours indirect (concordance des temps) :}
\begin{center}
\begin{tabular}{|l|l|l|}\hline
\rowcolor{popblueL} Discours direct & Verbe introducteur (passé) & Discours indirect \\ \hline
\esp{« La cooperación \textbf{es} una alianza »} & \esp{dijo / afirmó / explicó} & \esp{que la cooperación \textbf{era} una alianza} \\ \hline
\esp{« Las subvenciones \textbf{han financiado} »} & \esp{explicó / informó} & \esp{que las subvenciones \textbf{habían financiado}} \\ \hline
\esp{« La inversión \textbf{creará} empleo »} & \esp{afirmó / prometió} & \esp{que la inversión \textbf{crearía} empleo} \\ \hline
\esp{« Los jóvenes \textbf{necesitan} créditos »} & \esp{pidió / denunció} & \esp{que los jóvenes \textbf{necesitaban} créditos} \\ \hline
\end{tabular}
\end{center}
\end{itemize}
\end{aretenir}

\begin{methode}[Rédiger un compte rendu de débat / synthèse]
\begin{enumerate}
\item \textbf{Contextualiser} en une phrase : lieu, date, thème, participants (\esp{Ayer nuestra clase participó en el Foro... sobre el empleo de los jóvenes.}).
\item \textbf{Reporter les idées directrices} au discours indirect : identifier l'argument de chaque intervenant et appliquer la concordance des temps (présent $\rightarrow$ imparfait ; passé composé $\rightarrow$ plus-que-parfait ; futur $\rightarrow$ conditionnel).
\item \textbf{Articuler} les arguments avec des connecteurs (\esp{sin embargo, además, por lo tanto}) pour montrer l'opposition ou la complémentarité des points de vue.
\item \textbf{Conclure} par votre opinion personnelle en utilisant une structure d'opinion (\esp{En mi opinión, Me parece que, Es necesario que} + subjonctif).
\item \textbf{Respecter la contrainte de longueur} (huit lignes environ) et la correction linguistique (accords, accents, ponctuation).
\end{enumerate}
\end{methode}

\courssub{Proposition de production et corrigé commenté}

\begin{exoresolu}[Discours modèle : la jeune entrepreneure]
Énoncé : rédigez le discours de \esp{la joven emprendedora} en cinq phrases, avec quatre mots du lexique et un subjonctif.
\tcblower
\textbf{Production modèle :}

\esp{Buenas tardes, señoras y señores delegados. Me llamo Nadège y soy una joven emprendedora de Douala. En mi opinión, la cooperación entre España y Camerún solo será útil si crea empleo real para los jóvenes. Los acuerdos bilaterales y las subvenciones son importantes, pero sin embargo el dinero no llega a los pequeños empresarios de los barrios. Es necesario que una parte de cada inversión española financie directamente los proyectos de los jóvenes, por ejemplo en la agricultura o en la tecnología. Por lo tanto, pido menos burocracia, más créditos y una verdadera alianza entre nuestros dos países. Muchas gracias.}

\textbf{Commentaire des choix de langue :}
\begin{itemize}
\item Le discours s'ouvre par une formule de politesse adaptée au contexte formel (\esp{señoras y señores delegados}) et une présentation brève : c'est la première étape de toute prise de parole publique.
\item Le lexique imposé est bien intégré en contexte : \esp{cooperación}, \esp{acuerdos bilaterales}, \esp{subvenciones}, \esp{inversión} — quatre mots-clefs, la contrainte est respectée.
\item La structure d'opinion \esp{En mi opinión} introduit la thèse, suivie d'une proposition concessive avec \esp{sin embargo}, ce qui montre une argumentation nuancée et non un simple catalogue.
\item Le subjonctif est correctement employé après \esp{Es necesario que} : \esp{financie} (subjonctif présent de \esp{financiar}), car on exprime une exigence, un souhait, et non un fait constaté.
\item La conclusion (\esp{Por lo tanto, pido...}) reprend le verbe \esp{pedir} à la première personne, marque forte de l'engagement du locuteur, et le mot \esp{alianza} fait écho à l'éditorial, ce qui boucle le discours.
\end{itemize}
\end{exoresolu}

\begin{exoresolu}[Compte rendu modèle du débat (trace écrite)]
Énoncé : rédigez le compte rendu du forum en huit lignes, avec deux arguments rapportés au discours indirect.
\tcblower
\textbf{Production modèle :}

\esp{Ayer nuestra clase participó en el Foro Económico Hispano-Camerunés sobre el empleo de los jóvenes. El representante del gobierno camerunés explicó que los acuerdos bilaterales ya habían financiado escuelas técnicas y carreteras. El diplomático español afirmó que la inversión de las empresas de su país crearía miles de puestos de trabajo. Sin embargo, el responsable de la ONG denunció que las subvenciones no llegaban siempre a los barrios más pobres. Finalmente, la joven emprendedora pidió créditos y menos burocracia para los jóvenes. La clase votó por la propuesta de la emprendedora porque era la más concreta. En mi opinión, la cooperación es útil si respeta las prioridades de los cameruneses. Es necesario que los jóvenes participen de verdad en estos proyectos.}

\textbf{Commentaire :} le compte rendu suit la méthode du commentaire de texte : contexte d'abord, idées directrices ensuite, conclusion personnelle à la fin. Remarquez les concordances des temps au discours indirect : \esp{explicó que... habían financiado} (plus-que-parfait pour une action antérieure), \esp{afirmó que... crearía} (conditionnel pour un futur rapporté), \esp{denunció que... no llegaban} (imparfait pour un présent rapporté). L'opinion finale utilise à nouveau le subjonctif après \esp{Es necesario que}, ce qui consolide la règle grammaticale dans un contexte d'emploi authentique.
\end{exoresolu}

\courssub{Bilan de l'activité}

Cette simulation de forum a permis de réinvestir l'ensemble des acquis de la leçon E16 dans une tâche finale de communication authentique. Sur le plan de la \cle{compréhension}, l'éditorial de \emph{La Tribuna del Foro} a servi de point de départ pour identifier une thèse et des arguments, exactement comme dans un commentaire de texte d'examen. Sur le plan du \cle{lexique}, les six mots-clefs de l'aide au développement ont été compris, prononcés et réemployés dans des discours personnels, ce qui garantit leur mémorisation durable. Sur le plan \cle{grammatical}, les structures d'opinion et le subjonctif après \esp{es necesario que} ont été mobilisés à l'oral puis à l'écrit. Enfin, la trace écrite a entraîné le passage au discours indirect, compétence essentielle pour l'épreuve de traduction et de synthèse du baccalauréat. L'élève est désormais capable de lire un texte économique, d'en dégager les idées directrices et de défendre en espagnol une position argumentée sur les enjeux du développement en Afrique.

\newpage

\coursec{Méthodes et exercices résolus}

\coursec{Lección E16 : Lectura y comentario : « El mundo laboral hoy » ; ayudas, subvenciones, acuerdos bilaterales y multilaterales}

\courssub{Méthode 1 : Analyser et commenter un texte économique}

\begin{methode}[Commentaire de texte économique : démarche en 5 étapes]
\textbf{Étape 1 : Lecture globale et identification} — Lire le texte deux fois. Noter : \textbf{source} (journal, rapport, discours), \textbf{date}, \textbf{auteur}, \textbf{type de document} (article, rapport d'ONG, communiqué officiel). Repérer le \textbf{thème central} (ici : monde du travail et aide au développement) et la \textbf{thèse} défendue.

\textbf{Étape 2 : Découpage en idées directrices} — Diviser le texte en \textbf{paragraphes logiques} (introduction, développement, conclusion). Pour chaque paragraphe, formuler en une phrase l'\textbf{idée-force} (ex. : « Paragraphe 2 : Le chômage des jeunes en Afrique subsaharienne dépasse 20 \% »).

\textbf{Étape 3 : Analyse lexicale et stylistique} — Relever le \textbf{lexique spécialisé} (\cle{subvención}, \cle{acuerdo bilateral}, \cle{inversión extranjera}, \cle{cooperación técnica}) et les \textbf{connecteurs logiques} (\esp{por consiguiente}, \esp{no obstante}, \esp{en cuanto a}). Noter le \textbf{registre} (formel, objectif, argumentatif) et les \textbf{figures de style} (chiffres, exemples concrets, citations d'experts).

\textbf{Étape 4 : Évaluation critique} — Vérifier la \textbf{cohérence} des arguments, la \textbf{fiabilité} des données (sources citées ?), les \textbf{angles morts} (quels acteurs absents ?). Confronter au \textbf{contexte camerounais} (ex. : programmes PAD, accords UE-Cameroun, rôle de la diaspora).

\textbf{Étape 5 : Rédaction structurée} — \textbf{Introduction} : accroche + présentation du texte + annonce du plan. \textbf{Développement} : 2 ou 3 parties (selon idées directrices), chaque partie = idée + citation courte + analyse + ouverture contextuelle. \textbf{Conclusion} : synthèse + évaluation portée + perspective personnelle.
\end{methode}

\begin{textebox}[El mundo laboral hoy — Texte support original]
El mercado laboral africano se encuentra en una encrucijada histórica. Según la OIT, más del 60 \% de la población activa en África subsahariana trabaja en la economía informal, sin protección social ni contratos estables. En Camerún, por ejemplo, miles de jóvenes titulados universitarios engrosan las filas del desempleo o aceptan empleos precarios en el sector informal —vendedores ambulantes, conductores de \emph{bendskin}, pequeños comerciantes en los mercados de Yaoundé o Douala.

Ante este panorama, la cooperación internacional juega un papel decisivo. Los \textbf{acuerdos bilaterales} entre Camerún y países como España, Francia o China financian proyectos de formación profesional y de apoyo al emprendimiento juvenil. Las \textbf{subvenciones} de la Unión Europea, canalizadas a través de programas como el \emph{Programa Indicativo Plurianual}, permiten la creación de centros de formación técnica en regiones rurales. Por su parte, las \textbf{ONG} internacionales y locales complementan la acción estatal: forman a mujeres en \textbf{agroindustria}, apoyan cooperativas de cacao certificado y facilitan el acceso al microcrédito.

Sin embargo, la \textbf{ayuda al desarrollo} no está exenta de críticas. Algunos economistas señalan que la \textbf{inversión extranjera directa} privilegia a menudo la extracción de materias primas (petróleo, madera, cacao) sin generar empleo local cualificado. Los \textbf{acuerdos multilaterales} con el FMI o el Banco Mundial imponen a veces condicionalidades que reducen el gasto público en educación y salud. La \textbf{cooperación Sur-Sur}, impulsada por la Unión Africana y socios como Brasil o la India, ofrece una alternativa basada en el intercambio de experiencias y la transferencia de tecnología adaptada.

El reto principal sigue siendo la \textbf{adecuación entre formación y mercado de trabajo}. Las universidades deben dialogar con las empresas para diseñar currículos que respondan a las necesidades reales: energías renovables, TIC, transformación agroalimentaria. Sólo así la ayuda externa dejará de ser una muleta para convertirse en palanca de un desarrollo endógeno y sostenible.
\end{textebox}

\begin{exoresolu}[Compréhension et commentaire guidé du texte « El mundo laboral hoy »]
\textbf{Glossaire des termes clés :}
\begin{center}
\begin{tabularx}{\linewidth}{|X|X|X|}
\hline
\rowcolor{popblueL} \textbf{Español} & \textbf{Français} & \textbf{Remarque} \\
\hline
encrucijada histórica & carrefour historique & métaphore : moment décisif \\
economía informal & économie informelle & sans contrat, sans protection sociale \\
engrosar las filas & grossir les rangs & locution verbale \\
titulados universitarios & diplômés universitaires & \cle{titular} = titulaire \\
Sin embargo & cependant & locution adverbiale \\
canalizadas a través de & canalisées via & participe passé comme adjectif \\
agroindustria & agro-industrie & mot composé, accent sur \emph{ú} \\
microcrédito & microcrédit & prêt de faible montant \\
condicionalidades & conditionalités & conditions imposées par donateurs \\
adecuación & adéquation, correspondance & nom féminin \\
muleta & béquille & image : aide temporaire \\
palanca & levier & image : force motrice \\
desarrollo endógeno & développement endogène & venu de l'intérieur \\
\hline
\end{tabularx}
\end{center}

\textbf{Questions de compréhension (type Bac) :}
\begin{enumerate}
\item Quel pourcentage de la population active travaille dans l'économie informelle en Afrique subsaharienne selon l'OIT ?
\item Citez trois exemples d'emplois précaires mentionnés pour les jeunes camerounais.
\item Quels sont les deux types d'accords internationaux évoqués ? Donnez un exemple pour chacun.
\item Quel rôle jouent les ONG selon le texte ?
\item Quelles sont les deux critiques principales adressées à l'aide au développement ?
\item Expliquez l'expression « \esp{la ayuda externa dejará de ser una muleta para convertirse en palanca} ».
\item Quelle solution propose le texte pour améliorer l'adéquation formation-emploi ?
\end{enumerate}

\tcblower

\textbf{Solution détaillée :}

\textbf{Question 1.} « \esp{Más del 60 \% de la población activa en África subsahariana trabaja en la economía informal} » → \textbf{Plus de 60 \%}.

\textbf{Question 2.} « \esp{vendedores ambulantes, conductores de \emph{bendskin}, pequeños comerciantes} » → \textbf{Vendeurs ambulants, conducteurs de moto-taxi (bendskin), petits commerçants}.

\textbf{Question 3.}
\begin{itemize}
\item \textbf{Accords bilatéraux} : exemple \esp{acuerdos bilaterales entre Camerún y países como España, Francia o China}.
\item \textbf{Accords multilatéraux} : exemple \esp{acuerdos multilaterales con el FMI o el Banco Mundial}.
\end{itemize}

\textbf{Question 4.} Les ONG « \esp{forman a mujeres en agroindustria, apoyan cooperativas de cacao certificado y facilitan el acceso al microcrédito} » → \textbf{Forment les femmes en agro-industrie, soutiennent des coopératives de cacao certifié, facilitent l'accès au microcrédit}.

\textbf{Question 5.}
\begin{enumerate}
\item L'investissement étranger direct privilégie l'extraction de matières premières sans créer d'emploi local qualifié.
\item Les accords multilatéraux (FMI, Banque mondiale) imposent des conditionalités réduisant les dépenses publiques en éducation et santé.
\end{enumerate}

\textbf{Question 6.} \esp{Muleta} = béquille (appui temporaire, passif) ; \esp{palanca} = levier (force qui démultiplie l'action). L'aide ne doit plus être une simple assistance mais un \textbf{levier} permettant un développement autonome (\esp{endógeno}) et durable.

\textbf{Question 7.} « \esp{Las universidades deben dialogar con las empresas para diseñar currículos que respondan a las necesidades reales} » → \textbf{Dialogue universités-entreprises pour adapter les programmes aux besoins réels (énergies renouvelables, TIC, agroalimentaire)}.

\textbf{Commentaire synthétique (modèle) :}
Le texte « El mundo laboral hoy » dresse un tableau lucide du marché du travail africain, caractérisé par la prédominance de l'informel (60 \% selon l'OIT) et le chômage des diplômés, illustré par le cas camerounais (\emph{bendskin}, commerce de rue). Il analyse ensuite les réponses internationales : aide bilatérale (formation, entrepreneuriat), subventions européennes (centres techniques ruraux), action des ONG (femmes, coopératives, microcrédit). Une partie critique souligne les limites : investissement extractif, conditionalités néolibérales, et propose la coopération Sud-Sur comme alternative. La conclusion plaide pour l'adéquation formation-emploi via le dialogue université-entreprise.

\textbf{Points forts du commentaire :} structure claire (constat → acteurs → critiques → solution), vocabulaire précis (\cle{economía informal}, \cle{adecuación}, \cle{desarrollo endógeno}), exemples contextuels (Cameroun, cacao, bendskin), évaluation critique équilibrée.
\end{exoresolu}

\courssub{Méthode 2 : Maîtriser le lexique de l'aide au développement et du monde du travail}

\begin{methode}[Lexique thématique : stratégies d'acquisition et d'emploi]
\textbf{1. Classer par champs sémantiques} — Créer un \textbf{tableau à trois colonnes} : \textit{Acteurs} (Estado, ONG, empresa, universidad, cooperativa, diáspora), \textit{Instruments} (subvención, préstamo, microcrédito, inversión, acuerdo bilateral/multilateral, convenio), \textit{Résultats visés} (empleo, formación, emprendimiento, desarrollo sostenible, inclusión social, equidad).

\textbf{2. Travailler les familles de mots} — À partir d'un radical, dériver : \esp{cooperar} (verbe) → \esp{cooperación} (nom) → \esp{cooperante} (adj./nom) → \esp{cooperativo/a} (adj.) → \esp{cooperativismo} (mouvement). Faire de même pour \esp{desarrollar/desarrollo/desarrollista}, \esp{invertir/inversión/inversor}, \esp{subvencionar/subvención/subvencionado}.

\textbf{3. Maîtriser les collocations (associations figées)} — Apprendre par cœur : \esp{firmar un acuerdo}, \esp{conceder una subvención}, \esp{canalizar fondos}, \esp{imponer condicionalidades}, \esp{generar empleo}, \esp{reducir la brecha}, \esp{fomentar el emprendimiento}, \esp{garantizar la protección social}.

\textbf{4. Attention aux faux amis et calques} — \esp{Empresa} ≠ emploi (c'est « entreprise ») ; \esp{Éxito} ≠ sortie (c'est « succès ») ; \esp{Compromiso} ≠ compromis (c'est « engagement ») ; \esp{Subvención} = subvention (pas « soumission ») ; \esp{Inversión} = investissement (pas « inversion ») ; \esp{Actual} = actuel (pas « actuel » au sens de « réel » → \esp{real}).

\textbf{5. S'entraîner en contexte} — Rédiger 5 phrases personnelles utilisant chaque mot-clé dans une phrase complète sur le Cameroun (ex. : \esp{El Gobierno camerunés ha firmado acuerdos bilaterales con España para la formación profesional.}).
\end{methode}

\begin{exoresolu}[Exercices d'application lexicale — Niveau Bac]
\textbf{Exercice A — Complétez les phrases avec le mot approprié (forme correcte) :}
\begin{enumerate}
\item El \_\_\_\_\_\_\_\_ entre Camerún y la UE permite financiar centros de formación rural. (acuerdo bilateral / subvención / inversión)
\item Las \_\_\_\_\_\_\_\_ trabajan con cooperativas de mujeres para mejorar la producción de cacao. (ONG / empresas / universidades)
\item La \_\_\_\_\_\_\_\_ extranjera directa a veces solo busca materias primas. (ayuda / cooperación / inversión)
\item Los jóvenes necesitan \_\_\_\_\_\_\_\_ para crear sus propios negocios. (microcrédito / condicionalidad / brecha)
\item La \_\_\_\_\_\_\_\_ Sur-Sur ofrece alternativas al modelo tradicional. (cooperación / subvención / acuerdo)
\end{enumerate}

\textbf{Exercice B — Associez chaque terme à sa définition :}
\begin{center}
\begin{tabular}{|l|l|}
\hline
1. Subvención & A. Prêt de faible montant pour micro-entrepreneurs \\
2. Microcrédito & B. Aide financière non remboursable pour un projet \\
3. Acuerdo multilateral & C. Accord entre plusieurs pays/organisations \\
4. Condicionalidad & D. Exigence imposée par un bailleur de fonds \\
5. Economía informal & E. Secteur non réglementé, sans protection sociale \\
\hline
\end{tabular}
\end{center}

\textbf{Exercice C — Traduisez en espagnol (thème) :}
\begin{enumerate}
\item Le gouvernement a signé un accord bilatéral avec la France.
\item Les ONG accordent des subventions aux coopératives agricoles.
\item L'investissement étranger doit créer de l'emploi local qualifié.
\item La coopération Sud-Sur favorise le transfert de technologie.
\item Il faut réduire l'écart entre la formation et le marché du travail.
\end{enumerate}

\textbf{Exercice D — Version : Traduisez en français :}
\esp{« Las universidades deben dialogar con las empresas para diseñar currículos que respondan a las necesidades reales: energías renovables, TIC, transformación agroalimentaria. Solo así la ayuda externa dejará de ser una muleta para convertirse en palanca de un desarrollo endógeno y sostenible. »}

\tcblower

\textbf{Solution détaillée :}

\textbf{Exercice A :}
\begin{enumerate}
\item \esp{acuerdo bilateral} (contexte : entre deux pays)
\item \esp{ONG} (sujet pluriel, verbe \esp{trabajan})
\item \esp{inversión} (collocation \esp{inversión extranjera directa})
\item \esp{microcrédito} (contexte : créer entreprises)
\item \esp{cooperación} (collocation \esp{cooperación Sur-Sur})
\end{enumerate}

\textbf{Exercice B :} 1–B, 2–A, 3–C, 4–D, 5–E.

\textbf{Exercice C (Thème) :}
\begin{enumerate}
\item \esp{El Gobierno ha firmado un acuerdo bilateral con Francia.}
\item \esp{Las ONG conceden subvenciones a las cooperativas agrícolas.}
\item \esp{La inversión extranjera debe crear empleo local cualificado.}
\item \esp{La cooperación Sur-Sur favorece la transferencia de tecnología.}
\item \esp{Hay que reducir la brecha entre la formación y el mercado laboral.} (\cle{brecha} = écart, fossé ; \cle{mercado laboral} = marché du travail)
\end{enumerate}

\textbf{Exercice D (Version) :}
« Les universités doivent dialoguer avec les entreprises pour concevoir des programmes qui répondent aux besoins réels : énergies renouvelables, TIC, transformation agroalimentaire. Ce n'est qu'ainsi que l'aide extérieure cessera d'être une béquille pour devenir le levier d'un développement endogène et durable. »

\textbf{Notes grammaticales :}
\begin{itemize}
\item \esp{deben dialogar} = \cle{deber} + infinitif (obligation morale) → « doivent ».
\item \esp{diseñar} = infinitif après \cle{para} (finalité).
\item \esp{respondan} = \textbf{subjonctif présent} (3e pers. pl.) après \esp{que} + verbe de volonté/nécessité (\esp{responder a las necesidades}).
\item \esp{dejará de ser} = \cle{dejar de} + infinitif (cesser de) + futur simple.
\item \esp{convertirse en} = verbe pronominal (devenir).
\end{itemize}
\end{exoresolu}

\courssub{Méthode 3 : Synthétiser les idées directrices d'un texte argumentatif}

\begin{methode}[La synthèse en espagnol : structure et connecteurs]
\textbf{Objectif :} Produire un résumé de \textbf{120–150 mots} (consigne Bac) qui restitue la \textbf{logique interne} du texte sans y ajouter d'opinion personnelle.

\textbf{Étapes :}
\begin{enumerate}
\item \textbf{Repérer les articulations} : introduction (thèse), arguments (paragraphe 1, 2, 3...), conclusion (ouverture).
\item \textbf{Éliminer les détails} : supprimer exemples chiffrés précis, citations, anecdotes (sauf si emblématiques).
\item \textbf{Reformuler} : passer du \textbf{discours direct/indirect} au \textbf{discours rapporté} (\esp{El autor señala que...}, \esp{El texto destaca...}, \esp{Se critica...}). Utiliser l'\textbf{impersonnel} (\esp{Se afirma que}, \esp{Se denuncia}).
\item \textbf{Chaîner avec des connecteurs logiques} :
\begin{itemize}
\item \textit{Énumération} : \esp{En primer lugar}, \esp{Además}, \esp{También}, \esp{Por otra parte}.
\item \textit{Opposition} : \esp{No obstante}, \esp{Sin embargo}, \esp{En cambio}, \esp{A pesar de que}.
\item \textit{Conséquence} : \esp{Por consiguiente}, \esp{Por tanto}, \esp{Como resultado}.
\item \textit{Conclusion} : \esp{En definitiva}, \esp{En conclusión}, \esp{Finalmente}.
\end{itemize}
\item \textbf{Vérifier le nombre de mots} : compter (mots espagnols = groupes séparés par espaces). Ajuster en développant ou condensant.
\end{enumerate}

\textbf{Structure type de la synthèse :}
\begin{quote}
\esp{El texto « El mundo laboral hoy » aborda la problemática del empleo en África subsahariana. En primer lugar, denuncia la precariedad del sector informal y el desempleo juvenil, ilustrado con el caso de Camerún. Además, analiza los instrumentos de la cooperación internacional: acuerdos bilaterales, subvenciones europeas y acción de ONG. No obstante, señala límites: inversión extractiva y condicionalidades del FMI. Por tanto, propone la cooperación Sur-Sur y, sobre todo, la adecuación formación-empresa como clave para un desarrollo endógeno.}
\end{quote}
\end{methode}

\begin{textebox}[Síntesis — El mundo laboral hoy]
El texto « El mundo laboral hoy » analiza los retos del mercado laboral africano y el papel de la cooperación internacional. En primer lugar, denuncia que más del 60 \% de la población activa subsahariana trabaja en la economía informal, sin protección social, y pone como ejemplo a Camerún, donde jóvenes titulados recurren a empleos precarios como el \emph{bendskin} o el comercio ambulante.

Ante esta situación, la ayuda externa moviliza diversos instrumentos: acuerdos bilaterales para la formación profesional, subvenciones de la UE para centros técnicos rurales y acción de ONG en agroindustria y microcrédito. Sin embargo, el texto critica que la inversión extranjera directa privilegia la extracción de materias primas sin crear empleo cualificado, y que los acuerdos multilaterales imponen condicionalidades que reducen el gasto social.

En consecuencia, se propone la cooperación Sur-Sur como alternativa y, fundamentalmente, la adecuación entre la formación universitaria y las necesidades reales del mercado (energías renovables, TIC, agroalimentación) para transformar la ayuda en palanca de un desarrollo endógeno y sostenible.
\end{textebox}

\begin{exoresolu}[Synthèse guidée — Sujet type Bac : « Résumez le texte en 130 mots environ »]
\textbf{Consigne :} « \esp{Redacte una síntesis del texto « El mundo laboral hoy » en 130 palabras (± 10 \%), utilizando tus propias palabras y respetando la estructura argumentativa del original.} »

\textbf{Comptage des mots :} 132 mots (dans la fourchette 117–143).

\tcblower

\textbf{Analyse de la synthèse (pourquoi elle marche) :}

\begin{enumerate}
\item \textbf{Respect de la structure} : 3 paragraphes = 3 parties du texte original (constat → acteurs/aides → critiques/solutions).
\item \textbf{Reformulation réussie} : pas de copier-coller ; \esp{denuncia que...} (au lieu de \esp{Según la OIT...}), \esp{pone como ejemplo} (au lieu de \esp{En Camerún, por ejemplo...}), \esp{moviliza diversos instrumentos} (regroupe accords, subventions, ONG).
\item \textbf{Connecteurs logiques variés} : \esp{En primer lugar}, \esp{Ante esta situación}, \esp{Sin embargo}, \esp{En consecuencia}, \esp{Fundamentalmente}.
\item \textbf{Vocabulaire précis conservé} : \esp{economía informal}, \esp{protección social}, \esp{acuerdos bilaterales}, \esp{subvenciones}, \esp{inversión extranjera directa}, \esp{condicionalidades}, \esp{cooperación Sur-Sur}, \esp{adecuación}, \esp{desarrollo endógeno}.
\item \textbf{Impersonnel et discours rapporté} : \esp{El texto analiza}, \esp{denuncia}, \esp{pone como ejemplo}, \esp{se propone}, \esp{se critica} (passif réfléchi).
\item \textbf{Longueur maîtrisée} : 132 mots. Si trop long → supprimer « illustré avec le cas de Camerún... » ; si trop court → développer « énergies renouvelables, TIC, transformation agroalimentaire ».
\end{enumerate}

\textbf{Erreurs à éviter :}
\begin{itemize}
\item Dépasser 150 mots (pénalité).
\item Copier des phrases entières du texte (zéro point pour « reformulation »).
\item Oublier une partie (ex. : pas de mention des critiques → perte de points).
\item Ajouter son opinion (\esp{En mi opinión...}, \esp{Creo que...}) → hors sujet.
\item Utiliser le « je » (\esp{Yo pienso que...}) → registre inapproprié.
\end{itemize}
\end{exoresolu}

\courssub{Méthode 4 : Traduire un texte économique (Version et Thème)}

\begin{methode}[Traduction spécialisée : registre économique et institutions]
\textbf{Principes généraux :}
\begin{enumerate}
\item \textbf{Identifier le registre} : formel, précis, passif/impersonnel fréquent, vocabulaire technique.
\item \textbf{Segmenter} : couper les phrases longues en unités de sens (sujet-verbe-compléments).
\item \textbf{Traiter la terminologie} : ne pas traduire mot à mot ; chercher l'\textbf{équivalent fonctionnel} dans la langue cible (ex. : \esp{acuerdo bilateral} → « accord bilatéral » ; \esp{condicionalidades} → « conditionalités » ; \esp{brecha} → « écart / fossé »).
\item \textbf{Gérer la syntaxe} :
\begin{itemize}
\item \esp{Se + verbe} (passif réfléchi) → « on + verbe » ou passif français (\esp{Se firman acuerdos} → « On signe des accords / Des accords sont signés »).
\item \esp{Infinitif après préposition} (\esp{para diseñar}, \esp{sin generar}) → infinitif français (\esp{pour concevoir}, \esp{sans créer}).
\item \textbf{Subjonctif} après \esp{que} + volonté/nécessité/doute → indicatif ou subjonctif français selon le verbe régent.
\end{itemize}
\item \textbf{Soigner la ponctuation} : respecter les phrases longues à l'espagnol ; en français, préférer des phrases plus courtes ou des relativisées.
\item \textbf{Relire à voix haute} : vérifier la fluidité et l'exactitude terminologique.
\end{enumerate}

\textbf{Outils :} Dictionnaire bilingue économie (Routledge, Harrap's), sites officiels (UE, OIT, Banque mondiale) pour termes validés.
\end{methode}

\begin{exoresolu}[Traduction Version/Thème — Extraits du texte et paragraphes inédits]
\textbf{Version (Espagnol → Français) :}
Traduisez en français le paragraphe suivant :
\begin{quote}
\esp{« Los acuerdos multilaterales con el FMI o el Banco Mundial imponen a veces condicionalidades que reducen el gasto público en educación y salud. La cooperación Sur-Sur, impulsada por la Unión Africana y socios como Brasil o la India, ofrece una alternativa basada en el intercambio de experiencias y la transferencia de tecnología adaptada. »}
\end{quote}

\textbf{Thème (Français → Espagnol) :}
Traduisez en espagnol le paragraphe suivant :
\begin{quote}
« Le gouvernement camerounais a signé plusieurs accords bilatéraux avec des pays européens afin de financer la formation professionnelle des jeunes. Ces subventions permettent de créer des centres techniques dans les zones rurales. Cependant, les organisations non gouvernementales soulignent que l'aide ne suffit pas si les universités n'adaptent pas leurs programmes aux besoins du marché du travail. »
\end{quote}

\tcblower

\textbf{Solution Version :}
« Les accords multilatéraux avec le FMI ou la Banque mondiale imposent parfois des conditionalités qui réduisent les dépenses publiques en éducation et en santé. La coopération Sud-Sur, impulsée par l'Union africaine et des partenaires comme le Brésil ou l'Inde, offre une alternative fondée sur l'échange d'expériences et le transfert de technologie adaptée. »

\textbf{Justifications :}
\begin{itemize}
\item \esp{imponen} (indicatif présent, 3e pl.) → « imposent » (fait constaté).
\item \esp{condicionalidades} → terme technique conservé « conditionalités » (documents FMI/BM).
\item \esp{reducen} → « réduisent » (présent de narration).
\item \esp{impulsada por} (participe passé féminin singulier accordé avec \esp{cooperación}) → « impulsée par / portée par ».
\item \esp{basada en} → « fondée sur / basée sur ».
\item \esp{intercambio de experiencias} → « échange d'expériences » (collocation).
\item \esp{transferencia de tecnología adaptada} → « transfert de technologie adaptée » (adjectif postposé en espagnol, antéposé en français).
\end{itemize}

\textbf{Solution Thème :}
\esp{« El Gobierno camerunés ha firmado varios acuerdos bilaterales con países europeos con el fin de financiar la formación profesional de los jóvenes. Estas subvenciones permiten crear centros técnicos en las zonas rurales. No obstante, las organizaciones no gubernamentales señalan que la ayuda no basta si las universidades no adaptan sus programas a las necesidades del mercado laboral. »}

\textbf{Justifications grammaticales :}
\begin{enumerate}
\item \esp{El Gobierno camerunés} (sujet) + \esp{ha firmado} (passé composé = \cle{haber} + participe) → action passée achevée, lien avec présent.
\item \esp{varios acuerdos bilaterales} (COD) + \esp{con países europeos} (CC agent/compagnie).
\item \esp{con el fin de financiar} (locution finale) → « afin de financer » ; \esp{financiar} infinitif.
\item \esp{la formación profesional de los jóvenes} (nominalisation : « la formation professionnelle des jeunes »).
\item \esp{Estas subvenciones permiten crear} (\cle{permitir} + infinitif) → « permettent de créer ».
\item \esp{centros técnicos en las zonas rurales} (vocabulaire précis).
\item \esp{No obstante} (adverbe de concession, suivi de virgule) → « Cependant / Toutefois ».
\item \esp{las organizaciones no gubernamentales} (ONG) + \esp{señalan que} (verbe déclaratif + \esp{que} + \textbf{indicatif} car fait affirmé).
\item \esp{la ayuda no basta} (\cle{bastar} = suffire, 3e sg.) → « l'aide ne suffit pas ».
\item \esp{si las universidades no adaptan} (\esp{si} + indicatif présent = condition réelle possible) → « si les universités n'adaptent pas ».
\item \esp{sus programas a las necesidades del mercado laboral} (COD + \esp{a} préposition de \cle{adaptar algo a algo}).
\end{enumerate}

\textbf{Points de vigilance :}
\begin{itemize}
\item \esp{Gobierno} majuscule (institution) ; \esp{camerunés} minuscule (adjectif de nationalité).
\item \esp{ha firmado} (passé composé) et non \esp{firmó} (prétérit) car conséquence présente (centres existent).
\item \esp{No obstante} (formel) préféré à \esp{Pero} (trop oral).
\item \esp{señalan que} + indicatif (pas de subjonctif : certitude).
\item \esp{mercado laboral} (terme technique) et non \esp{mercado de trabajo} (synonyme mais moins formel ici).
\end{itemize}
\end{exoresolu}

\courssub{Méthode 5 : Rédiger un essai argumentatif / article d'opinion (Expression écrite)}

\begin{methode}[Essai argumentatif : structure « thèse-antithèse-synthèse » en espagnol]
\textbf{Consigne type Bac :} « \esp{Redacte un artículo de opinión (180–220 palabras) sobre el tema: « La cooperación internacional ¿es eficaz para el desarrollo de África? » } »

\textbf{Plan type en 4 paragraphes :}

\textbf{1. Introducción (25–35 mots) :}
\begin{itemize}
\item \textbf{Accroche} : fait marquant, chiffre, citation, question rhétorique.
\item \textbf{Sujet posé} : reformulation du sujet.
\item \textbf{Annonce du plan} : \esp{Analizaremos... ; Examinaremos... ; Finalmente, propondremos...}
\end{itemize}

\textbf{2. Desarrollo — Tesis (arguments POUR) :}
\begin{itemize}
\item Idée 1 + connecteur (\esp{En primer lugar}, \esp{Es innegable que...}) + exemple concret (Cameroun : centres de formation, routes, santé).
\item Idée 2 + connecteur (\esp{Además}, \esp{Por otra parte}) + exemple (microcrédit, ONG, coopération technique).
\item Vocabulaire : \esp{aportar}, \esp{facilitar}, \esp{fortalecer}, \esp{capacitar}, \esp{infraestructura}.
\end{itemize}

\textbf{3. Desarrollo — Antítesis (arguments CONTRE / limites) :}
\begin{itemize}
\item Idée 1 + connecteur (\esp{No obstante}, \esp{Sin embargo}, \esp{En cambio}) + contre-exemple (dette, conditionalités, projets inadaptés).
\item Idée 2 + connecteur (\esp{Tampoco hay que olvidar que...}) + critique (dépendance, gouvernance, corruption).
\item Vocabulaire : \esp{condicionar}, \esp{endeudar}, \esp{extraer}, \esp{imponer}, \esp{dependencia}.
\end{itemize}

\textbf{4. Conclusión — Síntesis (30–40 mots) :}
\begin{itemize}
\item \textbf{Bilan nuancé} : \esp{En definitiva}, \esp{En resumen}, \esp{En conclusión}.
\item \textbf{Position personnelle} : \esp{Considero que...}, \esp{Es fundamental que...}, \esp{Urge...}.
\item \textbf{Perspective d'avenir} : \esp{cooperación Sur-Sur}, \esp{adecuación formación-empleo}, \esp{desarrollo endógeno}.
\end{itemize}

\textbf{Connecteurs indispensables à mémoriser :}
\begin{center}
\begin{tabular}{|l|l|}
\hline
\esp{En primer lugar / En segundo lugar} & Premièrement / Deuxièmement \\
\esp{Por un lado... Por otro lado} & D'un côté... De l'autre \\
\esp{No obstante / Sin embargo} & Cependant / Toutefois \\
\esp{Por consiguiente / Por tanto} & Par conséquent / Donc \\
\esp{En definitiva / En conclusión} & En définitive / En conclusion \\
\hline
\end{tabular}
\end{center}
\end{methode}

\begin{textebox}[Artículo de opinión — ¿Es eficaz la cooperación internacional?]
¿Es eficaz la cooperación internacional para el desarrollo de África? La pregunta divide a expertos y ciudadanos. En primer lugar, es innegable que la ayuda externa ha permitido avances concretos: en Camerún, los acuerdos bilaterales con la UE han financiado centros de formación profesional rurales, y las ONG apoyan cooperativas de cacao certificado que mejoran los ingresos de miles de familias. Además, el microcrédito ha dado autonomía a mujeres emprendedoras en regiones aisladas.

No obstante, la eficacia no es automática. Sin embargo, muchos proyectos fracasan por falta de adecuación con las realidades locales: se imponen modelos extranjeros sin consultar a las comunidades. Por otra parte, la inversión extranjera directa suele centrarse en la extracción de petróleo o madera, generando pocos empleos cualificados. Las condicionalidades del FMI, a veces, obligan a recortar el gasto en educación y salud, debilitando el capital humano.

En definitiva, la cooperación solo es eficaz si respeta la soberanía y prioriza el desarrollo endógeno. Considero que la cooperación Sur-Sur y el diálogo universidad-empresa son claves. Urge transformar la ayuda en palanca, no en muleta, para que África sea dueña de su futuro.
\end{textebox}

\begin{exoresolu}[Rédaction d'un article d'opinion — Sujet Bac]
\textbf{Sujet :} \esp{« La cooperación internacional ¿es eficaz para el desarrollo de África? »} (180–220 mots)

\textbf{Comptage :} 198 mots (dans la fourchette).

\tcblower

\textbf{Analyse de la rédaction (grille d'évaluation Bac) :}

\begin{center}
\begin{tabularx}{\linewidth}{|X|X|X|}
\hline
\rowcolor{popblueL} \textbf{Critère} & \textbf{Réalisation dans le modèle} & \textbf{Points forts} \\
\hline
Respect du sujet & Réponse directe à la question, nuancée & Thèse/antithèse/synthèse claire \\
Structure & 4 paragraphes logiques, connecteurs variés & \esp{En primer lugar}, \esp{No obstante}, \esp{Por otra parte}, \esp{En definitiva} \\
Vocabulaire & Lexique spécialisé précis & \esp{adecuación}, \esp{desarrollo endógeno}, \esp{condicionalidades}, \esp{capital humano}, \esp{soberanía} \\
Grammaire & Temps justifiés, subjonctifs si nécessaire & \esp{ha permitido} (passé composé), \esp{se imponen} (indicatif), \esp{sea} (subj. après \esp{para que}) \\
Cohérence & Exemples camerounais authentiques & Centres ruraux, cacao, microcrédit, FMI \\
Registre & Formel, impersonnel, argumentatif & Pas de « je » avant conclusion, pas de familier \\
Longueur & 198 mots & Dans les clous \\
\hline
\end{tabularx}
\end{center}

\textbf{Variantes possibles pour l'antithèse :}
\begin{itemize}
\item \esp{Tampoco podemos ignorar la corrupción que desvía fondos...}
\item \esp{En cambio, la dependencia de la ayuda frena la recaudación fiscal interna...}
\end{itemize}

\textbf{Variantes pour la conclusion :}
\begin{itemize}
\item \esp{Es fundamental que los países africanos definan sus propias prioridades.}
\item \esp{El futuro pasa por una cooperación de igual a igual, no de donante a receptor.}
\end{itemize}


\end{exoresolu}

\courssub{Méthode 6 : Exprimer l'opinion et argumenter à l'oral (Production orale guidée)}

\begin{methode}[Prise de parole en continu : exposé / débat — 3 à 5 minutes]
\textbf{Situation :} « \esp{Presentas en clase tu análisis sobre la eficacia de la ayuda al desarrollo en Camerún, apoyándote en el texto y en tu conocimiento del contexto nacional.} »

\textbf{Structure de l'exposé (fiche bristol autorisée : mots-clés seulement) :}
\begin{enumerate}
\item \textbf{Accroche (30 s) :} \esp{Buenos días. ¿Sabían que más del 60 \% de los jóvenes africanos trabajan en la informalidad? Hoy analizaré si la cooperación internacional resuelve este reto en Camerún.}
\item \textbf{Constat (1 min) :} \esp{El texto « El mundo laboral hoy » muestra que...} (2–3 idées-chiffres : informel, diplômés au chômage, \emph{bendskin}).
\item \textbf{Acteurs et instruments (1 min) :} \esp{La ayuda llega vía acuerdos bilaterales (ej. España-Camerún), subvenciones UE, ONG...} (citer 1 exemple précis par acteur).
\item \textbf{Limites critiques (1 min) :} \esp{Pero hay sombras: inversión extractiva, condicionalidades FMI, proyectos sin sostenibilidad.}
\item \textbf{Proposition / Avenir (30 s) :} \esp{La clave es la adecuación formación-mercado. Las universidades deben hablar con empresas. La cooperación Sur-Sur ofrece modelos más adaptados.}
\item \textbf{Conclusion (30 s) :} \esp{En conclusión, la ayuda es útil si África manda. Gracias.}
\end{enumerate}

\textbf{Techniques orales :}
\begin{itemize}
\item \textbf{Regard} : balayer la classe, ne pas lire.
\item \textbf{Voix} : volume, débit modéré, pauses aux connecteurs.
\item \textbf{Appuis visuels} : carte mentale au tableau (mots-clés : \esp{INFORMEL}, \esp{ACUERDOS}, \esp{ONG}, \esp{LÍMITES}, \esp{SOLUCIONES}).
\item \textbf{Interaction} : prévoir 2 questions aux camarades (\esp{¿Qué opinan ustedes?}, \esp{¿Conocen algún proyecto de ONG en su barrio?}).
\end{itemize}

\textbf{Évaluation (grille) :} Contenu (5) + Structure (3) + Lexique/Grammaire (4) + Prononciation/Prosodie (3) + Interaction (2) = /17.
\end{methode}

\begin{exoresolu}[Simulation orale — Fiche élève + correction attendue]
\textbf{Consigne :} « \esp{Prepara una intervención oral de 3 minutos sobre: « El papel de las ONG en el desarrollo local en Camerún ». Utiliza el vocabulario de la lección.} »

\textbf{Fiche bristol (mots-clés autorisés) :}
\begin{center}
\begin{tabular}{|p{0.45\linewidth}|p{0.45\linewidth}|}
\hline
\textbf{INTRO} & \esp{Buenos días. Tema: ONG y desarrollo local en Camerún.} \\
\hline
\textbf{DATOS} & \esp{60\% informal (OIT). Jóvenes: bendskin, comercio.} \\
\hline
\textbf{ONG - ACCIONES} & \esp{Formación mujeres agroindustria. Cooperativas cacao certificado. Microcrédito.} \\
\hline
\textbf{EJEMPLOS CONCRETOS} & \esp{ONG local « Horizons Verts » (Dschang): 200 mujeres formadas. SNV (internacional): cacao exportación.} \\
\hline
\textbf{LÍMITES} & \esp{Proyectos cortos (2-3 años). Dependencia fondos externos. Falta coordinación Estado.} \\
\hline
\textbf{PROPUESTA} & \esp{Institucionalizar alianzas ONG-Estado-Universidad. Fondos nacionales (FONADER).} \\
\hline
\textbf{CONCLUSIÓN} & \esp{ONG son motor, no sustituto del Estado. Desarrollo endógeno = clave.} \\
\hline
\end{tabular}
\end{center}

\textbf{Production orale attendue (transcription modèle) :}
\begin{quote}
\esp{Buenos días a todos. Mi exposición trata sobre el papel de las ONG en el desarrollo local en Camerún. Como señala el texto « El mundo laboral hoy », más del 60 \% de la población activa trabaja en la economía informal. Ante esto, las ONG —tanto locales como internacionales— intervienen donde el Estado no llega suficientemente.

En primer lugar, forman a mujeres en agroindustria: por ejemplo, la ONG local « Horizons Verts » en Dschang ha capacitado a más de 200 mujeres en transformación de plátano y mango. En segundo lugar, apoyan cooperativas de cacao certificado, como hace SNV, lo que permite exportar a mejor precio. En tercer lugar, facilitan el microcrédito para pequeños emprendimientos rurales.

No obstante, existen límites. Los proyectos suelen ser cortos —dos o tres años— y dependen de fondos externos. A veces falta coordinación con las políticas públicas, lo que crea duplicados.

Propongo, por tanto, institucionalizar las alianzas entre ONG, Estado y universidades, y crear fondos nacionales como el FONADER para dar continuidad. En conclusión, las ONG son un motor esencial, pero no deben sustituir al Estado. El desarrollo endógeno es la clave. Gracias.}
\end{quote}

\tcblower

\textbf{Analyse de la performance (auto-évaluation) :}
\begin{itemize}
\item \textbf{Temps :} ~3 min 15 s (bon).
\item \textbf{Vocabulaire :} \esp{capacitar}, \esp{cooperativa}, \esp{certificado}, \esp{microcrédito}, \esp{institucionalizar}, \esp{desarrollo endógeno} — \textbf{excellent}.
\item \textbf{Grammaire :} \esp{ha capacitado} (passé composé), \esp{permite exportar} (infinitif), \esp{suelen ser} (présent habituel), \esp{falta coordinación} (impersonnel), \esp{deben sustituir} (obligation) — \textbf{correct}.
\item \textbf{Connecteurs :} \esp{En primer lugar}, \esp{En segundo lugar}, \esp{No obstante}, \esp{Por tanto}, \esp{En conclusión} — \textbf{maîtrisés}.
\item \textbf{Prononciation :} attention à \esp{NG} [eneˈxe] (ONG), \esp{certificado} [θer.ti.fi.ˈka.ðo] / [ser.ti.fi.ˈka.ðo], \esp{endógeno} [en.ˈdo.xe.no] (accent sur \emph{ó}).
\item \textbf{Amélioration possible :} ajouter une statistique vérifiable (ex. \esp{« Según el MINADER, el cacao certificado representa ya el 15 \% de las exportaciones »}) pour renforcer la crédibilité.
\end{itemize}
\end{exoresolu}

\courssub{Les 5 erreurs qui coûtent des points au Bac}

\begin{attention}[Les 5 erreurs fatales — Spécial « Monde du travail \& Aide au développement »]
\textbf{1. Accents oubliés ou mal placés (orthographe) :}
\begin{itemize}
\item \esp{cooperación} (accent sur \emph{ó}) ≠ \esp{cooperacion} ×
\item \esp{inversión} (accent sur \emph{ó}) ≠ \esp{inversion} ×
\item \esp{subvención} (accent sur \emph{ó}) ≠ \esp{subvencion} ×
\item \esp{acuerdo} (diphtongue \emph{ue}) ≠ \esp{acuerdó} × (pas d'accent sur monosyllabe)
\item \esp{país} (accent sur \emph{í} : hiatus) ≠ \esp{pais} × (signifie « païs » → pays ; sans accent = « païs » incorrect)
\item \esp{agroindustria} (accent sur \emph{ú} : esdrújula) ≠ \esp{agroindustria} ×
\end{itemize}
\cle{Règle} : Tous les mots oxytons (accent sur dernière syllabe) terminés par \emph{n, s, voyelle} portent l'accent écrit. \esp{cooperación, inversión, subvención, país, también, además}. Les mots esdrújules (accent sur l'antépénultième) portent toujours l'accent : \esp{agroindustria, microcrédito, desarrollo, endógeno}.

\textbf{2. Faux amis « entreprise / emploi / compromis » :}
\begin{center}
\begin{tabularx}{\linewidth}{|X|X|X|}
\hline
\rowcolor{poporangeL} \textbf{Français} & \textbf{Espagnol correct} & \textbf{Piège (faux ami)} \\
\hline
Entreprise & \esp{empresa} & \esp{empresa} ≠ emploi \\
Emploi / Travail & \esp{empleo / trabajo} & \esp{trabajo} ≠ travail (au sens ouvrage) \\
Compromis & \esp{compromiso} (engagement) & \esp{compromiso} ≠ compromis (accord) \\
Subvention & \esp{subvención} & \esp{subvención} ≠ soumission \\
Investissement & \esp{inversión} & \esp{inversión} ≠ inversion (renversement) \\
Actual / Actuel & \esp{actual} / \esp{real} & \esp{actual} = actuel (présent) ; \esp{real} = réel \\
\hline
\end{tabularx}
\end{center}

\textbf{3. Calques syntaxiques du français (structure) :}
\begin{itemize}
\item × \esp{La ayuda es útil para desarrollar África.} → ✓ \esp{La ayuda es útil para el desarrollo de África.} (nominalisation : \cle{para} + nom, pas infinitif après préposition pour but général).
\item × \esp{Los acuerdos bilaterales son firmados por Camerún y España.} (passif « à la française ») → ✓ \esp{Camerún y España firman acuerdos bilaterales.} (actif) ou \esp{Se firman acuerdos bilaterales entre Camerún y España.} (passif réfléchi naturel).
\item × \esp{Es importante que las universidades adaptan sus programas.} → ✓ \esp{Es importante que las universidades \textbf{adapten} sus programas.} (\textbf{subjonctif} après \esp{Es importante que}).
\item × \esp{Si la cooperación será eficaz, África se desarrollará.} → ✓ \esp{Si la cooperación \textbf{es} eficaz, África se desarrollará.} (\textbf{indicatif présent} après \esp{si} + futur, pas de futur dans le \esp{si}).
\end{itemize}

\textbf{4. Accords (genre/nombre) oubliés sur participes et adjectifs :}
\begin{itemize}
\item × \esp{Las subvenciones son canalizado por la UE.} → ✓ \esp{canalizad\textbf{as}} (fém. pl.).
\item × \esp{La cooperación Sur-Sur es impulsado por la UA.} → ✓ \esp{impulsad\textbf{a}} (fém. sg.).
\item × \esp{Los acuerdos multilaterales son firmado.} → ✓ \esp{firmad\textbf{os}} (masc. pl.).
\item × \esp{Una inversión extranjera directos.} → ✓ \esp{direct\textbf{a}} (fém. sg.).
\end{itemize}
\cle{Règle} : Participe passé = adjectif → s'accorde en genre et nombre avec le sujet (si passif) ou le COD antéposé.

\textbf{5. Choix du temps / aspect verbal (passé composé vs prétérit / imparfait) :}
\begin{itemize}
\item Contexte : action passée \textbf{achevée}, datée ou à conséquence présente → \textbf{Passé composé} (\esp{ha firmado}, \esp{han financiado}).
\item Contexte : action passée \textbf{ponctuelle}, narrée, sans lien présent → \textbf{Prétérit} (\esp{firmó}, \esp{financió}).
\item Contexte : description, habitude, arrière-plan dans le passé → \textbf{Imparfait} (\esp{había}, \esp{era}, \esp{trabajaba}).
\item \textbf{Piège Bac :} « Le gouvernement \textbf{a signé} (conséquence : centres existent) des accords » → \esp{ha firmado} (pas \esp{firmó}).
\end{itemize}

\textbf{Astuce de révision :} Avant de rendre ta copie, \textbf{relis à l'envers} (dernière phrase → première) : l'œil repère mieux les accords manquants et les accents oubliés.
\end{attention}

\newpage

\coursec{Exercices}

\courssub{Je vérifie mes connaissances}

\begin{exercice}[QCM : le lexique de la coopération]
Pour chaque question, entoure la bonne réponse.
\begin{enumerate}
\item Une \esp{subvención} est :
\begin{enumerate}
\item une dette que l'État doit rembourser ;
\item une aide financière accordée sans obligation de remboursement ;
\item un impôt sur les entreprises étrangères.
\end{enumerate}
\item Un \esp{acuerdo bilateral} est signé entre :
\begin{enumerate}
\item deux pays ; \item trois continents ; \item une ONG et une banque.
\end{enumerate}
\item Une \esp{ONG} est :
\begin{enumerate}
\item une entreprise multinationale ;
\item une organisation non gouvernementale ;
\item un ministère espagnol.
\end{enumerate}
\item \esp{La inversión extranjera} désigne :
\begin{enumerate}
\item l'argent placé dans un pays par des acteurs étrangers ;
\item l'exportation de cacao ;
\item le chômage des jeunes diplômés.
\end{enumerate}
\end{enumerate}
\end{exercice}

\begin{exercice}[Vrai ou faux : justifie par une citation]
D'après le texte \esp{El mundo laboral hoy} étudié en classe, dis si les affirmations suivantes sont vraies (V) ou fausses (F) et justifie chaque réponse par une courte citation du texte en espagnol.
\begin{enumerate}
\item Le monde du travail n'a pas changé ces dernières années.
\item La coopération internationale joue un rôle dans le développement économique de l'Afrique.
\item Les accords multilatéraux ne concernent que l'Europe.
\item La formation professionnelle est un enjeu important pour les jeunes.
\end{enumerate}
\end{exercice}

\begin{exercice}[Vocabulaire : trouve le mot juste]
Complète chaque phrase avec le mot du lexique qui convient : \esp{ayuda al desarrollo}, \esp{cooperación}, \esp{inversión}, \esp{acuerdo multilateral}, \esp{subvención}.
\begin{enumerate}
\item España y Guinea Ecuatorial mantienen relaciones de \ldots\ldots\ldots\ldots\ldots\ldots{} desde hace muchos años.
\item El gobierno concedió una \ldots\ldots\ldots\ldots\ldots\ldots{} a los jóvenes agricultores de la región del Centre.
\item La \ldots\ldots\ldots\ldots\ldots\ldots{} extranjera crea empleo en las ciudades africanas.
\item Varios países firmaron un \ldots\ldots\ldots\ldots\ldots\ldots{} para luchar contra el cambio climático.
\item La \ldots\ldots\ldots\ldots\ldots\ldots{} debe llegar directamente a la población más pobre.
\end{enumerate}
\end{exercice}

\begin{exercice}[Conjugaison : le subjonctif après \esp{para que}]
Complète les phrases avec le verbe entre parenthèses au subjonctif présent.
\begin{enumerate}
\item Las ONG trabajan sobre el terreno para que la ayuda \ldots\ldots\ldots\ldots\ldots\ldots{} (llegar) a la población.
\item El gobierno invierte en escuelas para que los jóvenes \ldots\ldots\ldots\ldots\ldots\ldots{} (tener) una buena formación.
\item Firmamos este acuerdo para que las empresas españolas \ldots\ldots\ldots\ldots\ldots\ldots{} (poder) invertir en Camerún.
\item Es importante que tú \ldots\ldots\ldots\ldots\ldots\ldots{} (saber) hablar varias lenguas en el mundo laboral.
\item Los países cooperan para que no \ldots\ldots\ldots\ldots\ldots\ldots{} (haber) más pobreza.
\end{enumerate}
\end{exercice}

\courssub{Je m'entraîne}

\begin{exercice}[Traduction : thème]
Traduis les phrases suivantes en espagnol.
\begin{enumerate}
\item Les ONG travaillent sur le terrain pour que l'aide arrive à la population.
\item Le Cameroun a signé un accord bilatéral avec l'Espagne.
\item Le gouvernement accorde des subventions aux jeunes entrepreneurs.
\item L'investissement étranger est important pour le développement de l'Afrique.
\item Il faut que la coopération internationale lutte contre le chômage des jeunes.
\end{enumerate}
\end{exercice}

\begin{exercice}[Texte à trous : un accord de coopération]
Complète le texte suivant avec les mots de la liste : \esp{subvención}, \esp{inversión}, \esp{ONG}, \esp{acuerdo multilateral}, \esp{cooperación}, \esp{desarrollo}.

\begin{textebox}[Un acuerdo para África]
Ayer se firmó en Madrid un \ldots\ldots\ldots\ldots\ldots\ldots{} entre varios países europeos y africanos. Este acuerdo prevé una \ldots\ldots\ldots\ldots\ldots\ldots{} económica importante para financiar proyectos de \ldots\ldots\ldots\ldots\ldots\ldots{} rural. Además, cada país africano recibirá una \ldots\ldots\ldots\ldots\ldots\ldots{} para formar a los jóvenes agricultores. Varias \ldots\ldots\ldots\ldots\ldots\ldots{} internacionales controlarán que el dinero llegue a la población. La \ldots\ldots\ldots\ldots\ldots\ldots{} entre los dos continentes es hoy más necesaria que nunca.
\end{textebox}
\end{exercice}

\begin{exercice}[Transformation de phrases]
Transforme les phrases selon la consigne entre parenthèses, sans changer le sens.
\begin{enumerate}
\item \esp{El gobierno ayuda a los jóvenes.} (transforme avec \esp{Es importante que} + subjonctif)
\item \esp{Las empresas invierten en África para crear empleo.} (transforme avec \esp{para que})
\item \esp{España firmó un acuerdo con Camerún.} (mets la phrase au passé composé : \esp{pretérito perfecto})
\item \esp{La ayuda internacional es suficiente.} (mets la phrase à la forme négative avec ton opinion : \esp{No creo que} + subjonctif)
\item \esp{Los jóvenes buscan trabajo.} (ajoute \esp{Es posible que} et fais les changements nécessaires)
\end{enumerate}
\end{exercice}

\begin{exercice}[Appariements : mots et définitions]
Relie chaque mot espagnol (colonne A) à sa définition (colonne B).

\begin{center}
\begin{tabularx}{\linewidth}{|X|X|}
\hline
\rowcolor{popblueL} \textbf{A. Palabra} & \textbf{B. Definición} \\
\hline
1. la subvención & a. dinero que se coloca en un país para crear riqueza \\
\hline
2. la ONG & b. tratado firmado por dos Estados \\
\hline
3. la inversión & c. dinero dado por el Estado sin obligación de devolverlo \\
\hline
4. el acuerdo bilateral & d. conjunto de acciones entre países para ayudarse \\
\hline
5. la cooperación & e. organización independiente que ayuda a la población \\
\hline
6. el desempleo & f. situación de las personas que no tienen trabajo \\
\hline
\end{tabularx}
\end{center}
\end{exercice}

\courssub{Je me prépare au Bac}

\begin{exercice}[Compréhension de texte type Bac]
Lis attentivement le texte suivant, puis réponds aux questions en espagnol, avec des phrases complètes.

\begin{textebox}[España y Camerún firman un acuerdo de cooperación]
Ayer, en el Palacio de la Unidad de Yaundé, los gobiernos de España y de Camerún firmaron un nuevo acuerdo bilateral de cooperación económica. Según este acuerdo, España concederá una subvención destinada a la formación profesional de los jóvenes cameruneses, especialmente en los sectores de la agricultura y de la tecnología.

El ministro camerunés declaró que esta ayuda al desarrollo permitirá crear miles de empleos en las regiones del Centro y del Litoral. Por su parte, la embajadora española insistió en la importancia de la inversión en la educación: « Sin formación, no hay desarrollo posible », afirmó.

Varias ONG locales participarán en el proyecto para garantizar que la ayuda llegue directamente a la población rural. Sin embargo, algunos especialistas piensan que la ayuda internacional no es suficiente: creen que África necesita sobre todo más inversión extranjera y un comercio más justo con Europa.

Este acuerdo forma parte de una cooperación más amplia entre la Unión Europea y los países africanos, que incluye también acuerdos multilaterales sobre el medio ambiente y la salud.
\end{textebox}

\textbf{Questions de compréhension} (6 points) :
\begin{enumerate}
\item ¿Dónde se firmó el acuerdo y entre qué países?
\item ¿A qué sectores está destinada la subvención española?
\item ¿Qué papel tendrán las ONG en este proyecto?
\item ¿Qué piensan algunos especialistas de la ayuda internacional?
\end{enumerate}

\textbf{Questions de langue} (4 points) :
\begin{enumerate}
\item Trouve dans le texte : un synonyme de \esp{trabajo} et le contraire de \esp{posible}.
\item Mets à l'infinitif les verbes suivants tirés du texte : \esp{firmaron}, \esp{declaró}, \esp{afirmó}, \esp{creen}.
\item Complète avec le subjonctif : \esp{Es importante que la ayuda \ldots\ldots\ldots\ldots\ldots\ldots{} (llegar) a la población rural.}
\item Traduis en français la phrase : \esp{Sin formación, no hay desarrollo posible.}
\end{enumerate}
\end{exercice}

\begin{exercice}[Thème et version]
\textbf{A. Thème} (5 points). Traduis en espagnol :
\begin{enumerate}
\item La coopération entre l'Espagne et l'Afrique est très importante.
\item Il faut que les jeunes reçoivent une bonne formation professionnelle.
\item Les accords multilatéraux aident les pays pauvres à se développer.
\item Le chômage des jeunes est un grand problème au Cameroun.
\item Nous espérons que cet accord créera beaucoup d'emplois.
\end{enumerate}

\textbf{B. Version} (5 points). Traduis en français le texte suivant :

\begin{textebox}[La ayuda al desarrollo]
La ayuda al desarrollo no consiste solamente en dar dinero a los países pobres. Consiste también en compartir conocimientos, formar a los jóvenes y construir escuelas y hospitales. Las ONG desempeñan un papel fundamental porque trabajan directamente con la población. Sin embargo, para que la ayuda sea eficaz, es necesario que los gobiernos luchen contra la corrupción y que los países ricos respeten sus promesas.
\end{textebox}
\end{exercice}

\begin{exercice}[Expression écrite type Bac]
Choisis \textbf{un} des deux sujets suivants et rédige un texte en espagnol de \textbf{80 à 100 mots}. Soigne la présentation, la grammaire (temps, subjonctif) et le vocabulaire de la leçon.

\textbf{Sujet 1.} \esp{¿Piensas que la ayuda internacional es suficiente para el desarrollo de África? Justifica tu respuesta con argumentos y ejemplos.}

\textbf{Sujet 2.} \esp{El chômage des jeunes es un problema grave en muchos países africanos. En tu opinión, ¿qué soluciones pueden aportar los gobiernos, las empresas y la cooperación internacional?}

\emph{Conseils méthodologiques :} rédige une courte introduction (présentation du sujet et de ton opinion), un développement avec deux ou trois arguments illustrés d'exemples (Cameroun, Guinée équatoriale, accords avec l'Espagne), et une conclusion. Utilise des connecteurs : \esp{en primer lugar}, \esp{además}, \esp{sin embargo}, \esp{por eso}, \esp{en conclusión}. Emploie au moins une fois le subjonctif après \esp{para que} ou \esp{es necesario que}.
\end{exercice}

\newpage

\coursec{Corrigés des exercices}

\coursec{Corrigés des exercices — Leçon E16 : El mundo laboral hoy}

\begin{corrige}[Exercice 1 — QCM : le lexique de la coopération]
\begin{enumerate}
\item \textbf{b.} Une \esp{subvención} est une aide financière accordée sans obligation de remboursement. C'est la différence essentielle avec un prêt (\esp{préstamo}), qui, lui, doit être remboursé.
\item \textbf{a.} Un \esp{acuerdo bilateral} est signé entre deux pays. Le préfixe \esp{bi-} signifie « deux » ; \esp{multilateral} concerne plusieurs pays.
\item \textbf{b.} Une \esp{ONG} est une organisation non gouvernementale (\esp{organización no gubernamental}).
\item \textbf{a.} \esp{La inversión extranjera} désigne l'argent placé dans un pays par des acteurs étrangers (entreprises, États, fonds).
\end{enumerate}
\end{corrige}

\begin{attention}[Point de vigilance — Exercice 1]
Ne confondez pas \esp{subvención} et \esp{préstamo} : la subvention n'est pas remboursable, le prêt l'est. Attention aussi au faux ami : \esp{inversión} = « investissement », et non « inversion ».
\end{attention}

\begin{corrige}[Exercice 2 — Vrai ou faux : justifie par une citation]
\begin{enumerate}
\item \textbf{Faux.} Le texte affirme au contraire : \esp{« El mundo laboral ha cambiado mucho en los últimos años »} (ou toute citation équivalente du texte étudié montrant la transformation du travail).
\item \textbf{Vrai.} Justification : \esp{« La cooperación internacional contribuye al desarrollo económico de África »} (citation du passage consacré à l'aide au développement).
\item \textbf{Faux.} Le texte précise que \esp{« los acuerdos multilaterales reúnen a países de varios continentes »} : ils ne concernent donc pas seulement l'Europe.
\item \textbf{Vrai.} Justification : \esp{« La formación profesional es un desafío importante para los jóvenes »} (citation du passage sur l'emploi des jeunes).
\end{enumerate}
\emph{Remarque méthodologique :} au Bac, la justification par une citation doit être \textbf{courte, exacte} (recopiée sans erreur d'accent) et \textbf{pertinente} : elle doit prouver directement votre réponse.
\end{corrige}

\begin{attention}[Point de vigilance — Exercice 2]
Veillez à recopier les accents exactement (\esp{ú}, \esp{ó}, \esp{á}, \esp{é}, \esp{í}, \esp{ñ}). Une citation tronquée ou mal accentuée perd les points de langue.
\end{attention}

\begin{corrige}[Exercice 3 — Vocabulaire : trouve le mot juste]
\begin{enumerate}
\item España y Guinea Ecuatorial mantienen relaciones de \textbf{cooperación} desde hace muchos años. (« L'Espagne et la Guinée équatoriale entretiennent des relations de coopération depuis de nombreuses années. »)
\item El gobierno concedió una \textbf{subvención} a los jóvenes agricultores de la región del Centro. (« Le gouvernement a accordé une subvention aux jeunes agriculteurs de la région du Centre. »)
\item La \textbf{inversión} extranjera crea empleo en las ciudades africanas. (« L'investissement étranger crée de l'emploi dans les villes africaines. »)
\item Varios países firmaron un \textbf{acuerdo multilateral} para luchar contra el cambio climático. (« Plusieurs pays ont signé un accord multilatéral pour lutter contre le changement climatique. »)
\item La \textbf{ayuda al desarrollo} debe llegar directamente a la población más pobre. (« L'aide au développement doit parvenir directement à la population la plus pauvre. »)
\end{enumerate}
\emph{Astuce :} vérifiez toujours l'accord de l'article (\esp{una subvención}, \esp{la inversión}, \esp{un acuerdo}) : il confirme le genre du mot cherché.
\end{corrige}

\begin{corrige}[Exercice 4 — Conjugaison : le subjonctif après \esp{para que}]
\begin{enumerate}
\item \ldots{} para que la ayuda \textbf{llegue} a la población. (\esp{llegar} : radical \esp{llegu-} devant \esp{-e} pour garder le son [gué] ; orthographe \esp{llegue}, jamais \esp{legeue}.)
\item \ldots{} para que los jóvenes \textbf{tengan} una buena formación. (\esp{tener}, irrégulier : \esp{tenga, tengas, tenga, tengamos, tengáis, tengan}.)
\item \ldots{} para que las empresas españolas \textbf{puedan} invertir en Camerún. (\esp{poder} : changement de voyelle \esp{o} $\rightarrow$ \esp{ue} : \esp{pueda, puedas, pueda, podamos, podáis, puedan}.)
\item Es importante que tú \textbf{sepas} hablar varias lenguas. (\esp{saber}, irrégulier : \esp{sepa, sepas, sepa, sepamos, sepáis, sepan}.)
\item \ldots{} para que no \textbf{haya} más pobreza. (\esp{haber} impersonnel : \esp{haya} ; jamais \esp{haiga}.)
\end{enumerate}
\end{corrige}

\begin{attention}[Point de vigilance — Exercice 4]
Après \esp{para que}, le subjonctif est obligatoire lorsque les deux sujets sont différents (\esp{las ONG} / \esp{la ayuda}). Si le sujet est identique, on emploie \esp{para} + infinitif : \esp{Trabajo para ayudar a mi familia.} Attention aux radicaux irréguliers : \esp{tengan} (et non \esp{tenan}), \esp{haya} (et non \esp{haba}).
\end{attention}

\begin{corrige}[Exercice 5 — Traduction : thème]
\begin{enumerate}
\item \esp{Las ONG trabajan sobre el terreno para que la ayuda llegue a la población.} (Subjonctif \esp{llegue} obligatoire après \esp{para que}.)
\item \esp{Camerún ha firmado un acuerdo bilateral con España.} (\esp{Camerún} prend un accent sur le \esp{ú} en espagnol ; \esp{ha firmado} : passé composé, auxiliaire \esp{haber} + participe.)
\item \esp{El gobierno concede subvenciones a los jóvenes emprendedores.} (Attention : « entrepreneur » se dit \esp{emprendedor} ou \esp{empresario}, jamais \esp{entrepreneur}.)
\item \esp{La inversión extranjera es importante para el desarrollo de África.} (\esp{África} avec accent sur le \esp{í} ; pas d'article devant \esp{África}.)
\item \esp{Es necesario que la cooperación internacional luche contra el desempleo de los jóvenes.} (« Il faut que » = \esp{es necesario que} + subjonctif : \esp{luche}. « Chômage » = \esp{desempleo} ou \esp{paro}.)
\end{enumerate}
\end{corrige}

\begin{attention}[Faux amis et calques — Exercice 5]
« Important » se traduit \esp{importante} (pas de problème), mais « développement » = \esp{desarrollo} (un seul \esp{l} à la fin, deux \esp{r} au milieu : \esp{desa\textbf{rr}ollo}). « Lutter contre » = \esp{luchar contra}, sans préposition supplémentaire.
\end{attention}

\begin{corrige}[Exercice 6 — Texte à trous : un accord de coopération]
\begin{textebox}[Un acuerdo para África — texto completo]
Ayer se firmó en Madrid un \textbf{acuerdo multilateral} entre varios países europeos y africanos. Este acuerdo prevé una \textbf{inversión} económica importante para financiar proyectos de \textbf{desarrollo} rural. Además, cada país africano recibirá una \textbf{subvención} para formar a los jóvenes agricultores. Varias \textbf{ONG} internacionales controlarán que el dinero llegue a la población. La \textbf{cooperación} entre los dos continentes es hoy más necesaria que nunca.
\end{textebox}
\emph{Traduction :} « Hier, un accord multilatéral a été signé à Madrid entre plusieurs pays européens et africains. Cet accord prévoit un investissement économique important pour financer des projets de développement rural. De plus, chaque pays africain recevra une subvention pour former les jeunes agriculteurs. Plusieurs ONG internationales veilleront à ce que l'argent parvienne à la population. La coopération entre les deux continents est aujourd'hui plus nécessaire que jamais. »
\emph{Indices de résolution :} \esp{un} + masculin singulier $\rightarrow$ \esp{acuerdo multilateral} ; \esp{una \ldots{} económica} $\rightarrow$ \esp{inversión} ; \esp{proyectos de \ldots{} rural} $\rightarrow$ \esp{desarrollo} ; \esp{varias} + féminin pluriel $\rightarrow$ \esp{ONG}.
\end{corrige}

\begin{corrige}[Exercice 7 — Transformation de phrases]
\begin{enumerate}
\item \esp{Es importante que el gobierno ayude a los jóvenes.} (\esp{Es importante que} + subjonctif : \esp{ayude}.)
\item \esp{Las empresas invierten en África para que se cree empleo.} (ou \esp{para que los jóvenes encuentren empleo} ; \esp{para que} + subjonctif car le sujet change.)
\item \esp{España ha firmado un acuerdo con Camerún.} (\esp{Pretérito perfecto} : \esp{haber} au présent + participe passé \esp{firmado}, invariable.)
\item \esp{No creo que la ayuda internacional sea suficiente.} (\esp{No creo que} exprime le doute $\rightarrow$ subjonctif \esp{sea}, irrégulier de \esp{ser}.)
\item \esp{Es posible que los jóvenes busquen trabajo.} (\esp{buscar} : \esp{c} $\rightarrow$ \esp{qu} devant \esp{-e} : \esp{busquen}.)
\end{enumerate}
\end{corrige}

\begin{attention}[Point de vigilance — Exercice 7]
\esp{Creo que} + indicatif (\esp{Creo que es suficiente}), mais \esp{No creo que} + subjonctif (\esp{No creo que sea suficiente}). C'est une règle classique du Bac : l'opinion affirmative garde l'indicatif, l'opinion niée exige le subjonctif.
\end{attention}

\begin{corrige}[Exercice 8 — Appariements : mots et définitions]
\begin{center}
\begin{tabularx}{\linewidth}{|X|X|}
\hline
\rowcolor{popblueL} \textbf{A. Palabra} & \textbf{B. Definición} \\
\hline
1. la subvención & \textbf{c.} dinero dado por el Estado sin obligación de devolverlo \\
\hline
2. la ONG & \textbf{e.} organización independiente que ayuda a la población \\
\hline
3. la inversión & \textbf{a.} dinero que se coloca en un país para crear riqueza \\
\hline
4. el acuerdo bilateral & \textbf{b.} tratado firmado por dos Estados \\
\hline
5. la cooperación & \textbf{d.} conjunto de acciones entre países para ayudarse \\
\hline
6. el desempleo & \textbf{f.} situación de las personas que no tienen trabajo \\
\hline
\end{tabularx}
\end{center}
\emph{Traduction des définitions :} c. « argent donné par l'État sans obligation de le rembourser » ; e. « organisation indépendante qui aide la population » ; a. « argent placé dans un pays pour créer de la richesse » ; b. « traité signé par deux États » ; d. « ensemble d'actions entre pays pour s'entraider » ; f. « situation des personnes qui n'ont pas de travail ».
\end{corrige}

\begin{corrige}[Exercice 9 — Compréhension de texte type Bac]
\textbf{Barème indicatif :} compréhension 6 points (1,5 point par question : 1 point pour le sens, 0,5 point pour la correction de la langue) ; langue 4 points (1 point par question).

\textbf{Questions de compréhension — réponses modèles :}
\begin{enumerate}
\item \esp{El acuerdo se firmó en el Palacio de la Unidad de Yaundé, entre los gobiernos de España y de Camerún.} (« L'accord a été signé au Palais de l'Unité de Yaoundé, entre les gouvernements de l'Espagne et du Cameroun. »)
\item \esp{La subvención española está destinada a la formación profesional de los jóvenes, especialmente en los sectores de la agricultura y de la tecnología.}
\item \esp{Las ONG participarán en el proyecto para garantizar que la ayuda llegue directamente a la población rural.}
\item \esp{Algunos especialistas piensan que la ayuda internacional no es suficiente: creen que África necesita sobre todo más inversión extranjera y un comercio más justo con Europa.}
\end{enumerate}

\textbf{Questions de langue :}
\begin{enumerate}
\item Synonyme de \esp{trabajo} : \esp{empleo}. Contraire de \esp{posible} : \esp{imposible} (présent dans \esp{no hay desarrollo posible} $\rightarrow$ le contraire est \esp{imposible}).
\item \esp{firmaron} $\rightarrow$ \esp{firmar} ; \esp{declaró} $\rightarrow$ \esp{declarar} ; \esp{afirmó} $\rightarrow$ \esp{afirmar} ; \esp{creen} $\rightarrow$ \esp{creer}.
\item \esp{Es importante que la ayuda \textbf{llegue} a la población rural.} (Subjonctif présent de \esp{llegar}, orthographe \esp{gu} devant \esp{e}.)
\item « Sans formation, il n'y a pas de développement possible. » (ou : « Sans formation, aucun développement n'est possible. »)
\end{enumerate}
\end{corrige}

\begin{attention}[Points de vigilance — Exercice 9]
Répondez par des \textbf{phrases complètes} en espagnol, en reprenant les termes de la question. Ne recopiez pas un paragraphe entier : sélectionnez l'information. Aux infinitifs, attention : \esp{creen} vient de \esp{creer} (et non \esp{crear}).
\end{attention}

\begin{corrige}[Exercice 10 — Thème et version]
\textbf{Barème indicatif :} 1 point par phrase (0,5 pour le sens, 0,5 pour la grammaire et l'orthographe) pour le thème ; 5 points pour la version (fidélité 3 points, qualité du français 2 points).

\textbf{A. Thème — correction :}
\begin{enumerate}
\item \esp{La cooperación entre España y África es muy importante.}
\item \esp{Es necesario que los jóvenes reciban una buena formación profesional.} (« Il faut que » + subjonctif : \esp{reciban}, de \esp{recibir}.)
\item \esp{Los acuerdos multilaterales ayudan a los países pobres a desarrollarse.} (\esp{ayudar a alguien a} + infinitif ; \esp{desarrollarse} : « se développer », verbe pronominal.)
\item \esp{El desempleo de los jóvenes es un gran problema en Camerún.} (\esp{gran}, apocope de \esp{grande} devant un nom masculin singulier.)
\item \esp{Esperamos que este acuerdo cree muchos empleos.} (Après \esp{esperar que} affirmatif, le subjonctif \esp{cree} est la forme la plus sûre au Bac. \esp{crear} : subjonctif \esp{cree} \textbf{sans accent} — attention : \esp{creé} avec accent est le prétérit.)
\end{enumerate}

\textbf{B. Version — traduction modèle :}

\begin{textebox}[Traduction version — Espagnol]
La ayuda al desarrollo no consiste solo en dar dinero a los países pobres. Consiste también en compartir conocimientos, formar a los jóvenes y construir escuelas y hospitales. Las ONG desempeñan un papel fundamental porque trabajan directamente con la población. Sin embargo, para que la ayuda sea eficaz, es necesario que los gobiernos luchen contra la corrupción y que los países ricos cumplan sus promesas.
\end{textebox}
\end{corrige}

\begin{attention}[Points de vigilance — Exercice 10]
\esp{desempeñar un papel} = « jouer un rôle » (faux ami : \esp{desempeñar} ne signifie pas « désemplir »). \esp{para que la ayuda sea eficaz} : subjonctif traduit par « pour que l'aide soit efficace ». \esp{eficaz} = « efficace » (en espagnol un seul \esp{f}, un \esp{z}). \esp{cumplan} (subjonctif de \esp{cumplir}) pour « respectent/tiendront ».
\end{attention}

\begin{corrige}[Exercice 11 — Expression écrite type Bac]
\textbf{Barème indicatif (sur 10 ou 20 selon l'épreuve) :} respect du sujet et de la longueur (80--100 mots) : 2 points ; richesse du vocabulaire de la leçon : 3 points ; correction grammaticale (accords, temps, subjonctif) : 3 points ; organisation (introduction, développement, conclusion, connecteurs) : 2 points.

\textbf{Proposition de rédaction modèle (Sujet 1) :}

\begin{textebox}[Modelo de redacción — Sujet 1]
En mi opinión, la ayuda internacional es necesaria, pero no es suficiente para el desarrollo de África.

En primer lugar, esta ayuda permite financiar proyectos importantes: escuelas, hospitales y formación profesional de los jóvenes. Por ejemplo, las subvenciones europeas ayudan a muchos agricultores cameruneses. Además, las ONG trabajan sobre el terreno para que la ayuda llegue directamente a la población.

Sin embargo, África necesita sobre todo más inversión extranjera y un comercio más justo. Es necesario que los gobiernos luchen contra la corrupción para que el dinero sea realmente útil.

En conclusión, la ayuda internacional es un primer paso, pero el verdadero desarrollo depende también de los africanos y de sus gobiernos.
\end{textebox}
\emph{(environ 100 mots)}

\textbf{Proposition de rédaction modèle (Sujet 2) :}

\begin{textebox}[Modelo de redacción — Sujet 2]
El desempleo de los jóvenes es un problema grave en muchos países africanos, como Camerún. ¿Qué soluciones existen?

En primer lugar, los gobiernos deben invertir en la formación profesional para que los jóvenes aprendan oficios útiles: agricultura, tecnología, comercio. Además, es importante que el Estado conceda subvenciones a los jóvenes emprendedores que quieren crear su empresa.

Por su parte, las empresas pueden ofrecer prácticas y empleos a los jóvenes diplomados. Finalmente, la cooperación internacional, con acuerdos bilaterales como los de España y Guinea Ecuatorial, puede financiar proyectos de desarrollo.

En conclusión, la lucha contra el desempleo exige la colaboración de todos: gobiernos, empresas y socios internacionales.
\end{textebox}
\emph{(environ 100 mots)}
\end{corrige}

\begin{attention}[Points de vigilance — Exercice 11]
\begin{itemize}
\item Respectez la fourchette de mots : un texte trop court ou trop long est pénalisé.
\item Employez au moins un subjonctif correct (\esp{para que la ayuda llegue}, \esp{es necesario que los gobiernos luchen}, \esp{para que los jóvenes aprendan}, \esp{es importante que el Estado conceda}).
\item Vérifiez les accents : \esp{opinión}, \esp{formación}, \esp{inversión}, \esp{además}, \esp{conclusión}, \esp{países}, \esp{agrícolas}, \esp{diplomados}.
\item Évitez les calques du français : « selon moi » = \esp{en mi opinión} (et non \esp{según yo}) ; « créer des emplois » = \esp{crear empleo} ; « sur le terrain » = \esp{sobre el terreno} ; « jeunes diplômés » = \esp{jóvenes diplomados} (et non \esp{graduados} seul, même si possible).
\end{itemize}
\end{attention}

\newpage

\coursec{Fiche bilan}

\begin{popfigure}
\begin{tikzpicture}[
  mindmap,
  concept color=popblue,
  level 1 concept/.append style={level distance=3.5cm, sibling angle=360/7},
  level 2 concept/.append style={level distance=2.5cm},
  every node/.style={font=\small, align=center},
  branch1/.style={concept color=poporange},
  branch2/.style={concept color=popgreen},
  branch3/.style={concept color=poppurple},
  branch4/.style={concept color=poppink},
  branch5/.style={concept color=popgold},
  branch6/.style={concept color=popdark},
  branch7/.style={concept color=popblue!70!popgreen}
]
\node[concept, text=white, scale=1.2] {El mundo laboral hoy}
  child[branch1] { node[concept, align=center]{Texto soporte\\ \emph{Didáctica V, U4}} }
  child[branch2] { node[concept, align=center]{Léxico clave\\ subvención, ONG,\\ cooperación, inversión} }
  child[branch3] { node[concept, align=center]{Comprensión\\ ideas directrices,\\ datos, argumentos} }
  child[branch4] { node[concept, align=center]{Método comentario\\ 1. Contexto 2. Análisis\\ 3. Síntesis 4. Opinión} }
  child[branch5] { node[concept, align=center]{Ayudas y acuerdos\\ bilaterales / multilaterales,\\ ayuda al desarrollo} }
  child[branch6] { node[concept, align=center]{Enjeux en África\\ empleo juvenil,\\ sector informal, cacao} }
  child[branch7] { node[concept, align=center]{Producción\\ síntesis, carta formal,\\ debate oral} };
\end{tikzpicture}
\legende{Carte mentale récapitulative de la leçon E16}
\end{popfigure}

\begin{bilan}
\end{bilan}

\courssub{Lexique clé : Mundo laboral y ayuda al desarrollo}
\begin{center}
\begin{tabularx}{\linewidth}{|X|X|X|}
\hline
\rowcolor{popblueL}
\textbf{Español} & \textbf{Français} & \textbf{Contexte / Collocation} \\
\hline
\cle{la subvención} & la subvention & \esp{conceder una subvención} (accorder une subv.) \\
\hline
\cle{la ayuda al desarrollo} & l'aide au développement & \esp{países donantes / receptores} \\
\hline
\cle{el acuerdo bilateral / multilateral} & l'accord bilatéral / multilatéral & \esp{firmar un acuerdo} \\
\hline
\cle{la inversión (extranjera directa)} & l'investissement (direct à l'étranger) & \esp{atrer inversión} \\
\hline
\cle{la cooperación (internacional)} & la coopération (internationale) & \esp{agencia de cooperación} \\
\hline
\cle{la ONG} (Organización No Gubernamental) & l'ONG & \esp{ONG de desarrollo} \\
\hline
\cle{el empleo / el paro} & l'emploi / le chômage & \esp{tasa de paro, empleo precario} \\
\hline
\cle{la economía informal} & l'économie informelle & \esp{trabajar en la informalidad} \\
\hline
\cle{la formación profesional} & la formation professionnelle & \esp{curso de formación} \\
\hline
\cle{el emprendimiento} & l'entrepreneuriat & \esp{espíritu emprendedor} \\
\hline
\end{tabularx}
\end{center}

\courssub{Points de grammaire \& structures utiles (rappel Bac)}
\begin{itemize}
\item \textbf{Subjonctif :} Après expressions de volonté, doute, nécessité liées à l'aide (\esp{es necesario que la ONG \emph{gestione} los fondos}, \esp{el gobierno quiere que se \emph{firme} el acuerdo}).
\item \textbf{Passif réfléchi (se + verbe) / Passif \"être\" (ser + part.) :} \esp{Se firman acuerdos} / \esp{Los acuerdos son firmados}. Très fréquent dans les textes économiques/institutionnels.
\item \textbf{Connecteurs logiques (argumentation) :} \esp{por tanto, en consecuencia, sin embargo, no obstante, además, en primer lugar, finalmente}.
\item \textbf{Discours indirect (style indirect) :} Changements de temps, personnes, indicateurs spatio-temporels (\esp{dice que} $\rightarrow$ \esp{dijo que}; \esp{mañana} $\rightarrow$ \esp{al día siguiente}).
\end{itemize}

\courssub{Méthode express : Commentaire de texte économique}
\begin{enumerate}
\item \textbf{Contexualiser :} Source, date, auteur, thème (économie/travail), public visé.
\item \textbf{Structurer :} Identifier les parties (introduction, développement : constats / causes / solutions, conclusion).
\item \textbf{Synthétiser :} Relever les \cle{idées directrices} (une par paragraphe) + données chiffrées / exemples.
\item \textbf{Analyser le lexique :} Champ lexical (travail, aide, accords), registres (formel, technique), figures de style.
\item \textbf{Évaluer / Réagir :} Pertinence, actualité (Cameroun, Guinée équatoriale), limites, ouverture personnelle.
\end{enumerate}


\begin{aretenir}[Checklist avant le Bac]
$\square$ Je sais définir en espagnol : \esp{subvención, ayuda al desarrollo, acuerdo bilateral, ONG, inversión, cooperación}. \\
$\square$ Je reconnais et j'utilise le vocabulaire du monde du travail : \esp{empleo, paro, economía informal, formación profesional, emprendimiento}. \\
$\square$ Je maîtrise la structure type d'un commentaire de texte (contexte, résumé structuré, analyse, évaluation). \\
$\square$ J'identifie les \cle{idées directrices} et les \cle{connecteurs logiques} dans un texte argumentatif. \\
$\square$ Je sais passer du style direct au style indirect (changements de temps, pronoms, indicateurs). \\
$\square$ J'utilise correctement la voix passive (\esp{se + verbe} / \esp{ser + participio}) pour décrire des mesures officielles. \\
$\square$ Je conjugue sans erreur les verbes clés au subjonctif présent : \esp{conceder, firmar, gestionar, requerir, ser, estar, haber}. \\
$\square$ Je sais rédiger une synthèse de 80--100 mots en espagnol (vocabulaire précis, neutre, connecteurs). \\
$\square$ Je suis capable de situer les enjeux de l'emploi au Cameroun et en Guinée équatoriale (jeunes, secteur informel, coopération Sud-Sud / Nord-Sud). \\
$\square$ Je prépare des arguments pour un débat oral : \esp{¿Ayuda o comercio? ¿Formalización de la economía?} \\
$\square$ Je relis mes productions pour corriger les \textbf{faux amis} (\esp{subvención} $\neq$ subversion, \esp{éxito} $\neq$ sortie, \esp{compromiso} $\neq$ compromis). \\
$\square$ Je vérifie les accents toniques et les signes \esp{¿ ¡} dans chaque phrase espagnole écrite.
\end{aretenir}

\findecours

\end{document}
